\documentclass[10pt]{article}

\newcommand\hmmax{0}
\newcommand\bmmax{2}
\usepackage{amsthm, amsmath,amssymb,bbm}
\DeclareMathAlphabet{\mathbbm}{U}{bbm}{m}{n}
\usepackage{natbib}
\usepackage{multirow}
\usepackage{subfigure}
\usepackage{makecell}
\usepackage{booktabs}
\usepackage{array}
\usepackage{tabularx}
\usepackage{tabulary}
\usepackage{caption}
\usepackage{booktabs}
\usepackage{url}
\usepackage{algorithm}
\usepackage{algpseudocode}
\usepackage{bm}
\usepackage{wrapfig}
\usepackage{lipsum}
\usepackage{mathrsfs}
\usepackage{dsfont}
\usepackage{titling}
%\usepackage{txfonts}
\usepackage{relsize}
\usepackage{enumitem}
\usepackage{setspace}
\usepackage[colorlinks,
linkcolor=tpurple,
anchorcolor=tpurple,
citecolor=tpurple
]{hyperref}

\usepackage{tikz}
\usetikzlibrary{positioning}

% arxiv 
\usepackage[nameinlink,capitalize]{cleveref}
% journal
%\usepackage[capitalize]{cleveref}

%%%%%%% Setup of User-Defined Colors

\definecolor{myred}{HTML}{ae1908}
\definecolor{myblue}{HTML}{05348b}
\definecolor{myorange}{HTML}{ec813b}
\definecolor{mylightblue}{HTML}{9acdc4}
\definecolor{tpurple}{HTML}{743096}
\definecolor{myyellow}{HTML}{e5a84b}
\definecolor{mygreen}{HTML}{6bb392}

\definecolor{revisecolor}{HTML}{000000}



%%%%%%% Setup of Colorlinks, Margin, Textwidth, Font, Numbering ...

%\setstretch{2}
%\usepackage{geometry}
% \geometry{
% a4paper,
 %total={170mm,257mm},
% left=28mm,
% top=30mm,
% }
\textwidth=6in
\usepackage{fullpage}
%\setstretch{2}

% caption size
\usepackage[font={footnotesize,it}]{caption}

\numberwithin{equation}{section}


%%%%%%% My macros
\usepackage{smile}
\newcommand{\indep}{\perp \!\!\! \perp}
\newcommand*{\supp}{\mathrm{supp}}
\usepackage{etoolbox}
\newtoggle{vaos}
\togglefalse{vaos}
\newcommand{\aosversion}[2]{\iftoggle{vaos}{#1}{#2}}

% Definition of Roman Numbers
\newcommand{\Rom}[1]{\text{\uppercase\expandafter{\romannumeral #1\relax}}}
\newcommand{\todo}[1]{{\color{tpurple} \colorbox{tpurple!30}{Todo} #1}}
\newcommand{\question}[1]{{\color{tpurple} \colorbox{tpurple!30}{Question} #1}}
\newcommand{\red}[1]{{\color{red} #1}}
\newcommand{\blue}[1]{{\color{blue} #1}}


\newcommand{\myred}[1]{{\color{myred} #1}}
\newcommand{\myblue}[1]{{\color{myblue} #1}}
\newcommand{\mylightblue}[1]{{\color{mylightblue} #1}}
\newcommand{\myorange}[1]{{\color{myorange} #1}}
\newcommand{\mypurple}[1]{{\color{mypurple} #1}}
\newcommand{\myyellow}[1]{{\color{myyellow} #1}}
\newcommand{\mygreen}[1]{{\color{mygreen} #1}}

\def\r#1{\textcolor{red}{\bf #1}}
\def\b#1{\textcolor{blue}{\bf #1}}
\begin{document}\appendix

\section*{Supplemental Material}

The supplemental materials collect the complete theoretical analysis (\cref{sec:theory:general}), all the population-level proofs (\cref{sec:proof-population}), finite sample proofs (\cref{sec:proof-nonasymptotic}), and omitted discussions in the main text (\cref{sec:discussion}).


% BEGIN INCLUDED SOURCE: tex/supplement_text.tex

\section{General Theoretical Analysis}

\label{sec:theory:general}

\subsection{General Notations and Conditions}


Define the pooled covariance matrix $\bar{\bSigma} = \sum_{e\in \mathcal{E}} \omega^{(e)} \bSigma^{(e)}$. We first state the standard assumptions used in linear regression, which are analogous to \cref{cond0-model-main}--\ref{cond3-sub-gaussian-eps-main}. \cref{cond1-well-condition} and \cref{cond3-sub-gaussian-eps} are just copies of \cref{cond1-well-condition-main} and \cref{cond3-sub-gaussian-eps-main}, respectively. \cref{cond0-model} allows for varying $(n^{(e)}, \omega^{(e)})$. \cref{cond2-subgaussain-x} replaces the $\bSigma$ in \cref{cond2-subgaussain-x-main} by $\bar{\bSigma}$.

\begin{condition}
\label{cond0-model}
	For each $e\in \mathcal{E}$, $(\bx_1^{(e)}, y_1^{(e)}), \ldots, (\bx_{n^{(e)}}^{(e)}, y_{n^{(e)}}^{(e)})$ are i.i.d. copies of $(\bx^{(e)}, y^{(e)}) \sim \mu^{(e)}$, where $\mu^{(e)}$ belongs to $\mathcal{U}_{\bbeta^*, \sigma^2}$ for some $\sigma^2$. The data from different environments are also independent. We have $\omega^{(e)} > 0$ for any $e\in \mathcal{E}$. 
\end{condition}

\begin{condition}
\label{cond1-well-condition}
	There exists some universal constants $\kappa_L \in (0,1]$ and $\kappa_U \in [1,\infty)$ such that
	\begin{align}
		\forall e\in \mathcal{E}, ~~~~~~~~~~ \kappa_L \bI_p \preceq \bSigma^{(e)} \preceq \kappa_U \bI_p.
	\end{align}
\end{condition}

\begin{condition}
\label{cond2-subgaussain-x}
	There exists some universal constant $\sigma_x \in [1,\infty)$ such that
	\begin{align}
		\forall e\in \mathcal{E}, \bv \in \mathbb{R}^p, ~~~~~~~~~~ \mathbb{E} \left[\exp\left\{\bv^\top \bar{\bSigma}^{-1/2}\bx^{(e)}\right\}\right] \le \exp\left(\frac{\sigma_x^2}{2} \cdot \|\bv\|_2^2\right).
	\end{align}
\end{condition}

\begin{condition}
\label{cond3-sub-gaussian-eps}
	There exists some universal constant $\sigma_\varepsilon \in \mathbb{R}^+$ such that,
	\begin{align}
		\forall e\in \mathcal{E}, \lambda \in \mathbb{R}, ~~~~~~~~~~ \mathbb{E} [e^{\lambda \varepsilon^{(e)}}] \le e^{\frac{1}{2} \lambda^2 \sigma_\varepsilon^2}.
	\end{align}
\end{condition}

We also define several quantities regarding sample size:
\begin{align}
    \label{eq:def-n-low-dim}
    n_* =\min_{e\in \mathcal{E}} \frac{n^{(e)}} {\omega^{(e)}} \qquad \bar{n} = \Bigg(\sum_{e\in \mathcal{E}}\frac{\omega^{(e)}}{n^{(e)}}\Bigg)^{-1} \qquad n_{\min} = \min_{e\in \mathcal{E}} n^{(e)} \qquad n_\dagger = \Bigg( \sum_{e\in \mathcal{E}} \frac{\omega^{(e)}}{(n^{(e)})^{3/2}} \Bigg)^{-1}.
\end{align}
It is easy to see that $n_* \ge \bar{n} \ge n_{\min}$. In the case of balanced data with equal weights, i.e., $n^{(e)} \equiv n$ and $\omega^{(e)} \equiv 1/|\mathcal{E}|$, we have $n_* = n\cdot |\mathcal{E}|$, $\bar{n}=n_{\min} = n$ and $n_\dagger=n^{3/2}$.

We also define the {\it pooled linear spurious variable} concerning the environment weights $\bomega$.

\begin{definition}[Pooled Linear Spurious Variables] We let $G_{\bomega}$ be the index set of all pooled linear spurious variables in environments $\mathcal{E}$ concerning the weights $\bomega$, that is,
		$G_{\bomega} = \{j \in [p]:  \sum_{e\in \mathcal{E}} \omega^{(e)} \mathbb{E}[x_j^{(e)}\varepsilon^{(e)}] \neq 0 \}$. We say a variable $x_j$ is a pooled linear spurious variable if $j\in G_{\bomega}$.
\end{definition}

We need the following condition related to the heterogeneity of environments in the presence of pooled linear spurious variables to recover $\bbeta^*$.

\begin{condition}[Identification]
\label{cond-ident}
    For any $S\subseteq [p]$ satisfying $S \cap G_{\bomega} \neq \emptyset$, there exists some $e, e'\in \mathcal{E}$ such that $\bbeta^{(e, S)} \neq \bbeta^{(e',S)}$, where $\bbeta^{(e,S)}$ is defined in \eqref{eq:local-minimizer}.
\end{condition} 

\subsection{Strong Convexity with respect to the True Parameter}

We are now ready to state the main population-level result, a generalized version of \cref{thm:global-strong-convexity-main} with varying environments coefficients $\bomega$.

\begin{theorem}[{Strong Convexity with respect to} $\bbeta^*$] 
\label{thm:global-strong-convexity}
Assume Conditions \ref{cond0-model}--\ref{cond1-well-condition} and \ref{cond-ident} hold. Then $\bbeta^*$ is the unique minimizer of $\mathsf{Q}(\bbeta;\gamma,\bomega)$ for large enough $\gamma$: for any $\epsilon \in (0, 1)$ and any $\gamma\ge \epsilon^{-1} \gamma^*$ with
\begin{align}
\label{eq:thm1-gamma-star}
	\gamma^* = (\kappa_L)^{-3} \sup_{S: S\cap G_{\bomega} \neq \emptyset} \left(\mathsf{b}_S/\bar{\mathsf{d}}_S\right),
\end{align} where
\begin{align}
\label{eq:def-bias-bias-difference}
 \mathsf{b}_S = \Bigg\|\sum_{e\in \mathcal{E}} \omega^{(e)} \mathbb{E} [\varepsilon^{(e)} \bx_S^{(e)}]\Bigg\|^2_2 \qquad \text{and} \qquad \bar{\mathsf{d}}_S = \sum_{e\in \mathcal{E}} \omega^{(e)} \big\|\bbeta^{(e,S)} - \bar{\bbeta}^{(S)}\big\|_2^2  
\end{align} 
 with $\bar{\bbeta}^{(S)}=\sum_{e'\in \mathcal{E}} \omega^{(e')} \bbeta^{(e',S)}$, we have
\begin{align}
\label{eq:thm1-jsc-condition}
	\mathsf{Q}(\bbeta;\gamma, \bomega) - \mathsf{Q}(\bbeta^*;\gamma, \bomega) \ge (1-\epsilon) \|\bar{\bSigma}^{1/2}(\bbeta - \bbeta^*)\|_2^2 + \kappa_L^2 (\gamma - \epsilon^{-1}\gamma^*) \bar{\mathsf{d}}_{\supp(\bbeta)}.
\end{align}
\end{theorem}

It is easy to see that \cref{thm:global-strong-convexity-main} is a direct corollary of the above \cref{thm:global-strong-convexity} with $\omega^{(e)} \equiv 1 / |\mathcal{E}|$.

\subsection{Non-asymptotic Analysis}

The first proposition characterizes the convergence of the pooled least squares.


\begin{proposition}[Properties of Pooled Least Squares]
\label{prop:pooled-least-squares}
	Assume Conditions \ref{cond0-model}--\ref{cond3-sub-gaussian-eps} hold. Then, there exists some $\bar{\bbeta}^{\mathsf{R}} \in \mathbb{R}^p$ satisfying $\frac{1}{\kappa_U} \left\|\sum_{e\in \mathcal{E}} \omega^{(e)} \mathbb{E} [\varepsilon^{(e)} \bx^{(e)}] \right\|_2 \le \|\bar{\bbeta}^{\mathsf{R}} - \bbeta^*\|_2 \le \frac{1}{\kappa_L} \left\|\sum_{e\in \mathcal{E}} \omega^{(e)} \mathbb{E} [\varepsilon^{(e)} \bx^{(e)}] \right\|_2$
 	such that, for any $\bbeta \in \mathbb{R}^p$,
	\begin{align}
	\label{eq:pooled-least-square-strong-convexity}
		\mathsf{R}(\bbeta) - \mathsf{R}(\bar{\bbeta}^{\mathsf{R}}) = \|\bar{\bSigma}^{1/2}(\bbeta - \bar{\bbeta}^{\mathsf{R}})\|_2^2.
	\end{align} Moreover, there exist universal constants $c_1$ and $c_2$ such that if $n^* \ge c_1 \sigma_x^4(p + t)$, then the pooled least      squares estimator $\hat{\bbeta}_{\mathsf{R}}$ minimizing \eqref{eq:pooled-l2-risk-empirical} satisfies
	\begin{align}
		\|\bar{\bSigma}^{1/2}(\hat{\bbeta}_{\mathsf{R}} - \bar{\bbeta}^{\mathsf{R}})\|_2 \le c_2 \sigma_x \left(\sigma_\varepsilon + \sigma_x \|\bar{\bSigma}^{1/2}(\bar{\bbeta}^{\mathsf{R}} - \bbeta^*)\|_2\right) \sqrt{\frac{p + t}{n_*}}
	\end{align} with probability $1-2e^{-t}$.
\end{proposition}

The following two theorems are generalized versions of \cref{thm:rate-low-dim-main} and \cref{thm:rate-faster-low-dim-main} with varying $n^{(e)}$ and $\omega^{(e)}$. 

\begin{theorem}[Non-asymptotic Variable Selection Property]
\label{thm:rate-low-dim} Define
\begin{align}
\label{eq:thm2-signal}
    \mathsf{s}_+ = \min_{j\in S^*} |\beta_j^*|^2 \qquad \text{and} \qquad \mathsf{s}_- = \min_{S\subseteq [p], S\cap G_{\bomega} \neq \emptyset} \bar{\mathsf{d}}_S
\end{align} 
Suppose Conditions \ref{cond0-model}--\ref{cond-ident} hold, and we choose $\gamma \ge 3\gamma^*\lor 1$ where $\gamma^*$ is defined in \cref{thm:global-strong-convexity}. There exists some universal constants $c_1$--$c_2$ that only depends on $(\kappa_U, \sigma_x, \sigma_\varepsilon)$ such that for any $t>0$, if $n_{\min} \ge p + \log(2|\mathcal{E}|) + t$, $\bar{n} \ge c_1(\gamma/\kappa_L)(p+ \log(|\mathcal{E}|) + t)\{\mathsf{s}_+^{-0.5}+\mathsf{s}_+^{-1}+(\gamma\kappa_L\mathsf{s}_-)^{-0.5}\}$, and $n_* \ge c_2(\gamma/\kappa_L)^2(p+t)\{\mathsf{s}_+^{-1} + (\gamma\kappa_L\mathsf{s}_-)^{-1} + 1\}$, then the EILLS estimator $\hat{\bbeta}_{\mathsf{Q}}$ minimizing \eqref{eq:eills-objective-q} satisfies
\begin{align}
\label{eq:low-dim-variable-selection}
    \mathbb{P}\left[S^* \subseteq \supp(\hat{\bbeta}_{\mathsf{Q}}) \subseteq (G_{\bomega})^c \right] \ge 1-7e^{-t}.
\end{align}
\end{theorem}

Combining \cref{thm:rate-low-dim} and \cref{prop:pooled-least-squares}, we can conclude that when $(n_{\min}, \bar{n}, n_*)$ are large enough, one can refit the solution found by the EILLS estimator to obtain a faster rate. To be specific, we can run another pooled least squares on $\bx_{\hat{S}}$ with $\hat{S}=\supp(\hat{\bbeta}_{\mathsf{Q}})$ and obtain the $\ell_2$-error $\{|(G_{\bomega})^c|/n_*\}^{1/2}$ given that $S^*\subseteq \hat{S}\subseteq (G_{\bomega})^c$. 

\begin{theorem}[Non-asymptotic $\ell_2$ Error Bound] 
\label{thm:rate-faster-low-dim}
Assume Conditions \ref{cond0-model}--\ref{cond-ident} hold, and we choose $\gamma \ge 3\gamma^*\lor 1$ where $\gamma^*$ is defined in \cref{thm:global-strong-convexity}. There exists some universal constants $c_1$--$c_4$  that only depend on $(\kappa_U, \sigma_x)$ such that for any $t>0$, if $n_{\min} \ge p + \log(2|\mathcal{E}|) + t$ and $n_* \ge c_1 (\gamma/\kappa_L)(p + t)$, then $\hat{\bbeta}_{\mathsf{Q}}$ minimizing \eqref{eq:eills-objective-q} satisfies
\begin{align}
\label{eq:low-dim-rate2}
    \frac{\|\hat{\bbeta}_{\mathsf{Q}} - \bbeta^*\|_2}{\sigma_\varepsilon (\gamma/\kappa_L)}  &\le c_2 \left(\sqrt{\frac{p+t}{n_*}} + \frac{p+\log(|\mathcal{E}|)+t}{\bar{n}}\right) + c_3 \frac{\sqrt{|S^*|}}{\bar{n}} \cdot \frac{\log(|\mathcal{E}||S^*|)+t}{\min_{j\in S^*}|\beta^*_j|}
\end{align} with probability at least $1-7e^{-t}$. Moreover, when the additional conditions in \cref{thm:rate-low-dim} hold, then
\begin{align}
\label{eq:low-dim-rate1}
    \frac{\|\hat{\bbeta}_{\mathsf{Q}} - \bbeta^*\|_2}{\sigma_\varepsilon (\gamma/\kappa_L)}  &\le c_2 \left(\sqrt{\frac{p_0+t}{n_*}} + \frac{p_0+\log(|\mathcal{E}|)+t}{\bar{n}}\right) \qquad \text{with}~~p_0 = |({G}_{\bomega})^c|
\end{align} occurs with probability at least $1-14e^{-t}$. The dependency of $(c_1,c_2,c_3)$ on $(\kappa_U,\sigma_x)$ can be found in~\eqref{eq:low-dim-faster-rate-c1} and~{\eqref{eq:low-dim-faster-rate-c2-4}} in \cref{sec:proof-thm3}. 
\end{theorem} 

In the high-dimensional regime, recall that $s^* = |S^*|$, $\beta_{\min} = \min_{j\in S^*} |\beta_j^*|$. We need the following condition regarding the sample size defined above in \eqref{eq:def-n-low-dim}.


\begin{condition}
\label{cond6-n-req}
    Suppose that $\log(|\mathcal{E}|)\le C\log p$ and that 
    
\noindent (1) $n_{\min} \ge c_1 (s^* + \log p)$,
    
\noindent (2) $n_* \ge c_2 (\gamma/\kappa_L)^2 \Big(s^* \log p \Big) \Big\{1 + 1\big/ \big(\kappa_L \beta_{\min}^2 \big)\Big\}$, 
    
\noindent (3) $\bar{n} \ge c_3 (\gamma/\kappa_L) (\log p) \left\{s^* + 1/ (\beta_{\min}^2)\right\}$ and $\bar{n} \ge c_3 (\gamma/\kappa_L) \sqrt{(s^*\log p)(s^* + \log p)} \big/ \big(\kappa_L \beta_{\min} \big)$, 
    
\noindent (4) $n_\dagger \ge c_4(\gamma/\kappa_L) \big(s^* \log p \big)\sqrt{s^* + \log p}\Big\{1+1\big/ \big(\sqrt{\kappa_L} \beta_{\min}\big)\Big\}$.
    
\noindent Here $c_1$--$c_4$ are universal positive constants that depend only on $(C,\sigma_x,\kappa_U, \sigma_\varepsilon)$. See details in~{\eqref{eq:req-n-high-dim}}.
\end{condition}

The above \cref{cond6-n-req} matches \cref{cond6-n-req-main} in the case of balanced data with equal weights. The following theorem claims that EILLS can attain variable selection consistency with proper choice of the model selection hyper-parameter $\lambda$.

\begin{theorem}[Variable Selection Consistency in High Dimensions]	
\label{thm:high-dim}  Assume Conditions \ref{cond0-model}--\ref{cond6-n-req} hold, and we choose $\gamma \ge 3\gamma^* \lor 1$ where $\gamma^*$ is defined in \cref{thm:global-strong-convexity}. Suppose further that the choice of $\lambda$ satisfies
\begin{align}
    c_1 \left\{ (\gamma/\kappa_L)^2 \frac{s^*\log p}{n_*} + \epsilon(\bar{n}, n_\dagger)\right\} \le \lambda \le c_2 \kappa_L \beta_{\min}^2,
\end{align} where $\epsilon(\bar{n}, n_\dagger) = (\gamma/\kappa_L)^2 (\log p) (s^* + \log p) \{\bar{n}^{-2}s^* + n_\dagger^{-2} (s^*)^2\log p + (n_\dagger \sqrt{s^* + \log p})^{-1}(\gamma/\kappa_L)^{-1} \}$, and  $c_1, c_2$ are some universal positive constants only depends on $(C,\kappa_U, \sigma_x,\sigma_\varepsilon)$. Then the $\ell_0$ regularized EILLS estimator $\hat{\bbeta}_{\mathsf{L}}$ solving \eqref{eq:eills-objective-l} satisfies 
\begin{align*}
    \mathbb{P}\left[\supp(\hat{\bbeta}_{\mathsf{L}}) = S^*\right] \ge 1-p^{-10}.
\end{align*}
\end{theorem}


\subsection{Examples of \cref{thm:global-strong-convexity}}

\label{sec:moreex}

\subsubsection{Linear SCM with Heterogeneous Variance}

We start with a toy example similar to \cref{ex-h-de}.

\begin{example}[Toy Example, Heterogeneous Variances]
\label{ex-h-var-toy}
	 Consider the following multi-environment linear model with $\mathcal{E}=\{1,2\}$ and $p=2$,
	 \begin{align*}
	 	x^{(e)}_1 &\gets \sqrt{v_1} \cdot \varepsilon_1 \\
	 	y^{(e)} &\gets 1\cdot x^{(e)}_1 + \myred{\sqrt{v_0^{(e)}}} \cdot \varepsilon_0 \\
	 	x^{(e)}_2 &\gets h \cdot x^{(e)}_1 + s \cdot y^{(e)} + \sqrt{v_2} \cdot \varepsilon_2	
	 \end{align*}
	for some parameters $s, h \in \mathbb{R}$ and $v_1, v_2 \in \mathbb{R}^+$ that is invariant across different environments, and one parameter $v_0^{(e)}$ that will vary for different environments. Here $\varepsilon_0, \varepsilon_1, \varepsilon_2$ are independent, standard normal distributed exogenous variables. We let $\omega^{(e)}\equiv 1/2$.

    {\normalfont In the example, $S^*=\{1\}$, $\bbeta^* = (1,0)^\top$, and $\varepsilon^{(e)} = (v_0^{(e)})^{1/2} \cdot \varepsilon_0$. Hence $G_{\bomega} = \{2\}$ whenever $s\neq 0$ since $\mathbb{E}[\varepsilon^{(e)} x_2^{(e)}] = s v_0^{(e)}$. %By some calculations, we have
	In a well-conditioned regime where $|h|+|s| \lesssim 1$ and $v_1 \asymp v_2 \asymp v_0^{(e)} \asymp 1$, we have
	\begin{align*}
		\sqrt{\frac{\mathsf{b}_{\{2\}}}{\bar{\mathsf{d}}_{\{2\}}}} %&= \frac{(v_1(s+h)^2 + v_2 + s^2 v_0^{(1)}) (v_1(s+h)^2 + v_2 + s^2 v_0^{(2)}) (v_0^{(1)} + v_0^{(2)})}{\big|(v_0^{(1)} - v_0^{(2)})(h(s+h)v_1 + v_2)\big|} \\
		&\asymp \frac{1}{|v_0^{(1)} - v_0^{(2)}|} \frac{1}{|h(s+h)v_1 + v_2|} ~~~~~~ \text{and} ~~~~~~ \sqrt{\frac{\mathsf{b}_{\{1,2\}}}{\bar{\mathsf{d}}_{\{1,2\}}}} \asymp \frac{1}{|v_0^{(1)} - v_0^{(2)}|}.
	\end{align*}
	Therefore, \cref{cond-ident} holds if 
	\begin{align}
	\label{eq:toy-example-identification}
		v_0^{(1)} \neq v_0^{(2)} ~~~~~~ \text{and} ~~~~~~ h (h+s) v_1 + v_2 \neq 0 
	\end{align} and the quantity $\gamma^*$ is of constant order if
	\begin{align}
	\label{eq:toy-example-constant-order-gamma-star}
		|v_0^{(1)} - v_0^{(2)}|	= \Omega(1) ~~~~~~ \text{and} ~~~~~~ \left|s - \left(-\frac{v_2}{v_1 h} - h\right)\right| = \Omega(1).
	\end{align} The first condition, $v^{(2)}_0$ should not be close to $v_0^{(1)}$, is easy to understand. This is because when $v_0^{(2)}$ is very close to $v_0^{(1)}$, the distribution of the two environments might be very similar, which will make it hard to distinguish. The second condition, the direct effect of $y$ on $x_2$ must be apart from $(-v_2 / (v_1 h) - h)$, is strange at first glance. Let us see what it implies when $s$ coincides with $(-v_2 / (v_1 h) - h)$. If $s=(-v_2 / (v_1 h) - h)$, then the heterogeneity of the environments (from $v_0^{(1)}$ to $v_0^{(2)}$) will not affect the solution $\bbeta^{(e,\{1,2\})}$ -- it will be fixed as $\bbeta^{(e,\{1,2\})} \equiv (0, s^{-1})^\top$. In other words, though the distributions of $(\bx^{(1)},y^{(1)})$ and $(\bx^{(2)},y^{(2)})$ differs a lot, the two identity equations hold
		\begin{align*}
			\mathbb{E} [y^{(e)}|x_1^{(e)}] \equiv 1\cdot x_1^{(e)} ~~~~~~ \text{and} ~~~~~~ \mathbb{E}[y^{(e)}|x_2^{(e)}] \equiv s^{-1}\cdot x_2^{(e)}.
		\end{align*} This will be problematic since the two sets $\{1\}$ and $\{2\}$ are all CE-invariant. In this case, one can still show that $\bbeta^*$ is the unique minimizer of \eqref{eq:eills-objective-q-population} if and only if
		$(v_0^{(1)} + v_0^{(2)})/2 < s^{-2}(h^2 v_1 + v_2) = v_1^2 h^2 / (v_1 h^2 + v_2)$, that is, whether the $L_2$ risk of $\bbeta^*$ is smaller than the $L_2$ risk of the spurious solution $(0, s^{-1})^\top$. But it is beyond the scope of our discussion. 
	
	Moreover, the condition \eqref{eq:toy-example-constant-order-gamma-star} is independent of $|s|$ that measures the magnitude of spuriousness. In particular, we do not need to use a very large $\gamma^*$ when $|s|$ is very small. This is because both $\bar{\mathsf{d}}_S$ and $\mathsf{b}_S$ grow linearly with $s^2$ around $s=0$, and the choice of $\gamma^*$ is their ratio.}
	\qed
\end{example}

Given the intuitions of the above \cref{ex-h-var-toy}, we are ready to present a clean condition for a generalization of the model in \cref{ex-h-var} with $p \ge 3$.

\begin{example}[General Linear SCM with intervention on the scale of $\varepsilon^{(e)}$.]
\label{ex-h-var}
	Consider the following two-environment model	 $\mathcal{E}=\{1,2\}$ for $p$-dimensional covariate $\bx$ and response variable $y$,
	\begin{align}
	\label{eq:model-linear-scm-varying-v0}
	\begin{split}
		\bx^{(e)} &= \bB \bx^{(e)} + \balpha y^{(e)} + \bvarepsilon_x,\\
		y^{(e)} &= (\bbeta^*)^\top \bx^{(e)} + \varepsilon_0^{(e)}
	\end{split}
	\end{align} where $\bB\in \mathbb{R}^p$, $\balpha \in \mathbb{R}^p$ and $\bbeta^*$ are the same across the two environments, the exogenous variables $(\bvarepsilon_x, \varepsilon_0^{(e)}) \in \mathbb{R}^{p}\times \mathbb{R}$ are independent, and the covariance matrix of $\bvarepsilon_x$, $\bD = \mathbb{V}[\bvarepsilon_x]$, is also fixed for the two environments. The only heterogeneity comes from the variance of the noise $\varepsilon_0^{(e)}$: $v_0^{(e)}=\mathbb{V}[\epsilon_{0}^{(e)}]$ is different for $e\in \{1,2\}$. We assume the induced graph is acyclic, and use $\bW = (\bI - \bB - \balpha (\bbeta^*)^\top)^{-1}$ to denote the total effect matrix for the covariate $\bx$, that is, $W_{i,j}$ is the total effect of the variable $x_j$ on variable $x_i$.
\end{example}

\begin{proposition} 
\label{prop:ex-h-var}
Consider the two-environment model in \cref{ex-h-var}, let $\omega^{(e)}\equiv 1/2$, suppose further that \begin{align}
\label{eq:ex2-req1}
	\max_{i\in [p]} D_{i,i} \lor \|\bW\|_2 \lor \|\balpha\|_2 \le c_1	
\end{align} for some universal constant $c_1>0$. For any $S$, define
\begin{align*}
\xi(S) = 1 - \balpha^\top \bW_{S,:}^\top \left(\bW_{S,:} \bD \bW_{S,:}^\top \right)^{-1} \bW_{S,:} \bD \bW_{T,:}^\top \bbeta_T^* ~~~~ \text{with} ~~~~T=S^*\setminus S.
\end{align*} If $\xi(S) \neq 0$ for any $S \subseteq [p]$, then \cref{cond-ident} holds, and the property \eqref{eq:thm1-jsc-condition} in \cref{thm:global-strong-convexity} holds with
\begin{align}
\label{eq:ex2-cond}
	\gamma^* = c_2 \left(\min_{i\in [p]} D_{i,i}\right)^{-7} &\times \left\{\frac{(v_0^{(1)} + v_0^{(2)})(1 + v_0^{(1)})(1 + v_0^{(2)})}{|v_0^{(1)} - v_0^{(2)}|} \times \sup_{S \subseteq [p]} \frac{1}{\left|\xi(S)\right|}\right\}^2
\end{align} for some constant $c_2>0$ depends only $c_1$.
\end{proposition}

When $\min_{i\in [p]} D_{i,i} \gtrsim 1$ and $v_0^{(e)} \lesssim 1$, \cref{prop:ex-h-var} suggests the following requirements 
\begin{align*}
	| v_0^{(1)} - v_0^{(2)} | = \Omega(1) ~~~~~~ \text{and} ~~~~~~ \inf_{S\subseteq [p]} |\xi(S)| = \Omega(1)
\end{align*} are needed to get a constant order $\gamma^*$, here $\xi(S)=1$ if $S^* \subseteq S$ or $S\cap G_{\bomega} = \emptyset$. Compared with the conditions in \eqref{eq:toy-example-constant-order-gamma-star}, the first condition is the same, and the condition $\inf_{S\subseteq [p]} |\xi(S)| = \Omega(1)$ generalizes the second condition of \eqref{eq:toy-example-constant-order-gamma-star} to the multivariate case with complicated and arbitrary variable dependencies. Similar to \cref{prop:ex-h-var}, the choice of $\gamma^*$ is independent of $\|\balpha\|_2$ when $\|\balpha\|_2 \lesssim 1$, which means one do not need to use a large $\gamma^*$ when the magnitude of spuriousness $\|\balpha\|_2$ is small, i.e., $\|\balpha\|_2 \approx 0$. 


\section{Proofs for Population-level Results}
\label{sec:proof-population}

\subsection{Proof of \cref{prop:oos-definition-property}}
\label{prop:oos-definition-property-proof}

We first prove the equality $\mathsf{R}_{\mathtt{oos}}(\bbeta^*;\mathcal{U}_{\bbeta^*, \sigma^2})=\sigma^2$. The definition of $\mathcal{U}_{\bbeta^*, \sigma^2}$ implies that,
\begin{align*}
    \mathsf{R}_{\mathtt{oos}}(\bbeta^*;\mathcal{U}_{\bbeta^*, \sigma^2}) = \sup_{\mu \in \mathcal{U}_{\bbeta^*, \sigma^2}} \mathbb{E}_{(\bx,y)\sim \mu} \left[ |y - \mathbb{E}[y|\bx_S^*] |^2 \right] = \sup_{\mu \in \mathcal{U}_{\bbeta^*, \sigma^2}} \mathrm{var}_\mu [y|\bx_{S^*}] \le \sigma^2.
\end{align*} Moreover, denote $\nu \sim \mathcal{N}(0, \sigma^2 I_{p})$ as a Gaussian distribution on $\mathbb{R}^{p}$. Let $\mu_*(d\bx, y) = \nu(\bx) \nu_y(d\by|\bx)$, where $\nu_y$ is also a Gaussian distribution with mean $(\bbeta^*)^\top \bx$ and variance $\sigma^2$, it is easy to show that $\mu_* \in \mathcal{U}_{\bbeta^*, \sigma^2}$ and $\mathbb{E}_{\mu_*} [|y - (\bbeta^*)^\top \bx|^2] = \sigma^2$.Combining with the above upper bound, we can conclude that $\mathsf{R}_{\mathtt{oos}}(\bbeta^*;\mathcal{U}_{\bbeta^*, \sigma^2})=\sigma^2$.

We next prove the upper bound part of \eqref{eq:optimal-beta-s-star}. For any $\mu$, we have
\begin{align*}
    \mathbb{E}_{(\bx,y)\sim \mu} \left[ |y - \bbeta^\top \bx|^2 \right] &= \mathbb{E}\left[\left(y - (\bbeta^*)^\top\bx + (\bbeta^*)^\top\bx - \bbeta^\top \bx\right)^2\right] \\
    &= \mathbb{E}\left[\left(y - (\bbeta^*)^\top\bx \right)^2\right] + \mathbb{E} \left[ \left((\bbeta^*)^\top\bx - \bbeta^\top \bx\right)^2 \right] \\ 
    & ~~~~~~~~~~~~~~~~~~ + 2 \mathbb{E}\left[\left(y - (\bbeta^*)^\top\bx \right) \left((\bbeta^*)^\top\bx - \bbeta^\top \bx\right) \right] \\
    &\overset{(a)}{\le} \sigma^2 + \lambda_{\max} \left(\mathbb{E}\left[\bx\bx^\top\right]\right) \|\bbeta - \bbeta^*\|^2_2 \\
    &~~~~~~~~~~~~~~~~~~~ + 2\sigma \sqrt{\lambda_{\max} \left(\mathbb{E}\left[\bx\bx^\top\right]\right) \|(\bbeta - \bbeta^*)_{[p]\setminus S^*}\|^2_2} \\
    &\le \mathsf{R}_{\mathtt{oos}}(\bbeta^*;\mathcal{U}_{\bbeta^*, \sigma^2}) + p^2\sigma^2 \|\bbeta - \bbeta^*\|_2^2 + 2\sigma^2 p  \|(\bbeta - \bbeta^*)_{[p]\setminus S^*}\|_2
\end{align*} where $(a)$ follows from the fact $\mathbb{E}[(y-(\bbeta^*)^\top)\bx|\bx_{S^*}]=0$ and Cauchy-Schwarz inequality.
For the lower bound part, 
Now that $\varepsilon = y - (\bbeta^*)^\top \bx$ is independent of $\bx$ under $\mu_*$  yields, for any $\bbeta \in \mathbb{R}^p$,
\begin{align*}
    \mathsf{R}_{\mathtt{oos}}(\bbeta;\mathcal{U}_{\bbeta^*, \sigma^2}) \ge \mathbb{E}_{(\bx,y)\sim\mu_*} \left[|y-\bbeta^\top \bx|^2\right] &= \mathbb{E}_{(\bx,y)\sim \mu_*} \left[\left(y - (\bbeta^*)^\top \bx +  (\bbeta^*)^\top \bx - \bbeta^\top \bx\right)^2\right] \\
    &= \mathbb{E}_{(\bx,y)\sim \mu_*} \left[\left(y - (\bbeta^*)^\top \bx\right)^2\right] + \sigma^2 \|\bbeta - \bbeta^*\|_2^2 \\
    &= \mathsf{R}_{\mathtt{oos}}(\bbeta^*;\mathcal{U}_{\bbeta^*,\sigma^2}) + \sigma^2 \|\bbeta - \bbeta^*\|_2^2,
\end{align*} which completes the proof. \qed
%\end{proof}

\subsection{Proof of \cref{prop:necessary-multiple-envs}}
\label{prop:necessary-multiple-envs-proof}
    Denote $s^* = |S^*|$. Without loss of generality, let $S^* = \{1,\cdots, s\}$ and $\beta_{s+1} \neq 0$. Consider the following matrix
    \begin{align*}
        \bar{\bSigma} = \begin{bmatrix}
            \bI_{s^*} &\bB &\bbeta^*_{S^*} \\
            \bB^\top &\bI_{p-s^*} + \bB^\top \bB & \balpha \\
            (\bbeta^*_{S^*})^\top & \balpha^\top & c\\
        \end{bmatrix}
    \end{align*} where $\bI_{q}$ is $q$ by $q$ identity matrix, $\bB \in \mathbb{R}^{s^*\times (p-s^*)}$ satisfies $B_{j,k} = 1\{k=1\} (\beta_j^* - \beta_j)/(\beta_k)$, $\balpha = \bB^\top \bbeta_{S^*} + (\bI_{p-s^*} + \bB^\top \bB) \bbeta_{T}$, $c=1+\|\balpha\|_2^2 + \|\bbeta_{S^*}^*\|_2^2$. It is easy to verify that our constructed $\bar{\bSigma}$ satisfies $\bar{\bSigma} \succeq \bI_{p+1}$. We then argue that the distribution $\mu$ that $(\bx,y) \sim \mathcal{N}(0,\bar{\bSigma})$ satisfies $\mu \in \mathcal{U}_{\bbeta^*,\sigma^2}$ for some large enough $\sigma^2$ and $\mathbb{E}_{\mu}[y|\bx] = \bbeta^\top \bx$. To this end, we only need to verify that $\mathbb{E}[y|\bx_{S^*}] = (\bbeta^*)^\top \bx$ and $\mathbb{E}[y|\bx] = \bbeta^\top \bx$. For the first condition, our construction of $\bar{\bSigma}$ and the conditional distribution of multivariate Gaussian ensures $\mathbb{E}[y|\bx_{S^*}] = (\bI_{s^*})^{-1} (\bbeta_{S^*}^*)^\top \bx_S^* = (\bbeta_{S^*}^*)^\top \bx_S^*$. Similarly, in order to show $\mathbb{E}[y|\bx] = \bbeta^\top \bx$, it suffices to verify that
    \begin{align*}
        \begin{bmatrix}
            \bI_{s^*} &\bB \\
            \bB^\top &\bI_{p-s^*} + \bB^\top \bB
        \end{bmatrix} \bbeta = \begin{bmatrix}
            \bbeta^*_{S^*} \\
            \balpha \\
        \end{bmatrix},
    \end{align*} which can be validated by plugging in our construction of $\bB$ and $\balpha$. This completes the proof. \qed

\subsection{Proof of \cref{thm:global-strong-convexity}}

It follows from the definition of $\mathsf{Q}(\bbeta; \gamma, \bomega)$ and \cref{cond0-model} that
\begin{align*}
	\mathsf{Q}(\bbeta; \gamma,\bomega) - \mathsf{Q}(\bbeta^*;\gamma, \bomega) = \mathsf{R}(\bbeta;\bomega) - \mathsf{R}(\bbeta^*; \bomega) + \gamma \mathsf{J}(\bbeta; \bomega) = \mathsf{T}_1(\bbeta) + \gamma \mathsf{T}_2(\bbeta).
\end{align*}
Denote $\bDelta = \bbeta - \bbeta^*$, $S=\supp(\bbeta)$, $T=S^*\setminus S$, and $\bar{S} = S\cup S^*$. it follows from the model $y^{(e)} = (\bbeta^*)^\top \bx^{(e)} + \varepsilon^{(e)}$ that
\begin{align*}
	\mathsf{T}_1(\bbeta) &= \sum_{e\in \mathcal{E}} \omega^{(e)} \left(\mathsf{R}^{(e)}(\bbeta) - \mathsf{R}^{(e)}(\bbeta^*)\right) \\
	&= \sum_{e\in \mathcal{E}} \omega^{(e)} \left(\mathbb{E} \left[|y^{(e)} - \bbeta^\top \bx^{(e)}|^2 \right]- \mathbb{E} \left[|y^{(e)} - (\bbeta^*)^\top \bx^{(e)}|^2 \right]\right) \\
	&= \sum_{e\in \mathcal{E}} \omega^{(e)} \left(\mathbb{E} \left[|(\bbeta^* - \bbeta)^\top \bx^{(e)} + \varepsilon^{(e)} |^2\right] - \mathbb{E}[|\varepsilon^{(e)}|^2 ]\right) \\
	&= \sum_{e\in \mathcal{E}} \omega^{(e)} \bDelta^\top \bSigma^{(e)} \bDelta - 2\bDelta^\top \mathbb{E}\left[\varepsilon^{(e)} \bx^{(e)} \right] \\
	&= \bDelta^\top \Bigg(\sum_{e\in \mathcal{E}} \omega^{(e)} \bSigma^{(e)} \Bigg) \bDelta - 2\bDelta^\top \Bigg(\sum_{e\in \mathcal{E}} \omega^{(e)} \mathbb{E}[\varepsilon^{(e)} \bx^{(e)}] \Bigg) = \bDelta_{\bar{S}}^\top \bar{\bSigma}_{\bar{S}} \bDelta_{\bar{S}} - 2 \bDelta^\top_{\bar{S}} \bar{\bu}_{\bar{S}}
\end{align*} 
It follows from the fact $\bar{\bu}_T=\bm{0}$ and Cauchy-Schwarz inequality that
\begin{align}
\label{eq:proof-thm1-eq1}
	2\bDelta^\top_{\bar{S}} \bar{\bu}_{\bar{S}} = 2\left(\sqrt{\epsilon} \bDelta_{\bar{S}} \bar{\bSigma}_{\bar{S}}^{1/2} \right) \left(\epsilon^{-1/2} \bar{\bSigma}_{\bar{S}}^{-1/2} \bar{\bu}_{\bar{S}} \right) \le \epsilon \bDelta_{\bar{S}}^\top \bar{\bSigma}_{\bar{S}} \bDelta_{\bar{S}} + \epsilon^{-1} \bar{\bu}_{\bar{S}}^\top \bar{\bSigma}_{\bar{S}}^{-1} \bar{\bu}_{\bar{S}}
\end{align}
Plugging \eqref{eq:proof-thm1-eq1} back yields
\begin{align}
\label{eq:proof-thm1-t1-lb}
	\mathsf{T}_1(\bbeta) \ge (1-\epsilon) \bDelta_{\bar{S}}^\top \bar{\bSigma}_{\bar{S}} \bDelta_{\bar{S}} - \epsilon^{-1} \bar{\bu}_{\bar{S}}^\top \bar{\bSigma}_{\bar{S}}^{-1} \bar{\bu}_{\bar{S}}.
\end{align}
At the same time, we also have
\begin{align*}
	\mathsf{T}_2(\bbeta) &= \sum_{e\in \mathcal{E}} \omega^{(e)} \left\|\mathbb{E} \left[(y^{(e)} - \bbeta^\top \bx^{(e)}) \bx_S^{(e)}\right]\right\|_2^2 \\
	&= \sum_{e\in \mathcal{E}} \omega^{(e)} \left\| \mathbb{E} \left[\varepsilon^{(e)} \bx_{\bar{S}}^{(e)} \right] + (\bbeta^*_S - \bbeta_S)^\top \mathbb{E} \left[\bx_S^{(e)} (\bx_S^{(e)})^\top\right] + (\bbeta_T^*)^\top \mathbb{E} \left[\bx_T^{(e)} (\bx_S^{(e)})^\top\right] \right\|^2 \\
	&= \sum_{e\in \mathcal{E}} \omega^{(e)} \left\| \bar{\bu}_{S} + \bSigma^{(e)}_{S,T} \bbeta_T^* - \bSigma^{(e)}_S \bDelta_S\right\|_2^2 \\
	&= \sum_{e\in \mathcal{E}} \omega^{(e)} \left\| \bSigma_S(\bbeta^{(e, S)} - \bbeta^* - \bDelta_S) \right\|^2.
\end{align*} Putting these pieces together with \cref{cond1-well-condition}, we have
\begin{align}
	\mathsf{T}_1(\bbeta) + \gamma \mathsf{T}_2(\bbeta) & \ge (1-\epsilon)  \|\bar{\bSigma}^{1/2}\bDelta_{\bar{S}}\|_2^2 - \varepsilon^{-1} \kappa_L^{-1} \|\bar{\bu}_{S}\|_2^2 + \gamma \sum_{e\in \mathcal{E}} \omega^{(e)} \kappa_L^2 \| \bbeta^{(e,S)} - \bbeta^* - \bDelta_S\|^2_2 \nonumber \\
	&\ge (1-\epsilon) \|\bar{\bSigma}^{1/2}\bDelta_{\bar{S}}\|_2^2 - \varepsilon^{-1} \kappa_L^{-1} \mathsf{b}_{S} + \inf_{\bv \in \mathbb{R}^{|S|}}\gamma \sum_{e\in \mathcal{E}} \omega^{(e)} \kappa_L^2 \| \bbeta^{(e,S)}_S - \bbeta^*_S - \bv\|^2_2 \label{eq:proof-thm1-mixed-lb}
\end{align}
Define 
\begin{align*}
	L(\bv) = \sum_{e\in \mathcal{E}} \omega^{(e)} \| \bbeta^{(e,S)}_S - \bbeta^*_S - \bv\|^2_2.
\end{align*} It is obvious that $L(\bv)$ is of a quadratic form and is strong convex with curvature $1$. So it has the unique global minimizer $\bv^*=\sum_{e\in \mathcal{E}}\omega^{(e)} \bbeta^{(e,S)} -\bbeta_S^*$ with minimum value
\begin{align}
\label{eq:j-to-bs}
	\inf_{\bv \in \mathbb{R}^{|S|}} L(\bv) = \sum_{e\in \mathcal{E}} \omega^{(e)} \|\bbeta_S^{(e,S)} - \bbeta_S^* - \bv^*\|_2^2 = \bar{\mathsf{d}}_S.
\end{align} Substituting it back into \eqref{eq:proof-thm1-mixed-lb} gives
\begin{align}
\label{eq:thm1-proof-decomp}
	\mathsf{T}_1(\bbeta) + \gamma \mathsf{T}_2(\bbeta) &\ge (1 - \epsilon) \|\bar{\bSigma}^{1/2}\bDelta\|_2^2 - \epsilon^{-1} \kappa_L^{-1} \mathsf{b}_S + \gamma \kappa_L^2 \mathtt{d}_{S} \\
	&= (1-\epsilon) \|\bar{\bSigma}^{1/2}\bDelta\|_2^2 - \epsilon^{-1} \kappa_L^{-1} \mathsf{b}_S + \epsilon^{-1} \kappa_L^2 \bar{\mathsf{d}}_S \gamma^* +  (\gamma - \epsilon^{-1} \gamma^*) \kappa_L^2 \bar{\mathsf{d}}_S \nonumber \\
	&\ge (1-\epsilon) \|\bar{\bSigma}^{1/2}\bDelta\|_2^2 + (\gamma - \epsilon^{-1} \gamma^*) \kappa_L^2 \bar{\mathsf{d}}_S. \nonumber
\end{align} This completes the proof. \qed

\subsection{Proof of \cref{prop:ex-h-var}}

We first calculate serval quantities of interests and determine the original curvature $\kappa_L$.

It follows from the model \eqref{eq:model-linear-scm-varying-v0} that
\begin{align*}
\bx^{(e)} = \bB \bx^{(e)} + 	\balpha (\bbeta^*)^\top \bx^{(e)} + \balpha \varepsilon^{(e)}_0 + \bvarepsilon_x.
\end{align*} Because the induced graph is acyclic, then there exists a permutation $\pi: [p] \to [p]$ such that
\begin{align*}
	\left[\bB + \balpha (\bbeta^*)^\top \right]_{i,j} = 0 ~~~ \text{if} ~~~ \pi(i) < \pi(j),
\end{align*} that is, there exists a permutation matrix $\bP$ such that the matrix 
\begin{align*}
	\tilde{\bV} = \bP [\bB + \balpha (\bbeta^*)^\top] \bP^\top
\end{align*} satisfies $\tilde{\bV}$ is a lower triangular matrix with all zeros on the diagonal. Hence, the inverse of the matrix $\bP \bI \bP^\top - \tilde{\bV}$ exists and is an upper triangular matrix with all ones on the diagonal, thus implying
\begin{align*}
	\bW = (\bI - \bB + \balpha (\bbeta^*)^\top)^{-1} ~~\text{exists with}~~ \nu_{\min}(\bW) \ge 1.
\end{align*} Therefore, we can represent the covariate $\bx^{(e)}$ in terms of all the exogenous variables $\bvarepsilon_x$ and $\varepsilon_0^{(e)}$ as
\begin{align}
\label{eq:proof-ex2-eq-x}
	\bx^{(e)} = \bW (\bvarepsilon_x + \balpha \varepsilon_0^{(e)}).
\end{align} Therefore, denote $\bu = \bW \balpha$, we further have
\begin{align}
\bSigma^{(e)} &= \bW \bD \bW^\top + v_0^{(e)} \bu \bu^\top,\\
\mathbb{E}\left[\varepsilon^{(e)} \bx^{(e)}\right] &= v_0^{(e)} \bu^{(e)}
\end{align}
Given these quantities, we can write down the bias for any $S$ as
\begin{align}
\label{eq:proof-ex2-b}
	\mathsf{b}_S = \left\|\frac{1}{2} \left(\mathbb{E}[\varepsilon^{(e)} \bx_S^{(e)}]\right) (v_0^{(1)} + v_0^{(2)}) \right\|_2^2 = \Bigg(\frac{v_0^{(1)} + v_0^{(2)}}{2}\Bigg)^2 \cdot \|\bu_S\|_2^2.
\end{align}
The calculation of bias-difference term $\bar{\mathsf{d}}_S$ is more involved. Denote $T=S^*\setminus S$, then one has
\begin{align*}
	\bbeta^{(e,S)}_S - \bbeta^*_S = (\bSigma_S^{(e)})^{-1} \mathbb{E}[\varepsilon^{(e)} \bx_S^{(e)}] + (\bSigma_S^{(e)})^{-1} \bSigma_{S,T}^{(e)} \bbeta_T^* = \mathsf{T}_1^{(e)} + \mathsf{T}_2^{(e)}.
\end{align*} Denoting $\bM_S = \bW_{S,:} \bD \bW_{S,:}^\top$, it follows from Sherman-Morrison formula that
\begin{align*}
	\mathsf{T}_1^{(e)} &= (\bM_S + v_0^{(e)} \bu_S \bu_S)^{-1} \bu_S v_0^{(e)}\\
	&= \left\{\bM_S^{-1} - \frac{\bM_S^{-1} \bu_S \bu_S^\top \bM_S^{-1} v_0^{(e)}}{1+\bu_S \bM_S^{-1} \bu_S v_0^{(e)}}\right\} \bu_S v_0^{(e)} \\
	&= \frac{v_0^{(e)}}{1 + \bu_S \bM_S^{-1} \bu_S v_0^{(e)}} \bM_S^{-1} \bu_S.
\end{align*}
Observe that $\bu_T = 0$. This yields
\begin{align*}
	\mathsf{T}_2^{(e)} &= \left(\bM_S + v_0^{(e)} \bu_S \bu_S^\top \right)^{-1} (\bW_{S,:} \bD \bW_{T,:}^\top) \bbeta_T^* \\
	&= \bM_S^{-1} (\bW_{S,:} \bD \bW_{T,:}^\top) \bbeta_T^* - \frac{\bM_S^{-1} \bu_S^{(e)}}{1+v_0^{(e)} \bu_S^\top\bM_S^{-1} \bu_S} v_0^{(e)} \bu_S^\top \bM_S^{-1} (\bW_{S,:} \bD \bW_{T,:}^\top) \bbeta_T^*
\end{align*}
Putting these pieces together, denote $\iota_S = \bu_S \bM_S^{-1} \bu_S$, we have 
\begin{align}
	\bar{\mathsf{d}}_S &= \sum_{e\in \{1,2\}} \frac{1}{2} \left\|\bbeta^{(e,S)}- \frac{\bbeta^{(1,S)} + \bbeta^{(2,S)}}{2}\right\|_2^2 \nonumber\\
				 &= \frac{1}{2} \left\|\bbeta^{(1,S)} - \bbeta^{(2,S)}\right\|_2^2 \nonumber\\
				 &= \frac{1}{2} \|\bM_S^{-1} \bu_S\|_2^2 \left|1 - \bu_S^\top \bM_S^{-1} (\bW_{S,:} \bD \bW_{T,:}^\top) \bbeta_T^*\right|^2 \left|\frac{v_0^{(1)}}{1+\iota_S v_0^{(1)}} - \frac{v_0^{(2)}}{1+\iota_S v_0^{(2)}} \right|^2 \nonumber\\
				 &\ge \frac{1}{2} c_1^{-2} \|\bu_S\|_2^2 \frac{|v_0^{(1)} - v_0^{(2)}|^2 \left|1 - \bu_S^\top \bM_S^{-1} (\bW_{S,:} \bD \bW_{T,:}^\top) \bbeta_T^*\right|^2 }{(1 + \iota_S v_0^{(1)})^2(1 + \iota_S v_0^{(2)})^2} \nonumber\\
				 &\ge \frac{C^{-1}}{2} \|\bu_S\|_2^2 \frac{|v_0^{(1)} - v_0^{(2)}|^2 \left|1 - \bu_S^\top \bM_S^{-1} (\bW_{S,:} \bD \bW_{T,:}^\top) \bbeta_T^*\right|^2 }{(1 + v_0^{(1)})^2(1 + v_0^{(2)})^2} \left\{\min_{i} D_{i,i} \right\}^{4} \label{eq:proof-ex2-lb-d}
\end{align} where the last inequality follows from the fact
\begin{align*}
	\iota_S \le \|\bM_S\|^{-1} \|\bu_S\|^2 \le \{\nu_{\min}(\bW)\}^2 \left\{\min_{i} D_{i,i} \right\}^{-1} \|\bW\|_2^2 \|\balpha\|_2^2 \lesssim \left\{\min_{i} D_{i,i} \right\}^{-1}.
\end{align*} Plugging the above quantities \eqref{eq:proof-ex2-b} and \eqref{eq:proof-ex2-lb-d} into \eqref{eq:thm1-gamma-star} completes the proof.


























































\section{Proofs for Non-asymptotic Results}
\label{sec:proof-nonasymptotic}

\subsection{Preliminaries}

We first define some concepts that will be used throughout the proof.

\begin{definition}[Sub-Gaussian Random Variable]
A random variable $X$ is a sub-Gaussian random variable with parameter $\sigma \in \mathbb{R}^+$ if 
\begin{align*}
    \forall \lambda \in \mathbb{R}, ~~~~~~ \mathbb{E} \left[\exp(\lambda X)\right] \le \exp\left(\frac{\lambda^2}{2}\sigma^2\right)
\end{align*}
\end{definition}

\begin{definition}[Sub-Exponential Random Variable]
    A random variable $X$ is a sub-Exponential random variable with parameter $(\nu,\alpha) \in \mathbb{R}^+ \times \mathbb{R}^+$ if
    \begin{align*}
        \forall |\lambda|<1/\alpha, ~~~~~~ \mathbb{E}\left[\exp(\lambda X)\right] \le \exp\left(\frac{\lambda^2}{2} \nu^2\right).
    \end{align*}
\end{definition}

It is easy to verify that the product of two sub-Gaussian random variables is a Sub-Exponential random variable, and the dependence of the parameters can be written as follows.

\begin{lemma}[Product of Two Sub-Gaussian Random Variables]
\label{lemma:product-sub-gaussian}
    Suppose $X_1$ and $X_2$ are two zero-mean sub-Gaussian random variables with parameters $\sigma_1$ and $\sigma_2$, respectively. Then $X_1 X_2$ is a sub-exponential random variable with parameter $(c_1 \sigma_1\sigma_2, c_2 \sigma_1\sigma_2)$, where $c_1,c_2>0$ are some universal constants.
\end{lemma}


We also have the following lemma stating the concentration inequality for the sum of independent sub-exponential random variables.
\begin{lemma}[Sum of Independent Sub-exponential Random Variables]
\label{lemma:sum-sub-exp}
Suppose $X_1,\ldots, X_N$ are independent sub-exponential random variables with parameters $\{(\nu_i, \alpha_i)\}_{i=1}^N$, respectively. There exists some universal constant $c_1$ such that the following holds,
\begin{align*}
    \mathbb{P}\left[\left|\sum_{i=1}^N (X_i - \mathbb{E}[X_i]) \right| \ge c_1 \left\{\sqrt{t\times \sum_{i=1}^N \nu_i^2} + t\times \max_{i\in [N]} \alpha_i \right\}\right] \le 2e^{-t}.
\end{align*}
\end{lemma}

We also define the following quantity regarding sample size
\begin{align}
\label{eq:def-n-omega}
n_{\bomega} = \Bigg(\sum_{e\in \mathcal{E}} \frac{(\omega^{(e)})^2}{n^{(e)}}\Bigg)^{-1}.
\end{align} It is easy to see that $n_{\bomega} \ge n_*$.

We also define
\begin{align}
\mathcal{B}(S) = \left\{\bbeta \in \mathbb{R}^p, \supp(\bbeta) = S, \|\bbeta\|_2=1\right\} \qquad \text{and} \qquad \mathcal{B}_s = \bigcup_{S\subseteq [p], |S| \le s} \mathcal{B}(S). \label{eq:ball}
\end{align}

\subsection{Proof of \cref{prop:pooled-least-squares}}

It follows from the definition of $\mathsf{R}(\bbeta)$ and $\mathsf{R}^{(e)}(\bbeta)$ that
\begin{align*}
    \mathsf{R}(\bbeta) &= \sum_{e\in \mathcal{E}} \omega^{(e)} \mathsf{R}^{(e)}(\bbeta) \\
    &= \sum_{e\in \mathcal{E}} \omega^{(e)} \left\{(\bbeta - \bbeta^*)^\top \bSigma^{(e)} (\bbeta - \bbeta^*) - 2(\bbeta - \bbeta^*)^\top \mathbb{E}[\bx^{(e)} \varepsilon^{(e)}] + \mathbb{E} |\varepsilon^{(e)}|^2 \right\}\\
    &= (\bbeta - \bbeta^*)^\top \bar{\bSigma} (\bbeta - \bbeta^*) - 2(\bbeta - \bbeta^*)^\top \sum_{e\in \mathcal{E}} \omega^{(e)} \mathbb{E} [\bx^{(e)} \varepsilon^{(e)}] + \sum_{e\in \mathcal{E}} \omega^{(e)} \mathbb{E} |\varepsilon^{(e)}|^2.
\end{align*}
We can see $\mathsf{R}(\bbeta)$ is of quadratic form of $\bbeta - \bbeta^*$. Since $\lambda_{\min}(\bSigma) > 0$, the minimizer of $\mathsf{R}(\bbeta)$ is unique and is
\begin{align}
\label{eq:beta-bar-r}
    \bar{\bbeta}^{\mathsf{R}} = \bbeta^* + (\bar{\bSigma})^{-1} \sum_{e\in \mathcal{E}} \omega^{(e)} \mathbb{E}[\bx^{(e)} \varepsilon^{(e)}].
\end{align}
Combining the above definition of $\bar{\bbeta}^{\mathsf{R}}$ together with \cref{cond1-well-condition} validates the upper bound and lower bound on $\|\bar{\bbeta}^{\mathsf{R}}-\bbeta^*\|_2$. Moreover, plugging in above \eqref{eq:beta-bar-r} into $\mathsf{R}(\bbeta)$ following by some calculations gives
\begin{align*}
    \mathsf{R}(\bbeta) - \mathsf{R}(\bar{\bbeta}^{\mathsf{R}}) = (\bbeta - \bbeta^*)^\top \bar{\bSigma} (\bbeta - \bbeta^*),
\end{align*} which completes the proof of claim \eqref{eq:pooled-least-square-strong-convexity}.

Moreover, for any $\bbeta \in \mathbb{R}^p$, we have
\begin{align*}
    \|\bar{\bSigma}^{1/2}(\bbeta - \bar{\bbeta}^{\mathsf{R}})\|_2^2 &= \mathsf{R}(\bbeta) - \mathsf{R}(\bar{\bbeta}^{\mathsf{R}}) \\
    &= \mathsf{R}(\bbeta) - \hat{\mathsf{R}}(\bbeta) + \hat{\mathsf{R}}(\bbeta) - \hat{\mathsf{R}}(\bar{\bbeta}^{\mathsf{R}}) + \hat{\mathsf{R}}(\bar{\bbeta}^{\mathsf{R}}) - \mathsf{R}(\bar{\bbeta}^{\mathsf{R}}).
\end{align*}

We argue that
\begin{align}
\label{eq:pooled-least-squares-high-prob-bound1}
\begin{split}
    \mathbb{P}[\mathcal{C}_{1,t}] = \mathbb{P} \Bigg\{ \forall \bbeta \in \mathbb{R}^p, ~~& \mathsf{R}(\bbeta) - \mathsf{R}(\bar{\bbeta}^{\mathsf{R}}) - \left(\hat{\mathsf{R}}(\bbeta) - \hat{\mathsf{R}}(\bar{\bbeta}^{\mathsf{R}}) \right) \\
    &~~~~~~ \le C \sigma_x \left(\sigma_\varepsilon + \sigma_x\|\bar{\bSigma}^{1/2}(\bbeta^* - \bar{\bbeta}^{\mathsf{R}})\|_2 \right) \sqrt{\frac{p+t}{n_*}} \|\bar{\bSigma}^{1/2}(\bbeta - \bar{\bbeta}^{\mathsf{R}})\|_2 \\
    &~~~~~~~~~~~~ + C \sigma_x^2 \sqrt{\frac{p + t}{n_*}} \|\bar{\bSigma}^{1/2}(\bbeta - \bar{\bbeta}^{\mathsf{R}})\|_2^2\Bigg\}\ge 1-2e^{-t}
\end{split}
\end{align} for any $t \in (0, n_* - p]$, which will be validated later. If such a claim holds, then under $\mathcal{C}_{1,t}$ which occurs with probability at least $1-2e^{-t}$, we have that, for any $\bbeta \in \mathbb{R}^p$
\begin{align*}
        \|\bar{\bSigma}^{1/2}(\bbeta - \bar{\bbeta}^{\mathsf{R}})\|_2^2 & \le \hat{\mathsf{R}}(\bbeta) - \hat{\mathsf{R}}(\bar{\bbeta}^{\mathsf{R}}) + C \sigma_x^2 \sqrt{\frac{p + t}{n_*}} \|\bar{\bSigma}^{1/2}(\bbeta - \bar{\bbeta}^{\mathsf{R}})\|_2^2 \\
        &~~~~~~ + C \sigma_x \left(\sigma_\varepsilon + \sigma_x\|\bar{\bSigma}^{1/2}(\bbeta^* - \bar{\bbeta}^{\mathsf{R}})\|_2 \right) \sqrt{\frac{p+t}{n_*}} \|\bar{\bSigma}^{1/2}(\bbeta - \bar{\bbeta}^{\mathsf{R}})\|_2.
\end{align*} Suppose further that $n^* \ge (p + t) 4C^2\sigma_x^4$, we can apply the above inequality with $\hat{\bbeta}_\mathsf{R}$ which satisfies $\hat{\mathsf{R}}(\hat{\bbeta}_\mathsf{R}) - \hat{\mathsf{R}}(\bar{\bbeta}^{\mathsf{R}})\le 0$ and obtain
\begin{align*}
\|\bar{\bSigma}^{1/2}(\hat{\bbeta}_\mathsf{R} - \bar{\bbeta}^{\mathsf{R}})\|_2^2 \le 2C \sigma_x \left(\sigma_\varepsilon + \sigma_x\|\bar{\bSigma}^{1/2}(\bbeta^* - \bar{\bbeta}^{\mathsf{R}})\|_2 \right) \sqrt{\frac{p+t}{n_*}} \|\bar{\bSigma}^{1/2}(\hat{\bbeta}_\mathsf{R} - \bar{\bbeta}^{\mathsf{R}})\|_2.
\end{align*} This completes the proof.

\noindent \emph{Proof of the Bound \eqref{eq:pooled-least-squares-high-prob-bound1}.} It follows from the definition of $\hat{\mathsf{R}}$ and the above identity that \begin{align*}
    &\mathsf{R}(\bbeta) - \mathsf{R}(\bar{\bbeta}^{\mathsf{R}}) - \left(\hat{\mathsf{R}}(\bbeta) - \hat{\mathsf{R}}(\bar{\bbeta}^{\mathsf{R}}) \right) \\
    &~~~~~~ = (\bbeta - \bar{\bbeta}^{\mathsf{R}})^\top \bar{\bSigma} (\bbeta - \bar{\bbeta}^{\mathsf{R}}) - (\bbeta - \bar{\bbeta}^{\mathsf{R}})^\top \Bigg(\sum_{e\in \mathcal{E}} \omega^{(e)}\hat{\bSigma}^{(e)}\Bigg) (\bbeta - \bar{\bbeta}^{\mathsf{R}}) \\
    &~~~~~~~~~~~~ - 2(\bbeta - \bar{\bbeta}^{\mathsf{R}})^\top \sum_{e\in \mathcal{E}} \omega^{(e)} \hat{\mathbb{E}} \left[\bx^{(e)} \left((\bar{\bbeta}^{\mathsf{R}})^\top \bx^{(e)} - y^{(e)}\right)\right].
\end{align*}
When $p + t \le n_*$, define the following event
\begin{align*}
    \mathcal{K}_{1,t} &= \left\{ \left\|\bI - \bar{\bSigma}^{-1/2} \sum_{e\in \mathcal{E}} \omega^{(e)}\hat{\bSigma}^{(e)} \bar{\bSigma}^{-1/2}\right\|_2 \le C_1 \sigma_x^2 \sqrt{\frac{p+t}{n_*}}\right\} \\
    \mathcal{K}_{2,t} &= \Bigg\{ \left\|\sum_{e\in \mathcal{E}} \omega^{(e)} \bar{\bSigma}^{-1/2} \hat{\mathbb{E}} \left[\bx^{(e)} \left((\bar{\bbeta}^{\mathsf{R}})^\top \bx^{(e)} - y^{(e)}\right)\right]\right\| \\    
    &~~~~~~~~~~~~~~~~~~~~~~~~~~\le C_2 \sigma_x \left(\sigma_\varepsilon + \sigma_x\|\bar{\bSigma}^{1/2}(\bbeta^* - \bar{\bbeta}^{\mathsf{R}})\|_2 \right) \sqrt{\frac{p+t}{n_*}}\Bigg\}
\end{align*} for some universal constant $C_1$--$C_2$ to be determined. It suffices to prove that $\mathbb{P}[\mathcal{K}_{1,t}] \land \mathbb{P}[\mathcal{K}_{2,t}] \ge 1-e^{-t}$ for any $t>0$. If the two claims are verified, then we can conclude that under the event $\mathcal{K}_{1,t} \cap \mathcal{K}_{2,t}$ that occurs with probability at least $1-2e^{-t}$, the following holds
\begin{align*}
&\mathsf{R}(\bbeta) - \mathsf{R}(\bar{\bbeta}^{\mathsf{R}}) - \left(\hat{\mathsf{R}}(\bbeta) - \hat{\mathsf{R}}(\bar{\bbeta}^{\mathsf{R}}) \right) \\
    &~~~~~~ = \left\|\bar{\bSigma}^{1/2}(\bbeta - \bar{\bbeta}^{\mathsf{R}})\right\|_2 \left\|\bI - \bar{\bSigma}^{-1/2} \sum_{e\in \mathcal{E}} \omega^{(e)}\hat{\bSigma}^{(e)} \bar{\bSigma}^{-1/2} \right\|_2 \left\|\bar{\bSigma}^{1/2}(\bbeta - \bar{\bbeta}^{\mathsf{R}})\right\|_2 \\
    &~~~~~~~~~~~~ + 2\left\|\bar{\bSigma}^{1/2}(\bbeta - \bar{\bbeta}^{\mathsf{R}})\right\|_2 \left\|\sum_{e\in \mathcal{E}} \omega^{(e)} \bar{\bSigma}^{-1/2} \hat{\mathbb{E}} \left[\bx^{(e)} \left((\bar{\bbeta}^{\mathsf{R}})^\top \bx^{(e)} - y^{(e)}\right)\right]\right\|_2 \\
    &~~~~~~ \le C' \sigma_x^2 \sqrt{\frac{p+t}{n_*}} \|\bar{\bSigma}^{1/2}(\bbeta - \bar{\bbeta}^{\mathsf{R}})\|_2^2  \\
    &~~~~~~~~~~~~~~~ + C'\sigma_x \left(\sigma_\varepsilon + \sigma_x\|\bar{\bSigma}^{1/2}(\bbeta^* - \bar{\bbeta}^{\mathsf{R}})\|_2 \right) \sqrt{\frac{p+t}{n_*}} \|\bar{\bSigma}^{1/2}(\bbeta - \bar{\bbeta}^{\mathsf{R}})\|_2,
\end{align*} which completes the proof of the claim \eqref{eq:pooled-least-squares-high-prob-bound1}.

\noindent {\sc Step 1. High Probability Bound for $\mathcal{K}_{1,t}$. } Let $\bA = \bI - \bar{\bSigma}^{-1/2} \sum_{e\in \mathcal{E}} \omega^{(e)}\hat{\bSigma}^{(e)} \bar{\bSigma}^{-1/2}$, similar to the derivation in \eqref{eq:step3-sup-argument-2}, we can construct $N$ pairs of $p$-dimensional unit vectors $(\bv_1, \bu_1),\ldots, (\bv_N, \bu_N)$ with $N\le 8100^p$ such that
\begin{align*}
    \|\bA\|_2 \le 4\sup_{k\in [N]} \bv_k^\top \bA \bu_k.
\end{align*} For fixed $(\bv, \bu)$, using the identity
\begin{align*}
    \bI = \sum_{e\in \mathcal{E}} \omega^{(e)} \bar{\bSigma}^{-1/2} \bSigma^{(e)} \bar{\bSigma}^{-1/2},
\end{align*} we obtain
\begin{align*}
    \bv^\top \bA \bu &= \sum_{e\in \mathcal{E}} \sum_{i=1}^{n^{(e)}} \frac{\omega^{(e)}}{n^{(e)}} \left\{\bu^\top \bar{\bSigma}^{-1/2} \bSigma^{(e)} \bar{\bSigma}^{-1/2} \bv - \bu^\top \bar{\bSigma}^{-1/2} \bx^{(e)}_i (\bx^{(e)}_i)^\top \bar{\bSigma}^{-1/2} \bv \right\} \\ 
    &= \sum_{e\in \mathcal{E}} \sum_{i=1}^{n^{(e)}} \frac{\omega^{(e)}}{n^{(e)}} \left(\mathbb{E}[X_{e,i}] - X_{e,i}\right)
\end{align*} for $X_{e,i} = \bu^\top \bar{\bSigma}^{-1/2} \bx^{(e)}_i (\bx^{(e)}_i)^\top \bar{\bSigma}^{-1/2} \bv$. It follows from \cref{cond2-subgaussain-x} that $X_{e,i}$ is the product of two zero-mean sub-Gaussian random variables with parameter $\sigma_x$. Then it follows from \cref{lemma:product-sub-gaussian} and \cref{lemma:sum-sub-exp} that the following holds
\begin{align*}
    \mathbb{P}\left[|\bv_k^\top \bA \bu_k| \ge C_3 \sigma_x \left(\sqrt{\frac{u}{n_{\bomega}}} + \frac{u}{n_*}\right)\right] \le 2e^{-u}
\end{align*} for all the $u>0$, where $C_3$ is some universal constant. Applying union bounds further gives
\begin{align*}
    \mathbb{P}\left[\sup_{k\in [N]} |\bv_k^\top \bA \bu_k| \ge C_3 \sigma_x \left(\sqrt{\frac{u}{n_{\bomega}}} + \frac{u}{n_*}\right)\right] \le 2Ne^{-u}.
\end{align*} So it concludes the proof by setting $u=t+\log(2N) = t+C_4 p$ with some $C_4 > 1$ and observing that
\begin{align*}
    \sqrt{\frac{t + C_4 p}{n_{\bomega}}} + \frac{t + C_4 p}{n_*} \le 2C_4 \sqrt{\frac{t + p}{n_*}}
\end{align*} provided $p + t\le n_*$.

\noindent {\sc Step 2. High Probability Bound for $\mathcal{K}_{2,t}$. } 

Let $\bb = \sum_{e\in \mathcal{E}} \omega^{(e)} \bar{\bSigma}^{-1/2} \hat{\mathbb{E}} \left[\bx^{(e)} \left((\bar{\bbeta}^{\mathsf{R}})^\top \bx^{(e)} - y^{(e)}\right)\right]$, the proof is similar to the first part of {\sc Step 5} in the proof of \cref{lemma:instance-dependent-bound-j-lowdim}. To be specific, there exist $N=90^p$ $p$-dimensional unit vectors $\bu_1,\ldots, \bu_N$ such that $\|\bb\|_2 \le 2\sup_{k\in [N]} \bu_k^\top \bb$. The key here is to decompose $\bu^\top \bb$ as the sum of zero-mean independent random variables. Denote
\begin{align*}
    X_{e,i} = \bu^\top \bar{\bSigma}^{-1/2} \bx^{(e)} \left((\bx^{(e)})^\top (\bar{\bbeta}^{\mathsf{R}} - \bbeta^*) - \varepsilon^{(e)}\right).
\end{align*} We have
\begin{align*}
    \bu^\top \bb = \sum_{e\in \mathcal{E}} \sum_{i=1}^{n^{(e)}} \frac{\omega^{(e)}}{n^{(e)}} X_{e,i} &= \sum_{e\in \mathcal{E}} \sum_{i=1}^{n^{(e)}} \frac{\omega^{(e)}}{n^{(e)}} (X_{e,i} - \mathbb{E}[X_{e,i}]) + \sum_{e\in \mathcal{E}} \omega^{(e)} \mathbb{E}[X_{e,1}] \\ &= \sum_{e\in \mathcal{E}} \sum_{i=1}^{n^{(e)}} \frac{\omega^{(e)}}{n^{(e)}} (X_{e,i} - \mathbb{E}[X_{e,i}])
\end{align*} where the last equality follows from the definition of $\bar{\bbeta}^{\mathsf{R}}$ \eqref{eq:beta-bar-r} that
\begin{align*}
    \sum_{e\in \mathcal{E}} \omega^{(e)} \mathbb{E}[X_{e,1}] &= \bu^\top \bar{\bSigma}^{-1/2} \bar{\bSigma} (\bar{\bbeta}^{\mathsf{R}} - \bbeta^*) - \bar{\bSigma}^{-1/2} \sum_{e\in \mathcal{E}} \omega^{(e)} \mathbb{E} [\bx^{(e)} \varepsilon^{(e)}] \\
    &= \bu^\top \bar{\bSigma}^{1/2} \left(\bar{\bbeta}^{\mathsf{R}} - \bbeta^* - \bar{\bSigma}^{-1} \sum_{e\in \mathcal{E}} \omega^{(e)} \mathbb{E} [\bx^{(e)} \varepsilon^{(e)}] \right) = 0.
\end{align*} Note that $\bu^\top \bar{\bSigma}^{-1/2} \bx^{(e)}$ is a zero-mean sub-Gaussian random variable with parameter $\sigma_x$, and $(\bx^{(e)})^\top (\bar{\bbeta}^{\mathsf{R}} - \bbeta^*) - \varepsilon^{(e)}$ is the sum of two zero-mean sub-Gaussian random variables with parameter $\sigma_\varepsilon$ and $\sigma_x \|\bSigma^{1/2}(\bar{\bbeta}^{\mathsf{R}}-\bbeta^*)\|_2$, respectively. It then follows from \cref{lemma:product-sub-gaussian} and \ref{lemma:sum-sub-exp} that
\begin{align*}
    \mathbb{P}\left[ |\bu_k^\top \bb| \ge C_5\sigma_x \left(\sigma_{\varepsilon} + \sigma_x \|\bSigma^{1/2}(\bar{\bbeta}^{\mathsf{R}}-\bbeta^*)\|_2\right) \left(\sqrt{\frac{u}{n_{\bomega}}} + \frac{u}{n_*}\right)\right] \ge 1-2e^{-u}.
\end{align*} Applying union bounds over all the $k$ and setting $u=t+\log(2N)$ completes the proof. \qed

\subsection{Proof of \cref{thm:rate-low-dim}}
\label{sec:proof-thm2}

\subsubsection{Key Proof Idea}

We will use the following decomposition in the proof of \cref{thm:rate-low-dim}.
\begin{align}
\label{eq:q-decomposition}
\begin{split}
\mathsf{Q}(\bbeta;\gamma,\bomega) - \mathsf{Q}(\bbeta^*;\gamma,\bomega)  =& ~\mathsf{R}(\bbeta) - \mathsf{R}(\bbeta^*) + \gamma \left(\mathsf{J}(\bbeta) - \mathsf{J}(\bbeta^*)\right) \\
	=& ~\left\{\mathsf{R}(\bbeta) - \mathsf{R}(\bbeta^*)\right\} - \left\{ \hat{\mathsf{R}}(\bbeta) - \hat{\mathsf{R}}(\bbeta^*)\right\} \\ 
		&~~~+ \left\{\hat{\mathsf{R}}(\bbeta) + \gamma \hat{\mathsf{J}}(\bbeta) - \hat{\mathsf{R}}(\bbeta^*) - \gamma^*\hat{\mathsf{J}}(\bbeta^*)\right\} \\
		&~~~+ \gamma \left[\left\{\mathsf{J}(\bbeta) - \mathsf{J}(\bbeta^*)\right\} - \left\{ \hat{\mathsf{J}}(\bbeta) - \hat{\mathsf{J}}(\bbeta^*)\right\}\right] \\
	=& ~\mathsf{T}_{\mathsf{R}}(\bbeta, \bbeta^*) + \mathsf{T}_{\mathsf{J}}(\bbeta, \bbeta^*) + \hat{\mathsf{Q}}(\bbeta;\gamma, \bomega) -  \hat{\mathsf{Q}}(\bbeta^*;\gamma, \bomega).
\end{split}
\end{align}

The analysis of the first term $\mathsf{T}_{\mathsf{R}}$ is standard. As stated in the following Lemma, we can provide an instance-dependent two-side bound for it.
\begin{lemma}[Instance-dependent Two-side Bound for ${\mathsf{R}}$]
\label{lemma:instance-dependent-bound-r-lowdim}
Assume \cref{cond0-model}--\ref{cond3-sub-gaussian-eps} hold. Define the event
\begin{align}
\begin{split}
    \mathcal{A}_{1,t} = \Bigg\{ \forall \bbeta \in \mathbb{R}^p, ~~ &~ \left|\mathsf{R}(\bbeta) - \mathsf{R}(\bbeta^*) - \hat{\mathsf{R}}(\bbeta) + \hat{\mathsf{R}}(\bbeta^*)\right| \\
    & ~~~~~~~ \le c_1 \left(\kappa_U \sigma_x^2 \delta_1 \|\bbeta - \bbeta_*\|_2^2 + \kappa_U^{1/2} \sigma_x \sigma_\varepsilon \delta_1 \|\bbeta - \bbeta_*\|_2 \right)\Bigg\}.
\end{split} 
\end{align} with $\delta_1 = \sqrt{\frac{p + t}{n_{\bomega}}} + \frac{p + t}{n_*}$ and some universal constants $c_1$. We have 
\begin{align*}
\mathbb{P}[\mathcal{A}_{1,t}] \ge 1-e^{-t}
\end{align*} for any $t>0$.
\end{lemma}
\begin{proof}[Proof of \cref{lemma:instance-dependent-bound-r-lowdim}]
    It concludes by applying \cref{lemma:instance-dependent-bound-r-highdim} with $s=p$.
\end{proof}

However, the analysis of the second term $\mathsf{T}_{\mathsf{J}}$ is more involved. The following Lemma provides an instance-dependent one-side bound because all the statistical analysis only cares about the upper bound of $\mathsf{T}_{\mathsf{J}}(\bbeta,\bbeta^*)$, or the lower bound of $-\mathsf{T}_{\mathsf{J}}(\bbeta,\bbeta^*)$.

\begin{lemma}[Instance-dependent One-side Bound for ${\mathsf{J}}$]
\label{lemma:instance-dependent-bound-j-lowdim}
Assume \cref{cond0-model}--\ref{cond3-sub-gaussian-eps} hold. Define the event
\begin{align}
\begin{split}
    \mathcal{A}_{2,t} = \Bigg\{ \forall \bbeta \in \mathbb{R}^p, ~~& \frac{1}{c_1} \left(\mathsf{J}(\bbeta) - \mathsf{J}(\bbeta^*) - \hat{\mathsf{J}}(\bbeta) + \hat{\mathsf{J}}(\bbeta^*)\right) \\
    &\le \|\bbeta - \bbeta_*\|_2^2 \times \kappa_U^2\sigma_x^2\sqrt{\frac{p+t}{n_*}} \\
    &~~~~~~ + \|\bbeta - \bbeta_*\|_2 \times \kappa_U^{3/2}\sigma_x^2 \sigma_\varepsilon \left( \sqrt{\frac{p + t}{n_*}} +\sigma_x \frac{p + \log(2|\mathcal{E}|)+t}{\bar{n}}\right)\\
    &~~~~~~ + \sqrt{\sum_{e\in \mathcal{E}} \omega^{(e)} \left\|\mathbb{E} [\bx_{\supp(\bbeta)}^{(e)} \varepsilon^{(e)}]\right\|_2^2} \times \kappa_U^{1/2} \sigma_x \sigma_\varepsilon \sqrt{\frac{p+t}{n_*}}  \\
    &~~~~~~ + |S^*\setminus \supp(\bbeta)| \times \kappa_U \sigma_x^2\sigma_\varepsilon^2  \frac{t+ \log(2|\mathcal{E}||S^*|)}{\bar{n}} + \kappa_U \sigma_x \sigma_\varepsilon^2 \frac{p+t}{n_*}
    \Bigg\}
\end{split}
\end{align} for some universal constant $c_1$. Then we have
\begin{align}
    \mathbb{P} [\mathcal{A}_{2,t}] \ge 1-6e^{-t}.
\end{align} for any $t\in (0, n_{\min} - \log(2|\mathcal{E}|) - p]$.

\end{lemma}
\begin{proof}[Proof of \cref{lemma:instance-dependent-bound-j-lowdim}]
    It concludes by applying \cref{lemma:instance-dependent-bound-j-highdim} with $s=p$ and observing that $n_{\bomega} \ge n_*$.
\end{proof}

We will use the following characterization of the population-level excess risk.

\begin{proposition}
\label{prop:population-lb}
    Under \cref{cond0-model}, \ref{cond1-well-condition} and \ref{cond-ident}. If $\gamma \ge 3\gamma^*$, then
    \begin{align*}
        \mathsf{Q}(\bbeta;\gamma,\bomega) - \mathsf{Q}(\bbeta^*;\gamma,\bomega) &\ge \frac{\kappa_L}{2} \|\bbeta - \bbeta^*\|_2^2 + \frac{\gamma}{3} \mathsf{J}(\bbeta;\bomega) \\
        &\ge \frac{\kappa_L}{2} \|\bbeta - \bbeta^*\|_2^2 + \frac{\gamma}{6} \kappa_L^2 \bar{\mathsf{d}}_{\supp(\bbeta)} + \frac{\gamma}{6} \mathsf{J}(\bbeta;\bomega).
    \end{align*}
\end{proposition}
\begin{proof}[Proof of \cref{prop:population-lb}] The proof is almost identical to the proof of \cref{thm:global-strong-convexity}. From \eqref{eq:j-to-bs} we know $\mathsf{J}(\bbeta;\bomega) \ge \kappa_L^2\bar{\mathsf{d}}_{\supp(\bbeta)}$. Following the notations in the proof of \cref{thm:global-strong-convexity}, we have
\begin{align*}
    \mathsf{Q}(\bbeta;\gamma,\bomega) - \mathsf{Q}(\bbeta^*;\gamma,\bomega) &\ge \mathsf{T}_1(\bbeta) + \gamma \mathsf{T}_2(\bbeta) \\
    &\ge \frac{\kappa_L}{2} \|\bbeta-\bbeta^*\|_2^2 - 2 \kappa_L^{-1} \mathsf{b}_{\supp(\bbeta)} + \frac{2}{3} \gamma \bar{\mathsf{d}}_{\supp(\bbeta)} + \frac{\gamma}{3} \mathsf{J}(\bbeta;\bomega)\\
    &\ge \frac{\kappa_L}{2} \|\bbeta-\bbeta^*\|_2^2 + \frac{\gamma}{3} \mathsf{J}(\bbeta;\bomega)
\end{align*} This completes the proof.
\end{proof}

With the help of the above claims, we are ready to prove \cref{thm:rate-low-dim}.

\subsubsection{Proof of the Variable Selection Property \eqref{eq:low-dim-variable-selection}}

It follows from our decomposition \eqref{eq:q-decomposition} that, for any $\bbeta \in \mathbb{R}^p$,
\begin{align}
\label{eq:low-dim-variable-selection-pf1}
    \hat{\mathsf{Q}}(\bbeta;\gamma, \bomega) - \hat{\mathsf{Q}}(\bbeta^*;\gamma, \bomega)
    \ge - \mathsf{T}_{\mathsf{J}}(\bbeta, \bbeta^*) - \mathsf{T}_{\mathsf{R}}(\bbeta, \bbeta^*) + \mathsf{Q}(\bbeta; \gamma, \bomega) - \mathsf{Q}(\bbeta^*; \gamma, \bomega).
\end{align}

The rest of the proof proceeds conditioned on that $n_{\min} > \log(2|\mathcal{E}|) + p + t$ such that the results of \cref{lemma:instance-dependent-bound-j-lowdim} is applicable. Recall that $\gamma \ge 1 + \kappa_L$. Applying \cref{lemma:instance-dependent-bound-r-lowdim} and \cref{lemma:instance-dependent-bound-j-lowdim}, we have that, under the event $\mathcal{A}_{1,t} \cap \mathcal{A}_{2,t}$, 
\begin{align}
\label{eq:low-dim-tj-tr-bound}
\begin{split}
    \mathsf{T}_{\mathsf{J}}(\bbeta, \bbeta^*) + \mathsf{T}_{\mathsf{R}}(\bbeta, \bbeta^*) &\le C\gamma \Bigg\{\|\bbeta - \bbeta^*\|^2_2 \times \kappa_U^2 \sigma_x^2 \sqrt{\frac{p + t}{n_*}} \\
    &~~~~~~~~~ + \|\bbeta - \bbeta_*\|_2 \times \kappa_U^{3/2}\sigma_x^2 \sigma_\varepsilon \left( \sqrt{\frac{p + t}{n_*}} +\sigma_x \frac{p + \log(2|\mathcal{E}|)+t}{\bar{n}}\right)\\
    &~~~~~~~~~ + \sqrt{\sum_{e\in \mathcal{E}} \omega^{(e)} \left\|\mathbb{E} [\bx_{\supp(\bbeta)}^{(e)} \varepsilon^{(e)}]\right\|_2^2} \times \kappa_U^{1/2} \sigma_x \sigma_\varepsilon \sqrt{\frac{p+t}{n_*}}  \\
    &~~~~~~~~~ + |S^*\setminus \supp(\bbeta)| \times \kappa_U \sigma_x^2\sigma_\varepsilon^2  \frac{t+ \log(2|\mathcal{E}||S^*|)}{\bar{n}} + \kappa_U \sigma_x \sigma_\varepsilon^2 \frac{p+t}{n_*} \Bigg\}\\
    &= C\gamma(\mathsf{I}_1 + \mathsf{I}_2 + \mathsf{I_3} + \mathsf{I_4}).
\end{split}
\end{align}
Applying Young's inequality gives
\begin{align}
\label{eq:low-dim-i2}
\mathsf{I}_2 \le \frac{\kappa_L}{16C\gamma} \|\bbeta - \bbeta^*\|^2_2 + \frac{8C\gamma}{\kappa_L} \kappa_U^3 \sigma_x^4 \sigma_\varepsilon^2 \left\{\frac{p + t}{n_*} + \sigma_x^2 \left(\frac{p + \log(2|\mathcal{E}|) + t}{\bar{n}}\right)^2\right\}.
\end{align} 
At the same time, it follows from \cref{lemma:removecondition7} that
\begin{align*}
    \mathsf{I}_3 &\le \kappa_U^{1/2} \sigma_x \sigma_\varepsilon \sqrt{\frac{p + t}{n_*}} \sqrt{2 \mathsf{J}(\bbeta;\bomega) + 2\kappa_U^2 \|\bbeta - \bbeta^*\|_2^2} \\
    &\le \frac{1}{6C} \mathsf{J}(\bbeta; \bomega) + 6C\kappa_U\sigma_x^2 \sigma_\varepsilon^2 \frac{p + t}{n_*} + \frac{\kappa_L}{16C\gamma} \|\bbeta - \bbeta^*\|_2^2 + \frac{16C\gamma}{\kappa_L}\kappa_U^3 \sigma_x^2 \sigma_\varepsilon^2 \frac{p + t}{n_*}
\end{align*}
Using the fact that
\begin{align*}
|S^*\setminus \supp(\bbeta)| \min_{j\in S^*} |\beta_j^*|^2 &= \sum_{j\in S^*, \beta_j=0} \min_{j\in S^*} |\beta_j^*|^2 \\
&\le \sum_{j\in S^*, \beta_j=0} |\beta_j - \beta^*_j|^2 \le \sum_{j=1}^p |\beta_j - \beta_j^*|^2 = \|\bbeta - \bbeta^*\|_2^2,
\end{align*} we obtain \begin{align}
&|S^*\setminus \supp(\bbeta)| \times \kappa_U \sigma_x^2\sigma_\varepsilon^2  \frac{t+ \log(2|\mathcal{E}||S^*|)}{\bar{n}} \nonumber \\
&~~~~~~~~~ = \left(|S^*\setminus \supp(\bbeta)| \min_{j\in S^*}|\beta_j^*|^2 \right) \times \left\{\frac{\kappa_U \sigma_x^2\sigma_\varepsilon^2}{ \min_{j\in S^*} |\bbeta_j^ *|^2}\frac{t+ \log(2|\mathcal{E}||S^*|)}{\bar{n}}\right\} \nonumber \\
&~~~~~~~~~ \le \|\bbeta - \bbeta_*\|_2^2  \left\{\frac{\kappa_U \sigma_x^2\sigma_\varepsilon^2}{ \min_{j\in S^*} |\bbeta_j^ *|^2}\frac{t+ \log(2|\mathcal{E}||S^*|)}{\bar{n}}\right\}, \label{eq:low-dim-bound-i4-1}
\end{align} 

Putting these pieces together, if 
\begin{align}
\label{eq:low-dim-cond-n-1}
    \bar{n} \ge \frac{16 C\gamma \kappa_U \sigma_x^2\sigma_\varepsilon^2 \{t + \log(2|\mathcal{E}||S^*|)\}}{ \kappa_L \min_{j\in S^*} |\beta_j^*|^2} ~~~~\text{and}~~~~ n_* \ge (p + t) \left(\frac{16C \kappa_U^2 \sigma_x^2}{\kappa_L }\gamma\right)^2,
\end{align} then we find that, under $\mathcal{A}_{1,t} \cap \mathcal{A}_{2,t}$,
\begin{align*}
    \mathsf{T}_{\mathsf{J}}(\bbeta, \bbeta^*) + \mathsf{T}_{\mathsf{R}}(\bbeta, \bbeta^*) &\le \frac{\kappa_L}{4} \|\bbeta - \bbeta^*\|_2^2 + \frac{\gamma}{6} \mathsf{J}(\bbeta;\omega) + C_1 \gamma\Bigg\{\frac{\gamma}{\kappa_L} \kappa_U^3 \sigma_x^4 \sigma_\varepsilon^2  \frac{p+t}{n_*} \\
    &~~~~~~~~~~~~~~~~ + \frac{\gamma}{\kappa_L} \kappa_U^3 \sigma_x^6 \sigma_\varepsilon^2 \left(\frac{p + \log(2|\mathcal{E}|) + t}{\bar{n}}\right)^2 \Bigg\}.
\end{align*} for another universal constant $C_1$.

Plugging the above inequality back into \eqref{eq:low-dim-variable-selection-pf1} and applying \cref{prop:population-lb} provided our choice of $\gamma \ge 3\gamma^*$, we establish that, if $\mathcal{A}_{1,t} \cap \mathcal{A}_{2,t}$ occurs, then for any $\bbeta \in \mathbb{R}^p$,
\begin{align*}
    \hat{\mathsf{Q}}(\bbeta) - \hat{\mathsf{Q}}(\bbeta^*) \ge &~\frac{\kappa_L}{4} \|\bbeta - \bbeta^*\|_2^2 + \frac{\gamma}{6} \kappa_L^2 \bar{\mathsf{d}}_{\supp(\bbeta)} \\
    &~~~~~~~~~~~~ - C_2 \Bigg\{ \frac{\gamma^2}{\kappa_L} \kappa_U^3 \sigma_x^6 \sigma_\varepsilon^2 \left(\frac{p + \log(2|\mathcal{E}|) + t}{\bar{n}}\right)^2 + \frac{\gamma^2}{\kappa_L} \kappa_U^3 \sigma_x^4 \sigma_\varepsilon^2 \frac{p + t}{n_*} \Bigg\}.
\end{align*}

The rest of the proof proceeds conditioned on the event $\mathcal{A}_{1,t} \cap \mathcal{A}_{2,t}$. We first argue that if
\begin{align}
\label{eq:low-dim-cond-n-2}
    \frac{n_*}{p + t} \ge \frac{18C_2 \kappa_U^3\sigma_x^4\sigma_\varepsilon^2}{(\kappa_L^3/\gamma)\mathsf{s}_-} ~~~~\text{and}~~~~ \frac{\bar{n}}{(p + \log(2|\mathcal{E}|) + t)} \ge \frac{\sqrt{18C_2} \kappa_U^{3/2} \sigma_x^3\sigma_\varepsilon}{\sqrt{(\kappa_L^3/\gamma) \mathsf{s}_{-}}},
\end{align} with $\mathsf{s}_{-} = \min_{S\cap G_{\bomega} \neq \emptyset} \bar{\mathsf{d}}_S$, then for any $\bbeta$ with $\supp(\bbeta) \cap G_{\bomega} \neq \emptyset$, the following holds,
\begin{align*}
    \hat{\mathsf{Q}}(\bbeta) - \hat{\mathsf{Q}}(\bbeta^*) \ge \frac{\kappa_L}{4} \|\bbeta - \bbeta^*\|_2^2 + \left(\frac{\gamma}{6} - 2\times \frac{\gamma}{18} \right) \kappa_L^2 \mathsf{s}_- > 0,
\end{align*} which means $\bbeta$ is not the empirical risk minimizer. This implies that $\supp(\hat{\bbeta}_{\mathsf{Q}}) \subseteq G_{\bomega}^c$. 

Meantime, denote $\mathsf{s}_+ = \min_{j\in S^*} |\beta_j^*|^2$. We then argue that if
\begin{align}
\label{eq:low-dim-cond-n-3}
    \frac{n_*}{p + t} \ge \frac{12C_2 \gamma^2 \kappa_U^3\sigma_x^4\sigma_\varepsilon^2}{\kappa_L^2 \mathsf{s}_+} ~~~~\text{and}~~~~ \frac{\bar{n}}{(p + \log(2|\mathcal{E}|) + t)} \ge \frac{\sqrt{12C_2}\gamma \kappa_U^{3/2} \sigma_x^3\sigma_\varepsilon}{  \sqrt{\kappa_L^2 \mathsf{s}_{+}}}
\end{align} then for any $\bbeta$ with $S^* \nsubseteq \supp(\bbeta)$, the following holds
\begin{align*}
    \hat{\mathsf{Q}}(\bbeta) - \hat{\mathsf{Q}}(\bbeta^*) \ge \frac{\kappa_L}{4} \|\bbeta - \bbeta^*\|_2^2 - 2\times \frac{\kappa_L}{12} \min_{j\in S^*} |\beta_j^*|^2 \ge \frac{\kappa_L}{12} \min_{j\in S^*} |\beta_j^*|^2 > 0,
\end{align*} which also means that $\bbeta$ is not the empirical risk minimizer, hence implying the $S^*\subseteq \supp(\hat{\bbeta}_{\mathsf{Q}})$. 

Combining the conditions on $n_*$ and $\bar{n}$ in \eqref{eq:low-dim-cond-n-1}, \eqref{eq:low-dim-cond-n-2} and \eqref{eq:low-dim-cond-n-3} together, we can conclude that if
\begin{align}
\label{eq:low-dim-req-n-bar}
    \frac{\bar{n}}{p+\log(|\mathcal{E}|) + t} \ge \underbrace{ (18C_2 \kappa_U^{3/2} \sigma_x^3 \sigma_\varepsilon) \lor (16C\kappa_U \sigma_x^2\sigma_\varepsilon^2)}_{c_1} \frac{\gamma}{\kappa_L} \left( \sqrt{\frac{1}{\mathsf{s}_+ \land (\gamma\kappa_L \mathsf{s}_{-})}} + \frac{1}{\mathsf{s}_+} \right)
\end{align} and
\begin{align}
\label{eq:low-dim-req-n-star}
    \frac{n_*}{p + t} \ge (\gamma/\kappa_L)^2 \left\{\underbrace{18C_2\kappa_U^3 \sigma_x^4 \sigma_\varepsilon^2}_{c_2} \frac{1}{\mathsf{s} \land (\gamma\kappa_L \mathsf{s}_{-})} + \underbrace{(16C^2)^2\kappa_U^4 \sigma_x^4}_{c_3}\right\}
\end{align}
then under $\mathcal{A}_{1,t}\cap \mathcal{A}_{2,t}$ that occurs with probability at least $1-e^{-7}$ due to \cref{lemma:instance-dependent-bound-j-lowdim}, \cref{lemma:instance-dependent-bound-r-lowdim} and union bound, the following holds
\begin{align}
\label{eq:proof-low-dim-variable-selection}
    S^*\subseteq \supp(\hat{\bbeta}_{\mathsf{Q}}) \subseteq G_{\bomega}^c.
\end{align} This completes the proof. \qed



\subsection{Proof of \cref{thm:rate-faster-low-dim}}
\label{sec:proof-thm3}

\subsubsection{Proof of the Rate \eqref{eq:low-dim-rate2}}
The proof is very similar to that of \cref{thm:rate-low-dim} but will use \cref{lemma:instance-dependent-bound-j-lowdim} and the strong convexity of $\mathsf{Q}$ around $\bbeta^*$ in a different way. The proof proceeds conditioned on $\mathcal{A}_{1,t} \cap \mathcal{A}_{2,t}$. It follows from \cref{lemma:instance-dependent-bound-j-lowdim} and \cref{lemma:instance-dependent-bound-r-lowdim} that the inequality in \eqref{eq:low-dim-tj-tr-bound} holds. We will bound the term $\mathsf{I}_4$ differently.

For first term of $\mathsf{I}_4$, it also follows from the fact $xy \le \frac{1}{2}(x^2 + y^2)$ that
\begin{align*}
&|S^*\setminus \supp(\bbeta)| \times \kappa_U \sigma_x^2\sigma_\varepsilon^2  \frac{t+ \log(2|\mathcal{E}||S^*|)}{\bar{n}} \\
&~~~~~~~~~ = \left(\sqrt{|S^*\setminus \supp(\bbeta)|} \min_{j\in S^*}|\beta_j^*| \sqrt{\kappa_L/16C\gamma}\right) \times \left\{\frac{2\sqrt{|S^*|}\kappa_U \sigma_x^2\sigma_\varepsilon^2}{\sqrt{\kappa_L/16C\gamma} \min_{j\in S^*} |\bbeta_j^ *|}\frac{t+ \log(2|\mathcal{E}||S^*|)}{\bar{n}}\right\} \\
&~~~~~~~~~ \le \frac{\kappa_L}{16C\gamma} |S^*\setminus \supp(\bbeta)| \min_{j\in S^*} |\bbeta_j^*|^2 + \frac{16C \gamma \kappa_U^2 \sigma_x^4\sigma_\varepsilon^2 |S^*|\{t + \log(2|\mathcal{E}||S^*|)\}^2}{\kappa_L (\min_{j\in S^*} |\beta_j^*|^2) \bar{n}^2} \\
&~~~~~~~~~ \le  \frac{\kappa_L}{16C\gamma} \|\bbeta - \bbeta^*\|_2^2 + 16C \kappa_U^2 \sigma_x^4\sigma_\varepsilon^2 \frac{(\gamma/\kappa_L) |S^*|}{\min_{j\in S^*} |\beta_j^*|^2} \left(\frac{t + \log(2|\mathcal{E}||S^*|)}{\bar{n}}\right)^2, 
\end{align*} 

At the same time, we also have $\mathsf{I}_1 \le \frac{\kappa_L}{16} \|\bbeta - \bbeta^*\|_2^2$ if 
\begin{align}
\label{eq:low-dim-faster-rate-c1}
n_* \ge \underbrace{\left(16C \kappa_U \sigma_x^2\right)}_{c_1} (p + t) \left(\gamma/\kappa_L\right)^2.
\end{align}

Putting these pieces together with the error bounds we have for $\mathsf{I}_2$--$\mathsf{I}_3$ in the proof of \cref{thm:rate-low-dim}, the following holds for all the $\bbeta \in \mathbb{R}^p$,
\begin{align*}
\mathsf{T}_{\mathsf{R}}(\bbeta,\bbeta^*) + \mathsf{T}_{\mathsf{J}}(\bbeta, \bbeta^*) &\le 4\times \frac{\kappa_L}{16} \|\bbeta - \bbeta^*\|_2^2 + \frac{\gamma}{6} \mathsf{J}(\bbeta;\omega) \\
&~~~~~~~~~~ + C_1 \frac{\gamma^2}{\kappa_L} \kappa_U^3 \sigma_x^4 \sigma_\varepsilon^2 \left\{\frac{p + t}{n_*} + \sigma_x^2 \left(\frac{p + \log(2|\mathcal{E}|) + t}{\bar{n}}\right)^2 \right\} \\
&~~~~~~~~~~ + C_1 \frac{\gamma^2}{\kappa_L} \kappa_U^2 \sigma_x^4\sigma_\varepsilon^2 \frac{|S^*|}{\min_{j\in S^*} |\beta_j^*|^2} \left(\frac{t + \log(2|\mathcal{E}||S^*|)}{\bar{n}}\right)^2
\end{align*}

Plugging $\hat{\bbeta}_{\mathsf{Q}}$ into the decomposition \eqref{eq:q-decomposition}, it then follows from \cref{prop:population-lb} that
\begin{align*}
    \frac{\kappa_L}{2}\|\hat{\bbeta}_{\mathsf{Q}} - \bbeta^*\|_2^2 + \frac{\gamma}{6} \mathsf{J}(\hat{\bbeta}_{\mathsf{Q}}; \bomega) + 0 &\le \mathsf{Q}(\hat{\bbeta}_{\mathsf{Q}};\gamma, \bomega) - \mathsf{Q}(\hat{\bbeta}_{\mathsf{Q}};\gamma, \bomega) \\
    &= \hat{\mathsf{Q}}(\hat{\bbeta}_{\mathsf{Q}};\gamma, \bomega) - \hat{\mathsf{Q}}(\hat{\bbeta}_{\mathsf{Q}};\gamma, \bomega) + \mathsf{T}_{\mathsf{R}}(\hat{\bbeta}_{\mathsf{Q}},\bbeta^*) + \mathsf{T}_{\mathsf{J}}(\hat{\bbeta}_{\mathsf{Q}}, \bbeta^*) \\
    &\le 0 + \mathsf{T}_{\mathsf{R}}(\hat{\bbeta}_{\mathsf{Q}},\bbeta^*) + \mathsf{T}_{\mathsf{J}}(\hat{\bbeta}_{\mathsf{Q}}, \bbeta^*)
\end{align*} Combining it with the upper bound on $\mathsf{T}_{\mathsf{R}}(\bbeta,\bbeta^*) + \mathsf{T}_{\mathsf{J}}(\bbeta, \bbeta^*)$ we derived, we conclude that
\begin{align}
\label{eq:low-dim-faster-rate-c2-4}
\begin{split}
    \frac{\|\hat{\bbeta}_{\mathsf{Q}} - \bbeta^*\|_2^2}{\sigma_\varepsilon^2} &~\le  \underbrace{4C_2\kappa_U^3 \sigma_x^6}_{c_2^2} \frac{\gamma^2}{\kappa_L^2} \Bigg( \frac{p + t}{n_*} + \frac{(p + \log(2|\mathcal{E}|) + t)^2}{\bar{n}^2}\Bigg)\\
    &~~~~~~~~~~ + \underbrace{4C_2\kappa_U^2 \sigma_x^4}_{c_3^2} \frac{(\gamma^2/\kappa_L^2)|S^*|}{\min_{j\in S^*} |\beta_j^*|^2} \left(\frac{t + \log(2|\mathcal{E}||S^*|)}{\bar{n}}\right)^2
\end{split}
\end{align} under $\mathcal{A}_{1,t}\cap \mathcal{A}_{2,t}$. This completes the proof. \qed

\subsubsection{Proof of the Rate \eqref{eq:low-dim-rate1}}


We assume that conditions \eqref{eq:low-dim-req-n-bar} and \eqref{eq:low-dim-req-n-star} are satisfied. Thus, the variable selection property \eqref{eq:proof-low-dim-variable-selection} also holds. Now that we only focus on the case of no pooled linear spurious variables. The covariate dimension is $p_0 = |(G_{\bomega})^c|$. Now we apply \cref{lemma:instance-dependent-bound-r-lowdim}-\ref{lemma:instance-dependent-bound-j-lowdim} with $\bx_{(G_{\bomega})^c}$. When $\mathcal{A}_{1,t} \cap \mathcal{A}_{2,t} \cap \mathcal{A}_{1,t}^{p_0} \cap \mathcal{A}_{2,t}^{p_0}$ occurs, the inequality 
\eqref{eq:low-dim-tj-tr-bound} with $p=p_0$ also holds for $\bbeta = \hat{\bbeta}_\mathsf{Q}$. Observe that $|S^* \setminus \supp(\bbeta)|=0$ given the variable selection property \eqref{eq:proof-low-dim-variable-selection}. The rest of the proof follows similarly to that for \eqref{eq:low-dim-rate2}.


\subsection{Proof of \cref{thm:high-dim}}

We first define some additional notations. Define $s^*=|S^*|$. Let 
\begin{align*}
    V(s,t) = s\log (ep/s) + s^* + t,
\end{align*} and define $0\times \log(1/0) = 0$.
Similar to the low-dimension counterpart, we also need the following two Lemmas.


\begin{lemma}[Instance-dependent One-side Bound for ${\mathsf{J}}$, High-dimensional]
\label{lemma:instance-dependent-bound-j-highdim}
Assume Conditions \ref{cond0-model} -- \ref{cond3-sub-gaussian-eps} hold. For any $1\le s\le p$, define the following event,
\begin{align}
\label{eq:instance-dependent-bound-j-highdim}
\begin{split}
    \mathcal{A}_{3,t}(s) = \Bigg\{ \forall \bbeta \in \mathcal{B}_s, ~~& \frac{1}{c_1} \left(\mathsf{J}(\bbeta) - \mathsf{J}(\bbeta^*) - \hat{\mathsf{J}}(\bbeta) + \hat{\mathsf{J}}(\bbeta^*)\right) \\
    &\le \|\bbeta - \bbeta^*\|_2^2 \times \kappa_U^2\sigma_x^2 \Bigg(\sqrt{\frac{V(s,t)}{n_{\bomega}}} + \frac{V(s,t)}{n_*}\Bigg) \\
    &~~~~~~ + \|\bbeta - \bbeta^*\|_2 \times \kappa_U^{3/2}\sigma_x^2 \sigma_\varepsilon \Bigg(\sqrt{\frac{V(s,t)}{n_{\bomega}}} + \frac{V(s,t)}{n_*}\Bigg)  \\
    &~~~~~~ + \|\bbeta - \bbeta^*\|_2 \times \kappa_U^{3/2}\sigma_x^3 \sigma_\varepsilon \frac{\sqrt{V(s,t+\log(2|\mathcal{E}|)) \cdot V(0, t+\log(2|\mathcal{E}|))}}{\bar{n}}\\
    &~~~~~~ + \|\bbeta - \bbeta^*\|_2 \times \kappa_U^{3/2}\sigma_x^3 \sigma_\varepsilon \frac{V(s,t+\log(2|\mathcal{E}|))\sqrt{V(0, t+\log(2|\mathcal{E}|))}}{n_\dagger} \\
    &~~~~~~ + \sqrt{\sum_{e\in \mathcal{E}} \omega^{(e)} \left\|\mathbb{E} [\bx_{\supp(\bbeta)}^{(e)} \varepsilon^{(e)}]\right\|_2^2} \times \kappa_U^{1/2} \sigma_x \sigma_\varepsilon \sqrt{\frac{V(|\supp(\bbeta) \setminus S^*|,t)}{n_*}}  \\
    &~~~~~~ + |S^*\setminus \supp(\bbeta)| \times \kappa_U \sigma_x^2\sigma_\varepsilon^2  \frac{t+ \log(2s^*|\mathcal{E}|)}{\bar{n}} \\
    &~~~~~~ + \kappa_U \sigma_x \sigma_\varepsilon^2 \frac{V(|\supp(\bbeta)\setminus S^*|,t)}{n_*}
    \Bigg\}
\end{split}
\end{align} for some universal constant $c_1$. Then we have $\mathbb{P} [\mathcal{A}_{3,t}(s)] \ge 1-6e^{-t}$ for any $t\in (0, n_{\min} - \log(2|\mathcal{E}|) - s^*]$.
\end{lemma}
\begin{proof}[Proof of \cref{lemma:instance-dependent-bound-j-highdim}]
    See \cref{subsec:proof-instance-dependent-bound-j-highdim}.
\end{proof}

\begin{lemma}[Instance-dependent Two-side Bound for ${\mathsf{R}}$, High-dimensional]
\label{lemma:instance-dependent-bound-r-highdim}
Assume Conditions \ref{cond0-model} -- \ref{cond3-sub-gaussian-eps} hold. Define the event
\begin{align}
\begin{split}
    \mathcal{A}_{4,t}(s) = \Bigg\{ \forall \bbeta \in \mathcal{B}_s, ~~ &~ \left|\mathsf{R}(\bbeta) - \mathsf{R}(\bbeta^*) - \hat{\mathsf{R}}(\bbeta) + \hat{\mathsf{R}}(\bbeta^*)\right| \\
    & ~~~~~~~ \le c_1 \left(\kappa_U \sigma_x^2 \delta_1 \|\bbeta - \bbeta^*\|_2^2 + \kappa_U^{1/2} \sigma_x \sigma_\varepsilon \delta_1 \|\bbeta - \bbeta^*\|_2 \right)\Bigg\}.
\end{split} 
\end{align} with $\delta_1 = \sqrt{\frac{V(s,t)}{n_{\bomega}}} + \frac{V(s,t)}{n_*}$ and some universal constants $c_1$. We have $\mathbb{P}[\mathcal{A}_{4,t}(s)] \ge 1-e^{-t}$ for any $t>0$.
\end{lemma}
\begin{proof}[Proof of \cref{lemma:instance-dependent-bound-r-highdim}]
    See \cref{subsec:proof-instance-dependent-bound-r-highdim}.
\end{proof}

We start with a lemma stating that if we choose some large $\lambda$, then the $\hat{S} = \supp(\hat{\bbeta}_{\mathsf{L}})$ will satisfy $|\hat{S}| \le 2s^*$ with high probability.

\begin{proposition}
\label{proposition:small-hat-s}
    Let $0<t\le n_{\min} - \log(2ep|\mathcal{E}|) - s^*$ be arbitrary. Assume Conditions \ref{cond0-model}--\ref{cond-ident} hold, and $\gamma \ge 3\gamma^*\lor 1$. Let $\zeta = \log(ep/s^*) + t$. Suppose further that $\log(|\mathcal{E}|) \le c_1 \log p$ for some universal constant $c_1>0$. There exist some universal constants $c_2$--$c_3$ depending only on $C$ such that if \begin{align*}
        \bar{n} &\ge c_2 \kappa_U\sigma_x^2 \gamma (t + \log p) \left\{ s^* +  \sigma_\varepsilon^2 / (\kappa_L \min_{j\in S^*} |\beta_j^*|^2)\right\}, \\
        n_* &\ge c_2 \kappa_U^4\sigma_x^4 (\gamma/\kappa_L)^2  \zeta s^*, \\
        n_\dagger &\ge c_2 \kappa_U^{3/2} \sigma_x^2 (\gamma/\kappa_L) \zeta s^* \sqrt{\zeta + s^*},
    \end{align*}
    and the choice of $\lambda$ satisfies
    \begin{align*}
        \lambda \ge c_3 \sigma_x^2 \sigma_\varepsilon^2 (\gamma/\kappa_L) \zeta \left\{\frac{\kappa_U^4\sigma_x^2 (\gamma/\kappa_L)}{n_*} + \frac{\kappa_U^{3/2} \sqrt{\zeta + s^*}}{n_\dagger} + \frac{\kappa_{U}^3\sigma_x^2 (\gamma/\kappa_L)(\zeta + s^*)}{\bar{n}^2}\right\}.
    \end{align*} Then the following holds
    \begin{align*}
        \mathbb{P}\left[|\hat{S}| \le 2s^*\right] \ge 1-3e^{-t}.
    \end{align*}
\end{proposition}
\begin{proof}[Proof of \cref{proposition:small-hat-s}] 
    See \cref{subsec:proof-small-hat-s}.
\end{proof}

We are ready to prove \cref{thm:high-dim}.

\begin{proof}[Proof of Theorem~\ref{thm:high-dim}]
We prove the Theorem in a more general form, including the tail probability $t$. Denote $\zeta = \log(ep/s^*) + t$. We first provide a detailed condition for different sample sizes presented in \cref{cond6-n-req},
\begin{align}
\label{eq:req-n-high-dim}
\begin{split}
    n_{\min} &\ge C^* (s^* + \log p) \\
    n_* &\ge C^* \Big\{\kappa_U^4\sigma_x^4 (\gamma/\kappa_L)^2 \zeta s^* \Big\} \lor \Big\{\kappa_U^3 \sigma_x^4 \sigma_\varepsilon^2 \zeta s^* (\gamma/\kappa_L)^2 (1/\kappa_L\beta_{\min}^2)\Big\} \\
    \bar{n} &\ge C^* \left\{\kappa_U\sigma_x^2 \gamma (t + \log p) \left\{ s^* +  \sigma_\varepsilon^2 / (\kappa_L \beta_{\min}^2)\right\} \right\}\lor \left\{\kappa_U^{3/2}\sigma_x^2\sigma_\varepsilon (\gamma/\kappa_L) \frac{\sqrt{s^* \zeta (s^* + \zeta)}}{\kappa_L \beta_{\min}}\right\} \\
    n_\dagger &\ge C^* \kappa_U^{3/2} \sigma_x^2 (\gamma/\kappa_L) s^* \zeta \sqrt{s^* + \zeta} \{1 + \sigma_\varepsilon/(\sqrt{\kappa_L} \beta_{\min})\}
\end{split}
\end{align} together with a detailed lower bound on the regularization parameter $\lambda$
\begin{align}
\label{eq:req-lambda-high-dim}
\begin{split}
    \lambda \ge C^*\kappa_U^3 \sigma_x^4 \sigma_\varepsilon^2 (\gamma/\kappa_L)^2 \Bigg\{ &\left(\frac{\kappa_U \zeta}{n_*} \lor \frac{s^* \zeta}{n_*}\right) + \frac{s^* \zeta(s^* + \zeta)}{\bar{n}^2} \\ 
    &~~~~ + \kappa_U^{-3/2} \sigma_x^{-2} (\gamma/\kappa_L)^{-1} \frac{\zeta\sqrt{s^* + \zeta}}{n_\dagger} + \frac{(s^*\zeta)^2(s^* + \zeta)}{n_\dagger^2} \Bigg\}.
\end{split}
\end{align}
Define the two events
\begin{align*}
    \mathcal{C}_{1,t} = \left\{|\hat{S}| \le 2s^*\right\} ~~~~\text{and}~~~~ \mathcal{C}_{2,t} = \mathcal{A}_{3,t}(2s^*) \cap \mathcal{A}_{4,t}(2s^*).
\end{align*} We can see that the conditions in \cref{lemma:instance-dependent-bound-j-highdim}, \cref{lemma:instance-dependent-bound-r-highdim},  and \cref{proposition:small-hat-s} are satisfied by \eqref{eq:req-n-high-dim} and \eqref{eq:req-lambda-high-dim} with large universal constant $C^*>0$. Denote $\phi_* = s^* + \log(ep/s^*) + t = s^* + \zeta$ and $\phi_s = s^* \log(ep/s^*) + t \le s^* \zeta$. We follow a similar strategy as that in \emph{Case 1} of \cref{proposition:small-hat-s}. Under $\mathcal{C}_{1,t} \cap \mathcal{C}_{2,t}$, it follows our choice of $\gamma$ and \cref{prop:population-lb} that, there exists some universal constant $C_1>0$,
\begin{align*}
    \hat{\mathsf{Q}}(\hat{\bbeta}_{\mathsf{L}}) - \hat{\mathsf{Q}}(\bbeta^*) &= \hat{\mathsf{Q}}(\hat{\bbeta}_{\mathsf{L}}) - \mathsf{Q}(\hat{\bbeta}_{\mathsf{L}}) + \mathsf{Q}(\hat{\bbeta}_{\mathsf{L}}) - \mathsf{Q}(\bbeta^*) + \mathsf{Q}(\bbeta^*) - \hat{\mathsf{Q}}(\bbeta^*) \\
    &\ge \frac{\kappa_L}{2} \|\hat{\bbeta}_{\mathsf{L}} - \bbeta^*\|_2^2 + \frac{\gamma}{6} \mathsf{J}(\hat{\bbeta}_{\mathsf{L}};\bomega) + \hat{\mathsf{R}}(\hat{\bbeta}_{\mathsf{L}}) - \mathsf{R}(\hat{\bbeta}_{\mathsf{L}}) + \mathsf{R} (\bbeta^*) - \hat{\mathsf{R}} (\bbeta^*) \\
    & ~~~~~~~~~~~~ - \gamma \left(\mathsf{J}(\hat{\bbeta}_{\mathsf{L}}) - \hat{\mathsf{J}}(\hat{\bbeta}_{\mathsf{L}}) - \mathsf{J}(\bbeta^*) + \hat{\mathsf{J}}(\bbeta^*)\right)\\
    &\overset{(a)}{\ge} \left(\frac{\kappa_L}{2} - C_1 \kappa_U^2\sigma_x^2 \gamma \sqrt{\frac{\phi_s}{n_*}}\right) \|\hat{\bbeta}_{\mathsf{L}} - \bbeta^*\|_2^2 \\
    & ~~~~~~~~~~~~ + \frac{\gamma}{6} \mathsf{J}(\hat{\bbeta}_{\mathsf{L}};\bomega) - C_1 \gamma \sqrt{\sum_{e\in \mathcal{E}} \omega^{(e)} \left\|\mathbb{E}[\bx_{\hat{S}}^{(e)} \varepsilon^{(e)}]\right\|_2^2} \kappa_U^{1/2} \sigma_x\sigma_\varepsilon \sqrt{\frac{|\hat{S}\setminus S^*|\log (ep/s^*)+s^* + t}{n_*}} \\
    & ~~~~~~~~~~~~ - C_1 \|\hat{\bbeta}_{\mathsf{L}} - \bbeta^*\|_2 \times \gamma \kappa_U^{3/2} \sigma_x^2 \sigma_\varepsilon \sqrt{\frac{\phi_s}{n_*}} \\
    & ~~~~~~~~~~~~ - C_1 \|\hat{\bbeta}_{\mathsf{L}} - \bbeta^*\|_2 \times \gamma \kappa_U^{3/2} \sigma_x^3 \sigma_\varepsilon \left(\frac{\sqrt{\phi_*} \sqrt{\phi_s}}{\bar{n}} + \frac{\phi_s \sqrt{\phi_*}}{n_\dagger} \right) \\
    & ~~~~~~~~~~~~ - C_1 \gamma |S^*\setminus \hat{S}| \kappa_U \sigma_x^2 \sigma_\varepsilon^2 \frac{t+ \log p}{\bar{n}} - C_1 \gamma \kappa_U \sigma_x \sigma_\varepsilon^2 \frac{\phi_s}{n_*} \\
    &= \mathsf{I}_1 + \mathsf{I}_2 + \mathsf{I}_3 + \mathsf{I}_4 + \mathsf{I}_5,
\end{align*} where $(a)$ follows from the fact that $\mathcal{C}_{1,t} \cap \mathcal{C}_{2,t}$ holds such that we can apply the result of \cref{lemma:instance-dependent-bound-j-highdim} and \cref{lemma:instance-dependent-bound-r-highdim} with $s=2s^*$ to $\hat{\bbeta}_{\mathsf{L}}$. Denote $\Diamond=0.5 (6C_1)^2 (\gamma/\kappa_L)^2 \kappa_U^3 \sigma_x^4\sigma_\varepsilon^2$. Following a similar strategy of lower bounds on $\mathsf{I}_1$ -- $\mathsf{I}_5$ and using the fact that $\phi_s \le 2s^*\zeta$ and
\begin{align*}
    \mathsf{I}_2 \ge - \frac{\kappa_L}{12} \|\bbeta - \bbeta^*\|_2^2 - 4\Diamond \frac{|\hat{S} \setminus S^*| \zeta}{n_*},
\end{align*} we obtain
\begin{align*}
    &~\hat{\mathsf{Q}}(\hat{\bbeta}_{\mathsf{L}}) - \hat{\mathsf{Q}}(\bbeta^*) + \lambda \|\hat{\bbeta}_{\mathsf{L}}\|_0 - \lambda \|\bbeta^*\|_0 \\
    &~~~~ = \hat{\mathsf{Q}}(\hat{\bbeta}_{\mathsf{L}}) - \hat{\mathsf{Q}}(\bbeta^*) + \lambda |\hat{S} \setminus S^*| - \lambda |S^* \setminus \hat{S}| \\
    &~~~~\ge \frac{\kappa_L}{12} \|[\hat{\bbeta}_{\mathsf{L}}]_{\hat{S}} - [\bbeta^*]_{\hat{S}}\|_2^2 \\
    &~~~~~~~~~~~~ + \frac{\kappa_L}{12} \|[\hat{\bbeta}_{\mathsf{L}}]_{S^*\setminus \hat{S}} - [\bbeta^*]_{S^*\setminus \hat{S}}\|_2^2 - \lambda |S^* \setminus \hat{S}| \\
    &~~~~~~~~~~~~ + \frac{\lambda}{2} |\hat{S} \setminus S^*| - 4\Diamond \cdot \frac{\zeta}{n_*} |\hat{S}\setminus S^*| \\ 
    &~~~~~~~~~~~~ + \frac{\lambda}{2} |\hat{S}\setminus S^*| - 2 \Diamond \cdot (2s^*) \zeta \left(\frac{1}{n_*} + \frac{\phi_*}{\bar{n}^2} + \frac{2s^* \phi_* \zeta^2}{n_\dagger^2}\right) \\
    &~~~~ \overset{(a)}{\ge} |S^* \setminus \hat{S}| \left(\frac{\kappa_L}{12} \min_{j\in S^*}|\beta^*_j|^2 - \lambda \right) \\
    &~~~~~~~~~~~~ + \frac{\lambda}{2} |\hat{S}\setminus S^*| - 4 \Diamond \cdot s^* \zeta \left(\frac{1}{n_*} + \frac{\phi_*}{\bar{n}^2} + \frac{2s^* \phi_* \zeta}{n_\dagger^2}\right) + \left(\frac{\lambda}{2} - 4\Diamond \cdot \frac{\zeta}{n_*}\right) |\hat{S}\setminus S^*|,
\end{align*} 
Here $(a)$ follows from the facts $\|[\hat{\bbeta}_{\mathsf{L}}]_{\hat{S}} - [\bbeta^*]_{\hat{S}}\|_2^2 \ge 0$, $\|[\hat{\bbeta}_{\mathsf{L}}]_{S^*\setminus \hat{S}} - [\bbeta^*]_{S^*\setminus \hat{S}}\|_2^2 \ge |S^* \setminus \hat{S}| \min_{j\in S^*}|\beta^*_j|$. We argue that if
\begin{align}
\label{eq:proof-thm-highdim-n-req}
    \frac{\kappa_L}{36} \min_{j\in S^*} |\beta_j^*|^2 \ge  4 \Diamond \cdot s^* \zeta \left(\frac{1}{n_*} + \frac{\phi_*}{\bar{n}^2} + \frac{2s^* \phi_* \zeta}{n_\dagger^2}\right),
\end{align} which is satisfied by \cref{cond6-n-req}, and $\lambda \le \frac{\kappa_L}{36} \min_{j\in S^*} |\beta_j^*|^2$, then $|S^*\setminus \hat{S}|=0$. This is because if $|S^*\setminus \hat{S}| \ge 1$, we will have
\begin{align*}
    \hat{\mathsf{Q}}(\hat{\bbeta}_{\mathsf{L}}) - \hat{\mathsf{Q}}(\bbeta^*) + \lambda \|\hat{\bbeta}_{\mathsf{L}}\|_0 - \lambda \|\bbeta^*\|_0 \ge |S^*\setminus \hat{S}| \frac{\kappa_L}{36} \min_{j\in S^*}|\beta^*_j|^2 > 0.
\end{align*} This is contrary to the fact that $\hat{\bbeta}_{\mathsf{L}}$ is the minimizer of \eqref{eq:eills-objective-l}. Moreover, we also have $|\hat{S} \setminus S^*| = 0$ if
\begin{align}
\label{eq:proof-thm-highdim-lambda-req2}
    \lambda \ge 9 \Diamond \cdot s^* \zeta \left(\frac{1}{n_*} + \frac{\phi_*}{\bar{n}^2} + \frac{2s^* \phi_* \zeta}{n_\dagger^2}\right).
\end{align} This is because if $|\hat{S} \setminus S^*| \ge 1$, we will have
\begin{align*}
    \hat{\mathsf{Q}}(\hat{\bbeta}_{\mathsf{L}}) - \hat{\mathsf{Q}}(\bbeta^*) + \lambda \|\hat{\bbeta}_{\mathsf{L}}\|_0 - \lambda \|\bbeta^*\|_0 \ge \frac{\lambda}{2} |\hat{S} \setminus S^*| - 4 \Diamond \cdot s^* \zeta \left(\frac{1}{n_*} + \frac{\phi_*}{\bar{n}^2} + \frac{2s^* \phi_* \zeta}{n_\dagger^2}\right) > 0,
\end{align*} which also contradicts the fact that $\hat{\bbeta}_{\mathsf{L}}$ minimizes \eqref{eq:eills-objective-l}. In conclusion, letting $\lambda \le \frac{\kappa_L}{36} \min_{j\in S^*} |\beta_j^*|^2$ and choosing some large enough constant $C^*$ such that \eqref{eq:proof-thm-highdim-n-req} and \eqref{eq:proof-thm-highdim-lambda-req2} are satisfied by \eqref{eq:req-n-high-dim} and \eqref{eq:req-lambda-high-dim}, we can then argue that under the event $\mathcal{C}_{1,t} \cap \mathcal{C}_{2,t}$, which occurs with probability $1-11e^{-t}$, one has $|\hat{S}\setminus S^*| + |S^*\setminus \hat{S}| = 0$, implying that $\hat{S} = S^*$. This completes the proof via setting $t=10\log (ep)$. 
\end{proof}

\subsection{Proof of \cref{lemma:instance-dependent-bound-j-highdim}}
\label{subsec:proof-instance-dependent-bound-j-highdim}

Let $S=\supp(\bbeta)$, $I = S\cap S^*$. It follows from the fact that $\mathsf{J}(\bbeta^*)=0$ and the definition of $\mathsf{J}$ that
\begin{align*}
	\mathsf{J}(\bbeta) - \mathsf{J}(\bbeta^*) - \hat{\mathsf{J}}(\bbeta) + \hat{\mathsf{J}}(\bbeta^*) = \sum_{e\in \mathcal{E}} \frac{\omega^{(e)}}{4} \left\{\|\nabla_S \mathsf{R}^{(e)}(\bbeta)\|^2_2 - \|\nabla_S \hat{\mathsf{R}}^{(e)}(\bbeta)\|_2^2 + \|\nabla_{S^*} \hat{\mathsf{R}}^{(e)}(\bbeta^*)\|_2^2\right\}.
\end{align*}
Observe that, for any fixed $e\in \mathcal{E}$, 
\begin{align}
&\frac{1}{8}\left(\|\nabla_S \mathsf{R}^{(e)}(\bbeta)\|^2_2 - \|\nabla_S \hat{\mathsf{R}}^{(e)}(\bbeta)\|_2^2 + \|\nabla_{S^*} \hat{\mathsf{R}}^{(e)}(\bbeta^*)\|_2^2\right) \nonumber \\
=&~ \frac{1}{8}\left\{\|\nabla_S \mathsf{R}^{(e)}(\bbeta)\|^2_2 - \|\nabla_S \hat{\mathsf{R}}^{(e)}(\bbeta) - \nabla_S \mathsf{R}^{(e)}(\bbeta) + \nabla_S \mathsf{R}^{(e)}(\bbeta)\|_2^2 + \|\nabla_{S^*} \hat{\mathsf{R}}^{(e)}(\bbeta^*)\|_2^2\right\} \nonumber\\
=&~ -\frac{1}{4} \{\nabla_S \mathsf{R}^{(e)}(\bbeta)\}^\top \left\{\nabla_S \hat{\mathsf{R}}^{(e)}(\bbeta) - \nabla_S {\mathsf{R}}^{(e)}(\bbeta)\right\} + \frac{1}{8} \|\nabla_{S^*\setminus S} \hat{\mathsf{R}}^{(e)}(\bbeta^*)\|_2^2 \nonumber\\
&~~~~~~~~~~~ - \frac{1}{8} \|\nabla_{S\setminus S^*} \hat{\mathsf{R}}^{(e)}(\bbeta) - \nabla_{S\setminus S^*} {\mathsf{R}}^{(e)}(\bbeta)\|_2^2 \nonumber\\
&~~~~~~~~~~~ + \frac{1}{8}\left\{\|\nabla_{I} \hat{\mathsf{R}}^{(e)}(\bbeta^*)\|_2^2 - \|\nabla_{I} \hat{\mathsf{R}}^{(e)}(\bbeta) - \nabla_{I} \mathsf{R}^{(e)}(\bbeta)\|_2^2\right\} \nonumber\\
\overset{(a)}{\le}&  \left(\bSigma^{(e)}_{S,:}(\bbeta - \bbeta^*) - \mathbb{E}[\bx_S^{(e)} \varepsilon^{(e)}]\right)^\top \Bigg\{\hat{\mathbb{E}}[\bx_S^{(e)} \varepsilon^{(e)}] - \mathbb{E}[\bx_S^{(e)} \varepsilon^{(e)}]  \nonumber\\
&~~~~~~~~~~~~~~~~ - \left\{\hat{\mathbb{E}}[\bx_S^{(e)} (\bx^{(e)})^\top (\bbeta - \bbeta^*)] - \bSigma^{(e)}_{S,:} (\bbeta- \bbeta^*)\right\}\Bigg\} \nonumber\\
&~~~~~~~~~~~~~~~~ + \left(\hat{\mathbb{E}}[\bx_{I}^{(e)} \varepsilon^{(e)}] \right)^\top \left(\hat{\mathbb{E}} [\bx_{I}^{(e)} (\bx^{(e)})^\top (\bbeta - \bbeta^*)] - \bSigma^{(e)}_{I,:} (\bbeta - \bbeta^*)\right) + \frac{1}{2} \left\|\hat{\mathbb{E}}[\bx_{S^*\setminus S}^{(e)} \varepsilon^{(e)}]\right\|_2^2 \nonumber\\
=&~  (\bbeta - \bbeta^*)^\top \bSigma^{(e)}_{:,S} \left\{\hat{\mathbb{E}}[\bx_S^{(e)} \varepsilon^{(e)}] - \mathbb{E}[\bx_S^{(e)} \varepsilon^{(e)}]\right\} \nonumber\\
&~~~~~~~~~ -\mathbb{E}[\bx^{(e)}_S \varepsilon^{(e)}]^\top \left\{\hat{\mathbb{E}}[\bx_S^{(e)} \varepsilon^{(e)}] - \mathbb{E}[\bx_S^{(e)} \varepsilon^{(e)}]\right\} \nonumber\\
&~~~~~~~~~ +(\bbeta - \bbeta^*)^\top \bSigma_{:,S}^{(e)} \left\{\hat{\mathbb{E}}[\bx_S^{(e)} (\bx^{(e)})^\top (\bbeta - \bbeta^*)] - \bSigma^{(e)}_{S,:} (\bbeta- \bbeta^*)\right\} \nonumber\\
&~~~~~~~~~ -\mathbb{E}[\bx_S^{(e)}\varepsilon^{(e)}]^\top \left\{\hat{\mathbb{E}}[\bx_S^{(e)} (\bx^{(e)})^\top (\bbeta - \bbeta^*)] - \bSigma^{(e)}_{S,:} (\bbeta- \bbeta^*)\right\} \nonumber\\
&~~~~~~~~~ +\left(\hat{\mathbb{E}}[\bx_{I}^{(e)} \varepsilon^{(e)}] \right)^\top \left(\hat{\mathbb{E}} [\bx_{I}^{(e)} (\bx^{(e)})^\top (\bbeta - \bbeta^*)] - \bSigma^{(e)}_{I,:} (\bbeta - \bbeta^*)\right) + \frac{1}{2} \left\|\hat{\mathbb{E}}[\bx_{S^*\setminus S}^{(e)} \varepsilon^{(e)}]\right\|_2^2\nonumber\\
=&~ \mathsf{T}^{(e)}_1(\bbeta) + \mathsf{T}^{(e)}_2(\bbeta) + \mathsf{T}^{(e)}_3(\bbeta) + \mathsf{T}^{(e)}_4(\bbeta) + \mathsf{T}^{(e)}_5(\bbeta) + \mathsf{T}_6^{(e)}(\bbeta), \label{eq:low-dim-id-decomposition}
\end{align}
where $(a)$ follows from the definition of $\nabla \mathsf{R}(\bbeta)$, $\nabla \hat{\mathsf{R}}(\bbeta)$ together with the facts
\begin{align*}
	&\frac{1}{8}\left\{\|\nabla_{I} \hat{\mathsf{R}}^{(e)}(\bbeta^*)\|_2^2 - \|\nabla_{I} \hat{\mathsf{R}}^{(e)}(\bbeta) - \nabla_{I} {\mathsf{R}}^{(e)}(\bbeta)\|_2^2\right\} \\
	=& ~\frac{1}{2} \sum_{j\in I} \left(\hat{\mathbb{E}}[x_j^{(e)} \varepsilon^{(e)}] \right)^2 - \left(- \hat{\mathbb{E}}[x_j^{(e)} \varepsilon^{(e)}] + \hat{\mathbb{E}}[x_j^{(e)} (\bx^{(e)})^\top (\bbeta - \bbeta^*)] - {\mathbb{E}}[x_j^{(e)} (\bx^{(e)})^\top (\bbeta - \bbeta^*)]\right)^2\\
	=& ~\frac{1}{2} \sum_{j \in I} \left(\hat{\mathbb{E}}[x_j^{(e)} (\bx^{(e)})^\top (\bbeta - \bbeta^*)] - {\mathbb{E}}[x_j^{(e)} (\bx^{(e)})^\top (\bbeta - \bbeta^*)]\right) \times \\
	&~~~~~~~~~~~~~~~~~~~~~ \left\{2\hat{\mathbb{E}}[x_j^{(e)} \varepsilon^{(e)}] -\left(\hat{\mathbb{E}}[x_j^{(e)} (\bx^{(e)})^\top (\bbeta - \bbeta^*)] - {\mathbb{E}}[x_j^{(e)} (\bx^{(e)})^\top (\bbeta - \bbeta^*)]\right)\right\} \\
	\le& ~ \sum_{j\in I} \hat{\mathbb{E}}[x_j^{(e)} \varepsilon^{(e)}] \times \left(\hat{\mathbb{E}}[x_j^{(e)} (\bx^{(e)})^\top (\bbeta - \bbeta^*)] - {\mathbb{E}}[x_j^{(e)} (\bx^{(e)})^\top (\bbeta - \bbeta^*)]\right) \\
	=& ~\left(\hat{\mathbb{E}}[\bx_{I}^{(e)} \varepsilon^{(e)}] \right)^\top \left(\hat{\mathbb{E}} [\bx_{I}^{(e)} (\bx^{(e)})^\top (\bbeta - \bbeta^*)] - \bSigma^{(e)}_{I,:} (\bbeta - \bbeta^*)\right),
\end{align*} and $-\frac{1}{8} \|\nabla_{S\setminus S^*} \hat{\mathsf{R}}^{(e)}(\bbeta) - \nabla_{S\setminus S^*} {\mathsf{R}}^{(e)}(\bbeta)\|_2^2 \le 0$.



The rest of the proof will be divided into several pieces deriving the instance-dependent high-probability upper bounds on $\sum_{e\in \mathcal{E}} \omega^{(e)} \mathsf{T}_k^{(e)}(\bbeta)$ for each $k\in [6]$.

%%%%%%%%%%%%%%%%%%%%%%%%%%%%%%%%%%%% Step 1
\noindent {\sc Step 1. Upper bound on $\mathsf{T}_1^{(e)}$.} Define the event
\begin{align}
\label{eq:low-dim-id-bound1}
	\mathcal{C}_{1,t} = \left\{\forall \bbeta \in \mathcal{B}_s, ~~\sum_{e\in \mathcal{E}} \omega^{(e)} \mathsf{T}_1^{(e)} \le C_1 \kappa_U^{3/2} \sigma_x \sigma_\varepsilon \|\bbeta - \bbeta^*\|_2 \Bigg(\sqrt{\frac{V(s,t)}{n_{\bomega}}} + \frac{V(s,t)}{n_*}\Bigg)\right\}
\end{align} for some constant $C_1$ to be determined. We will claim that $\mathbb{P}(\mathcal{C}_{1,t}) \ge 1-e^{-t}$ in this step. We further assume $\bbeta \neq \bbeta^*$ since the inequality holds trivially when $\bbeta = \bbeta^*$. It then follows from Cauchy-Schwarz inequality that
\begin{align*}
	\sup_{\bbeta \in \mathcal{B}_s, \bbeta \neq \bbeta^*} \frac{\sum_{e\in \mathcal{E}} \omega^{(e)} \mathsf{T}_1^{(e)}}{\|\bbeta - \bbeta^*\|_2} &= \sup_{\bbeta \in \mathcal{B}_s, \bbeta \neq \bbeta^*} \frac{(\bbeta - \bbeta^*)^\top}{\|\bbeta - \bbeta^*\|_2} \sum_{e\in \mathcal{E}} \omega^{(e)} \bSigma^{(e)}_{S\cup S^*,S} \left\{\hat{\mathbb{E}}[\bx_S^{(e)} \varepsilon^{(e)}] - \mathbb{E}[\bx^{(e)} \varepsilon^{(e)}]\right\} \\
	&\le \sup_{|S| \le s} \left\|\sum_{e\in \mathcal{E}} \omega^{(e)} \bSigma^{(e)}_{S\cup S^*,S} \left\{\hat{\mathbb{E}}[\bx_S^{(e)} \varepsilon^{(e)}] - \mathbb{E}[\bx^{(e)} \varepsilon^{(e)}]\right\}\right\|_2.
\end{align*}

For any $S\subseteq [p]$ with $|S|\le s$, let $\bv^{(S)}_1, \ldots, \bv^{(S)}_{N_S}$ be an $1/4-$covering of $\mathcal{B}(S \cup S^*)$, that is, for any $\bv\in \mathcal{B}(S\cup S^*)$, there exists some $\pi(\bv) \in [N_S]$ such that
\begin{align}
\label{eq:proof-cover1}
	\|\bv - \bv_{\pi(v)}^{(S)}\|_2 \le 1/4.
\end{align}
	It follows from standard empirical process result that $N_S \le 9^{|S\cup S^*|}$, then 
\begin{align}
\label{eq:calculate-cover-number}
\begin{split}
	N = \sum_{|S| \le s} N_S \le \sum_{|S|\le s} 9^{|S\cup S^*|} &\le \sum_{i=0}^s 9^{i+s^*} \binom{p}{i} \\
            &\le 9^{s^*} \times  \left(\frac{9p}{s}\right)^{s} \sum_{i=0}^s \left(\frac{s}{p}\right)^i \binom{p}{i} \le 9^{s^*} \times  \left(\frac{9p}{s}\right)^{s} \sum_{i=0}^p \left(\frac{s}{p}\right)^i \binom{p}{i} \\
            & \le 9^{s^*} \left(\frac{9p}{s}\right)^{s} \left(1 + \frac{s}{p}\right)^p \le 9^{s^*} \left(\frac{9ep}{s}\right)^{s}.
\end{split}
\end{align} 

At the same time, denote $\bxi = \sum_{e\in \mathcal{E}} \omega^{(e)} \bSigma^{(e)}_{S\cup S^*,S} \left\{\hat{\mathbb{E}}[\bx_S^{(e)} \varepsilon^{(e)}] - \mathbb{E}[\bx^{(e)} \varepsilon^{(e)}]\right\}$. For any $S\in [p]$ with $|S| \le s$, it follows from the variational representation of the $\ell_2$ norm that
\begin{align*}
\|\bxi\|_2 &= \sup_{\bv \in \mathcal{B}(S \cup S^*)} \bv^\top \bxi = \sup_{k \in [N_S]} (\bv_k^{(S)})^\top \bxi + \sup_{\bv \in \mathcal{B}(S\cup S^*)} (\bv - \bv_{\pi(v)}^{(S)})^\top \bxi \le \sup_{k \in [N_S]} (\bv_k^{(S)})^\top \bxi + \frac{1}{4}\|\bxi\|_2,
\end{align*} where the last inequality follows from the Cauchy-Schwarz inequality and our construction of covering in \eqref{eq:proof-cover1}. This implies $\|\bxi\|_2 \le 2 \sup_{k \in [N_S]} (\bv_k^{(S)})^\top \bxi$, thus
\begin{align}
\label{eq:proof-sup1}
\begin{split}
	&\sup_{|S| \le s} \left\|\sum_{e\in \mathcal{E}} \omega^{(e)} \bSigma^{(e)}_{S\cup S^*,S} \left\{\hat{\mathbb{E}}[\bx_S^{(e)} \varepsilon^{(e)}] - \mathbb{E}[\bx^{(e)} \varepsilon^{(e)}]\right\}\right\|_2  \\
	&~~~~~~ \le 2 \sup_{|S|\le s, k\in [N_S]} \underbrace{(\bv_{k}^{(S)})^\top \sum_{e\in \mathcal{E}} \omega^{(e)} \bSigma^{(e)}_{S\cup S^*,S} \left\{\hat{\mathbb{E}}[\bx_S^{(e)} \varepsilon^{(e)}] - \mathbb{E}[\bx^{(e)} \varepsilon^{(e)}]\right\}}_{Z(S,k)}. 
\end{split}
\end{align} 
For given fixed $\bv_k^{(S)} \in \mathcal{B}(S\cup S^*)$, $Z(S,k)$ can be written as the sum of independent zero-mean random variables as
\begin{align*}
	\sum_{e\in \mathcal{E}} \sum_{i=1}^{n^{(e)}} \frac{\omega^{(e)}}{n^{(e)}}\left(X_{e,i} - \mathbb{E}[X_{e,i}] \right) ~~~~ with ~~~~ X_{e,i} = \left((\bv^{(S)}_k)^\top \bSigma^{(e)}_{S\cup S^*,S} \bx_S^{(e)}\right) \left( \varepsilon^{(e)} \right).
\end{align*} Observe that $\varepsilon^{(e)}$ is a zero-mean sub-Gaussian random variable with parameter $\sigma_\varepsilon$ by \cref{cond3-sub-gaussian-eps}, and $(\bv^{(S)}_k)^\top \bSigma^{(e)}_{S\cup S^*, S}\bx_{S}$ is a zero-mean sub-Gaussian random variable with parameter 
\begin{align*}
	\sigma_1 = \|(\bv_{k}^{(S)})^\top \bSigma^{(e)}_{S\cup S^*, S} \bar{\bSigma}_{S,:}^{1/2} \|_2 \sigma_x \le \kappa_U^{3/2} \sigma_x
\end{align*} by \cref{cond2-subgaussain-x}. It then follows from \cref{lemma:product-sub-gaussian} and \cref{lemma:sum-sub-exp} that there exists some universal constant $C'$ such that
\begin{align*}
    \mathbb{P}\left[ |Z(S,k)| \ge C' \sigma_1 \sigma_\varepsilon \left\{\sqrt{\sum_{e\in \mathcal{E}} \sum_{i=1}^{n^{(e)}}  \left(\frac{\omega^{(e)}}{n^{(e)}}\right)^2 u} + \max_{e\in \mathcal{E}} \frac{\omega^{(e)}}{n^{(e)}} u\right\}\right] \le 2e^{-u}
\end{align*} for any $u>0$. Letting $u=t+\log(2N) \le 3 \left(t + s\log(ep/s) + s^*\right)$, we obtain
\begin{align*}
    \mathbb{P}\left[\sup_{|S|\le s, k} |Z(S,k)| \ge 3C' \sigma_1\sigma_\varepsilon \Bigg(\sqrt{\frac{V(s,t)}{n_{\bomega}}} + \frac{V(s,t)}{n_*}\Bigg)\right] \le  N\times 2e^{-\log(2N)-t} \le e^{-t}.
\end{align*} Combining with the argument \eqref{eq:proof-sup1} concludes the proof of the claim with $C_1=6C'$.


%%%%%%%%%%%%%%%%%%%%%%%%%%%%%%%%%%%% Step 2 
\noindent {\sc Step 2. Upper Bound on $\mathsf{T}_2^{(e)}$. } We claim that $\mathbb{P}(\mathcal{C}_{2,t}) \ge 1-e^{-t}$ for any $t>0$, where
\begin{align}
\label{eq:low-dim-id-bound2}
\begin{split}
\mathcal{C}_{2,t} = \Bigg\{\forall \bbeta, ~~~~ \sum_{e\in \mathcal{E}} \omega^{(e)} \mathsf{T}_2^{(e)}  \le &~C_2 \kappa_U^{1/2} \sigma_x \sigma_\varepsilon \sqrt{\frac{V(|S\setminus S^*|,t)}{n_*}} \times \sqrt{\sum_{e\in \mathcal{E}} \omega^{(e)} \left\|\mathbb{E} [\bx_S^{(e)} \varepsilon^{(e)}]\right\|_2^2} \\
	&~~~~ + C_2 \kappa_U \sigma_x \sigma_\varepsilon^2 \frac{V(|S\setminus S^*|,t)}{n_*} \Bigg\}
\end{split}
\end{align} for some universal constant $C_2$ to be determined. Note that L.H.S. and R.H.S. of the inequality in \eqref{eq:low-dim-id-bound2} both depend on $S$, which is the support set of $\bbeta$ and satisfies $|S|\le s$. For a fixed $S$, we can write down $\sum_{e\in \mathcal{E}} \omega^{(e)} \mathsf{T}_2^{(e)}$ as sum of independent random variables as
\begin{align*}
	Z(S)=\sum_{e\in \mathcal{E}} \omega^{(e)} \mathsf{T}_2^{(e)} = \sum_{e\in \mathcal{E}} \sum_{i=1}^{n^{(e)}} \frac{\omega^{(e)}}{n^{(e)}} \left(X_{e,i} - \mathbb{E}[X_{e,i}]\right) ~~~~\text{with}~~~~ X_{e,i} = \left(\mathbb{E}[\bx_S^{(e)} \varepsilon^{(e)}]^\top [\bx_{i}^{(e)}]_S \right) (\varepsilon^{(e)}_{i}).
\end{align*}
Observe that $X_{i,e}$ are independent sub-exponential random variables because of \cref{lemma:product-sub-gaussian} and our assumptions \cref{cond2-subgaussain-x}--\ref{cond3-sub-gaussian-eps}. It then follows from \cref{lemma:sum-sub-exp} that the following event,
\begin{align*}
	|Z(S)| &\le C'\kappa_U \sigma_x \sigma_\varepsilon \left\{\sqrt{\sum_{e\in \mathcal{E}} (\omega^{(e)})^2 \frac{1}{n^{(e)}} \left\|\mathbb{E}[\bx_S^{(e)} \varepsilon^{(e)}]\right\|_2^2} \times \sqrt{u} + \max_{e\in \mathcal{E}} \frac{\omega^{(e)}}{n^{(e)}} \left\|\mathbb{E}[\bx_S^{(e)}\varepsilon^{(e)}]\right\|_2^2 \times u\right\}\\
	&\le C'\kappa_U^{1/2} \sigma_x \sigma_\varepsilon \left\{ \sqrt{u\times \max_{e'\in \mathcal{E}} \frac{\omega^{(e')}}{n^{(e')}}} \sqrt{\sum_{e\in \mathcal{E}} \omega^{(e)} \left\|\mathbb{E}[\bx_S^{(e)} \varepsilon^{(e)}]\right\|_2^2} + \frac{u}{n_*} \max_{e\in \mathcal{E}}  \left\|\mathbb{E}[\bx^{(e)}\varepsilon^{(e)}]\right\|_2^2 \right\} \\
	&\le C'\kappa_U^{1/2} \sigma_x \sigma_\varepsilon \left\{ \sqrt{\frac{x}{n_*}} \sqrt{\sum_{e\in \mathcal{E}} \omega^{(e)} \left\|\mathbb{E}[\bx_S^{(e)} \varepsilon^{(e)}]\right\|_2^2}  + \kappa_U^{1/2} \sigma_\varepsilon \frac{x}{n_*} \right\}
\end{align*} occurs with probability at least $1-2e^{-x}$ for any $u>0$, where the last inequality follows from \cref{lemma:x-product-epsilon} and the definition of $n_*$ in \eqref{eq:def-n-low-dim}. Now, define the event 
\begin{align*}
\mathcal{K}_u(r) = \left\{ \forall S, |S\setminus S^*|=r, ~~~ |Z(S)| \le 2C'\kappa_U^{1/2} \sigma_x \sigma_\varepsilon \left\{ \sqrt{\frac{V(r,u)}{n_*}} \sqrt{\sum_{e\in \mathcal{E}} \omega^{(e)} \left\|\mathbb{E}[\bx_S^{(e)} \varepsilon^{(e)}]\right\|_2^2} + \kappa_U^{1/2} \sigma_\varepsilon \frac{V(r,u)}{n_*} \right\}\right\}
\end{align*} for any $r\ge 1$. The total number of $S$ satisfying $|S\setminus S^*|=r$ can be upper-bounded by
\begin{align*}
    N_r = 2^{s^*} \times \binom{p-s^*}{r} \le 2^{s^*} (ep/r)^r.
\end{align*} Applying union bound with $x=u+\log(2N_r) \le 2(u + r\log(ep/r) + s^*)$ then gives $\mathbb{P}[\mathcal{K}_u(r)] \ge 1-2N_r e^{-(u+\log(2N_r))} = 1-e^{-u}$. Therefore, we can argue that, under $\bigcap_{r=1}^p \mathcal{K}_{t+\log p}(r)$, the following holds
\begin{align*}
\forall \bbeta, ~~~ Z(S) \le 4C'\kappa_U^{1/2} \sigma_x \sigma_\varepsilon \left\{ \sqrt{\frac{V(|S\setminus S^*|,t)}{n_*}} \sqrt{\sum_{e\in \mathcal{E}} \omega^{(e)} \left\|\mathbb{E}[\bx_S^{(e)} \varepsilon^{(e)}]\right\|_2^2} + \kappa_U^{1/2} \sigma_\varepsilon \frac{V(|S\setminus S^*|,t)}{n_*} \right\}
\end{align*} by the fact that
\begin{align*}
    \forall r \ge 1, ~~~~~ V(r, t + \log p) = r\log(ep/r) + s^* + (\log p) + t \le 2r\log(ep/r) + s^* + t \le 2V(r, t).
\end{align*} This completes the proof of via setting $C_2 = 4C'$.

%%%%%%%%%%%%%%%%%%%%%%%%%%%%%%%%%%%% Step 23
\noindent {\sc Step 3. Upper Bound on $\mathsf{T}_3^{(e)}$.} We argue in this step that the event
\begin{align}
\label{eq:low-dim-id-bound3}
	\mathcal{C}_{3,t} = \left\{ \forall \bbeta \in \mathcal{B}_s, ~~~~ \sum_{e\in \mathcal{E}} \omega^{(e)} \mathsf{T}_3^{(e)} \le C_3\kappa_U^2 \sigma_x^2  \|\bbeta - \bbeta^*\|_2^2 \Bigg(\sqrt{\frac{V(s,t)}{n_{\bomega}}} + \frac{V(s,t)}{n_*}\Bigg) \right\}	
\end{align} occurs with probability at least $1-e^{-t}$ for any $t>0$, where $C_3$ is some universal constant to be determined. Without loss of generality, let $\bbeta \neq \bbeta^*$, then it suffices to establish an upper bound for
\begin{align}
\label{eq:step3-sup-argument-1}
\begin{split}
&~\sup_{\bbeta \neq \bbeta^*} \frac{(\bbeta - \bbeta^*)^\top}{\|\bbeta - \bbeta^*\|_2} \sum_{e\in \mathcal{E}} \omega^{(e)} \bSigma_{:,S}^{(e)} \left(\hat{\mathbb{E}}[\bx_S^{(e)} (\bx^{(e)})^\top] - \mathbb{E}[\bx_S^{(e)} (\bx^{(e)})^\top] \right) \frac{(\bbeta - \bbeta^*)}{\|\bbeta - \bbeta^*\|_2^2} \\
&~~~\le \sup_{|S| \le s} \left\| \sum_{e\in \mathcal{E}} \omega^{(e)} \bSigma_{S\cup S^*,S}^{(e)} \left(\hat{\mathbb{E}}[\bx_S^{(e)} (\bx^{(e)}_{S\cup S^*})^\top] - \mathbb{E}[\bx_S^{(e)} (\bx^{(e)}_{S\cup S^*})^\top] \right) \right\|_2 = \sup_{|S| \le s} \|\bA_S\|_2.
\end{split}
\end{align}
We follow a similar strategy as {\sc Step 1}. For any $S\in \mathcal{B}_s$, let $\{(\bv^{(e)}_{k}, \bu^{(e)}_{k})\}_{k=1}^{N_S} \in \mathcal{B}(S\cup S^*) \times \mathcal{B}(S\cup S^*) := \mathcal{B}^2(S\cup S^*)$ be a $1/4$-covering of $\mathcal{B}^2(S\cup S^*)$ in a sense that for any $(\bu, \bv) \in \mathcal{B}^2(S\cup S^*)$, there exists some $\pi(\bu, \bv) \in [N_S]$ such that
\begin{align*}
    \|\bu - \bu^{(S)}_{\pi(\bu, \bv)}\|_{2} + \|\bv - \bv^{(S)}_{\pi(\bu, \bv)}\|_{2} \le \frac{1}{4}.
\end{align*} It follows from standard empirical process theory that $N_S \le 9^{2|S\cup S^*|}$, then
\begin{align*}
    N = \sum_{S\subseteq [p], |S| \le s} N_S &\le \sum_{S\subseteq [p], |S| \le s} N_S 9^{2|S\cup S^*|} \le \sum_{i=0}^{s} 81^{i + s^*} \binom{p}{i}
    \le 81^{s^*} \left(\frac{81ep}{s}\right)^s,
\end{align*} where the last inequality follows from the same procedure as \eqref{eq:calculate-cover-number}.

At the same time, denote $\bu^\dagger = \bu^{(S)}_{\pi(\bu,\bv)}$ and $\bv^\dagger = \bv^{(S)}_{\pi(\bu,\bv)}$. It follows from the variational representation of the matrix $\ell_2$ norm that
\begin{align*}
    \|\bA_S\|_2 &= \sup_{(\bu, \bv) \in \mathcal{B}^2(S\cup S^*)} \bu^\top \bA_S \bv \\
    &\le \sup_{(\bu, \bv) \in \mathcal{B}^2(S\cup S^*)} (\bu^\dagger)^\top \bA_S \bv^\dagger + \sup_{(\bu, \bv) \in \mathcal{B}^2(S\cup S^*)} (\bu-\bu^\dagger)^\top \bA_S \bv^\dagger \\
    &~~~~~~~~ + \sup_{(\bu, \bv) \in \mathcal{B}^2(S\cup S^*)} (\bu-\bu^\dagger)^\top \bA_S (\bv - \bv^\dagger) + \sup_{(\bu, \bv) \in \mathcal{B}^2(S\cup S^*)} (\bu^\dagger)^\top \bA_S (\bv - \bv^\dagger) \\
    &\le  \sup_{k \in [N_S]} (\bu^{(S)}_k)^\top \bA_S \bv^{(S)}_k + \left(\frac{1}{4} + \frac{1}{4} + \frac{1}{16} \right) \|\bA_S\|_2, 
\end{align*} which implies $\|\bA_S\|_2 \le 4 \sup_{k \in [N_S]} (\bu^{(S)}_k)^\top \bA_S \bv^{(S)}_k$, thus
\begin{align}
\label{eq:step3-sup-argument-2}
    \sup_{|S| \le s} \|\bA_S\|_2 \le \sup_{|S| \le s, k \in [N_S]} (\bu_{k}^{(e)})^\top \bA_S (\bv_k^{(e)}) = \sup_{|S| \le s, k\in [N_S]} Z(S,k).
\end{align} For fixed $k$ and $S$, $Z(S,k)$ can be written as the sum of independent zero-mean random variables as
\begin{align*}
    Z(S,k) = \sum_{e\in \mathcal{E}} \sum_{i=1}^{n^{(e)}} \frac{\omega^{(e)}}{n^{(e)}} (X_{e,i} - \mathbb{E} [X_{e,i}]) \qquad \text{with} \qquad X_{e,i} = \left((\bu_{k}^{(e)})^\top \Sigma^{(e)}_{S\cup S^*, S} \bx_S^{(e)} \right) \left((\bx_{S\cup S^*}^{(e)})^\top \bv_k^{(e)}\right).
\end{align*} Here $X_{e,i}$ is the product of two zero-mean sub-Gaussian random variables with parameter $\kappa_U^{3/2}\sigma_x$ and $\kappa_U^{1/2}\sigma_x$, respectively. It then follows from \cref{lemma:product-sub-gaussian} and \cref{lemma:sum-sub-exp} that, for any $u>0$,
\begin{align*}
    \mathbb{P}\left[|Z(S,k)| \le C'\kappa_U^2 \sigma_x^2 \left(\sqrt{\frac{u}{n_{\bomega}}} + \frac{u}{n_*}\right)\right] \ge 1-2e^{-u},
\end{align*} where $C'$ is some universal constant. Finally, we apply union bound with $u=t+\log(2N) \le 6(t + s\log(ep/s) + s^*)$ and obtain
\begin{align*}
    \mathbb{P}\left[\sup_{|S| \le s, k\in [N_S]} |Z(S,k)| \le 6C'\kappa_U^2 \sigma_x^2 \Bigg(\sqrt{\frac{V(s,t)}{n_{\bomega}}} + \frac{V(s,t)}{n_*}\Bigg)\right] \ge 1-N\times 2e^{-\log(2N) + t} = 1-e^{-t}.
\end{align*} Set $C_3 = 24C'$. Combining the above inequality with the suprema arguments \eqref{eq:step3-sup-argument-2} and \eqref{eq:step3-sup-argument-1} completes the proof.



\noindent {\sc Step 4. Upper Bound on $\mathsf{T}_4^{(e)}$.} Our target in this step is to show that, for any $t>0$,
\begin{align}
\label{eq:low-dim-id-bound4}
\mathbb{P}(\mathcal{C}_{4,t}) = \mathbb{P}\left[\forall \bbeta \in \mathcal{B}_s, ~~ \sum_{e\in \mathcal{E}} \omega^{(e)} \mathsf{T}_4^{(e)} \le C_4 \kappa_U^{3/2} \sigma_x^2\sigma_\varepsilon \|\bbeta - \bbeta^*\|_2\Bigg(\sqrt{\frac{V(s,t)}{n_{\bomega}}} + \frac{V(s,t)}{n_*}\Bigg)\right] \ge 1-e^{-t}.
\end{align} The proof is very similar to that in {\sc Step 1}, we only sketch here and highlight the difference. Following a similar strategy, it suffices to derive an upper bound for the quantity
\begin{align*}
    \sup_{|S| \le s, k\in [N_S]} Z(S,k) ~~~~~ \text{with} ~~~~~ Z(S,k) = \sum_{e\in \mathcal{E}} \sum_{i=1}^{n^{(e)}} \frac{\omega^{(e)}}{n^{(e)}} (X_{e,i} - \mathbb{E}[X_{e,i}]) ~~\text{and}~~ N=\sum_{S\subseteq [p]} N_S \le 9^{s^*}(9ep/s)^s,
\end{align*} where $X_{e,i}=\left((\mathbb{E}[\bx_S^{(e)} \varepsilon^{(e)}])^\top [\bx^{(e)}_i]_S\right) \left([\bx_i^{(e)}]_{S\cup S^*}^\top \bv_k^{(S)}\right)$ is the product of two zero-mean sub-Gaussian random variables with parameters $\kappa_U\sigma_x\sigma_\varepsilon$ and $\kappa_{U}^{1/2}\sigma_x$, respectively. It then follows from \cref{lemma:sum-sub-exp} that, for any $u>0$,
\begin{align*}
    \mathbb{P}\left[|Z(S,k)| \le C' \kappa_U^{3/2} \sigma_x^2\sigma_\varepsilon \left(\sqrt{\frac{u}{n_{\bomega}}} + \frac{u}{n_*}\right)\right] \ge 1-2e^{-u}.
\end{align*} So it concludes via applying union bound with $u = t + \log (2N) \le 3V(s,t)$.

\noindent {\sc Step 5. Upper Bound on $\mathsf{T}_5^{(e)}$.} In this step, we claim that the following event
\begin{align}
\label{eq:low-dim-id-bound5}
    \mathcal{C}_{5,t} = \left\{\forall \bbeta \in \mathbb{R}^p, ~~ \sum_{e\in \mathcal{E}} \omega^{(e)} \mathsf{T}_5^{(e)} \le C_5\kappa_U^{3/2} \sigma_x^3 \sigma_\varepsilon \frac{p + \log(2|\mathcal{E}|)+t}{\bar{n}}\|\bbeta - \bbeta^*\|_2\right\}
\end{align} occurs with probability at least $1-e^{-t}$ if $s^*+t+\log(|2\mathcal{E}|) \le n_{\min}$. Define the event $\mathcal{K}_{5,u}(e)$ as
\begin{align*}
    \mathcal{K}_{5,u}(e) = \left\{\left\|\hat{\mathbb{E}}[\bx_{S^*}^{(e)}\varepsilon^{(e)}]\right\|_2 \le C_{1}'\kappa_U^{1/2} \sigma_x \sigma_\varepsilon \Bigg(\sqrt{\frac{s^* + u}{n^{(e)}}} + \frac{s^* + u}{n^{(e)}}\Bigg)\right\}
\end{align*} for some universal constant $C_1'$. Observe that $\|\mathbb{E}[\bx_{S^*}^{(e)}\varepsilon^{(e)}]\|_2 = \sup_{\bv\in \mathbb{R}^{s^*}, \|\bv\|_2=1} \bv^\top \hat{\mathbb{E}}[\bx_{S^*}^{(e)}\varepsilon^{(e)}]$,
and $\bv^\top \hat{\mathbb{E}}[\bx_S^{(e)}\varepsilon^{(e)}]$ with given fixed $\bv\in \mathbb{R}^{s^*}$ is the sum of independent (centered) products of two zero-mean sub-Gaussian random variables with parameter $\kappa_U^{1/2} \sigma_x$ and $\sigma_\varepsilon$, respectively. Then it follows from \cref{lemma:product-sub-gaussian} and \cref{lemma:sum-sub-exp} that, for any fixed $\bv\in \mathbb{R}^{s^*}$ and $x>0$
\begin{align*}
    \mathbb{P}\left[ \left|\bv^\top \hat{\mathbb{E}}[\bx_{S^*}^{(e)}\varepsilon^{(e)}] \right| \ge C'\kappa_U^{1/2}\sigma_x \sigma_\varepsilon\Bigg(\sqrt{\frac{x}{n^{(e)}}} + \frac{x}{n^{(e)}}\Bigg)\right] \le 2e^{-x}.
\end{align*} Following a similar argument as that of {\sc Step 1} and applying union bound gives $\mathbb{P}\left[\mathcal{C}_{5,u}(e)\right] \ge 1-e^{-u}$. 

At the same time, let 
\begin{align*}
    \mathcal{K}_{5,u}'(e) = \left\{ \forall \bbeta \in \mathcal{B}_s, ~~ \left\| \hat{\mathbb{E}}[\bx_{S^*}^{(e)} (\bx^{(e)})^\top] - \bSigma^{(e)}_{S^*,:} \right\|_2 \le C_2'\kappa_U \sigma_x^2 \Bigg(\sqrt{\frac{V(s,u)}{n^{(e)}}} + \frac{V(s,u)}{n^{(e)}}\Bigg)\right\}
\end{align*} for some universal constant $C_2'>0$. It is easy to verify that $\mathbb{P}(\mathcal{C}_{5,u}'(e)) \ge 1-e^{-u}$ for any given fixed $u>0$ and $e\in \mathcal{E}$. 

Under the event $\mathcal{K}_u = \bigcap_{e\in \mathcal{E}} \{\mathcal{K}_{5,u}(e) \cap \mathcal{K}'_{5,u}(e)\}$, which occurs with probability $1-2|\mathcal{E}| e^{-u}$, we obtain
\begin{align*}
    \forall \bbeta \in \mathcal{B}_s, ~~\sum_{e\in \mathcal{E}} \omega^{(e)} \mathsf{T}_5^{(e)} &= \sum_{e\in \mathcal{E}} \omega^{(e)} \left(\hat{\mathbb{E}}[\bx_I^{(e)}\varepsilon^{(e)}]\right)^\top \left(\hat{\mathbb{E}}[\bx_{I}^{(e)} (\bx^{(e)})^\top] - \bSigma^{(e)}_{I,:}\right) \left(\bbeta - \bbeta^*\right) \\
    &\le \sum_{e\in \mathcal{E}} \omega^{(e)} \left\|\hat{\mathbb{E}}[\bx_I^{(e)}\varepsilon^{(e)}]\right\|_2 \left\|\hat{\mathbb{E}}[\bx_{I}^{(e)} (\bx^{(e)})^\top] - \bSigma^{(e)}_{I,:}\right\|_2 \|\bbeta - \bbeta^*\|_2 \\
    &\le \sum_{e\in \mathcal{E}} \omega^{(e)} \left\|\hat{\mathbb{E}}[\bx_{S^*}^{(e)}\varepsilon^{(e)}]\right\|_2 \left\|\hat{\mathbb{E}}[\bx_{S^*}^{(e)} (\bx^{(e)})^\top] - \bSigma^{(e)}_{S^*,:}\right\|_2 \|\bbeta - \bbeta^*\|_2 \\ 
    &\overset{(a)}{\le} C'' \kappa_U^{3/2} \sigma_x^3 \sigma_\varepsilon \sum_{e\in \mathcal{E}} \omega^{(e)} \sqrt{\frac{|S^*| + u}{n^{(e)}}} \Bigg(\sqrt{\frac{V(s, u)}{n^{(e)}}} + \frac{V(s,u)}{n^{(e)}} \Bigg) ~ \|\bbeta - \bbeta^*\|_2 \\
    &\overset{(b)}{\le} C'' \kappa_U^{3/2} \sigma_x^3 \sigma_\varepsilon \frac{\sqrt{V(s,u) V(0,u)}}{\bar{n}} \|\bbeta - \bbeta^*\|_2 \\
    &~~~~~~~~ + C'' \kappa_U^{3/2} \sigma_x^3 \sigma_\varepsilon \frac{V(s,u)\sqrt{V(0,u)}}{n_\dagger} \|\bbeta - \bbeta^*\|_2
\end{align*} provided $s^* + u \le n_{\min}$. Here $(a)$ follows from the definition of the events $\mathcal{K}_{5,u}(e)$ and $\mathcal{K}_{5,u}'(e)$ and the fact that $x \le \sqrt{x}$ when $x\in [0,1]$, $(b)$ follows directly from and the definition of $\bar{n}$ in \eqref{eq:def-n-low-dim} and the definition of $n_\dagger$ in \eqref{eq:def-n-low-dim}. This completes the proof via letting $u = \log(2|\mathcal{E}|) + t$.


\noindent {\sc Step 6. Upper Bound on $\mathsf{T}_6^{(e)}$.} The goal of this step is to derive a high-probability bound for the event
\begin{align}
\label{eq:low-dim-id-bound6}
    \mathcal{C}_{6,t} = \left\{\forall \bbeta \in \mathbb{R}^p, ~~ \sum_{e\in \mathcal{E}} \omega^{(e)} \mathsf{T}_6^{(e)} \le C_6 \kappa_U \sigma_x^2\sigma_\varepsilon^2  \frac{t+ \log(2|\mathcal{E}||S^*|)}{\bar{n}} |S^*\setminus S| \right\}
\end{align} for any $t \in (0, n_{\min} - \log(2|\mathcal{E}||S^*|)]$. Note that both the L.H.S. and R.H.S. of the inequality in \eqref{eq:low-dim-id-bound6} depends on $\bbeta$, or more precisely, $S=\supp(\bbeta)$. Denoting $\delta = C_6 \kappa_U\sigma_x^2\sigma_\varepsilon^2 \{t+ \log(2|\mathcal{E}||S^*|)\}/\bar{n}$, we have the following decomposition
\begin{align*}
    \mathcal{C}_{6,t} &= \bigcup_{T\subseteq S^*} \left\{\forall \bbeta \in \mathbb{R}^p, S^*\setminus \supp(\bbeta) = T, ~~ \sum_{e\in \mathcal{E}} \omega^{(e)} \mathsf{T}_6^{(e)} \le \delta |T| \right\}\\
    &= \bigcup_{T\subseteq S^*} \left\{\forall \bbeta \in \mathbb{R}^p, S^*\setminus \supp(\bbeta) = T, ~~ \sum_{j\in T} \sum_{e\in \mathcal{E}} \omega^{(e)} \left|\hat{\mathbb{E}}[x_j^{(e)}\varepsilon^{(e)}]\right|^2 \le \delta |T| \right\} \\
    &= \bigcup_{T\subseteq S^*} \mathcal{K}(T).
\end{align*}

At the same time, given fixed $j\in S^*$ and $e\in \mathcal{E}$, it follows from \cref{lemma:sum-sub-exp} that,
\begin{align*}
    \mathbb{P}[\mathcal{K}_{6,x}(e, j)] = \mathbb{P}\left[ \left|\hat{\mathbb{E}}[x_j^{(e)} \varepsilon^{(e)}]\right| \le C'\kappa_U^{1/2}\sigma_x\sigma_\epsilon \Bigg(\sqrt{\frac{x}{n^{(e)}}} + \frac{x}{n^{(e)}}\Bigg)\right] \ge 1- 2e^{-x}.
\end{align*} for some universal constant $C'$. We claim that
\begin{align}
\label{eq:low-dim-id-step6-claim1}
    \mathcal{K}(T) \subseteq \bigcup_{j\in S^*, e\in \mathcal{E}} \mathcal{K}_{6,t + \log(2s^*|\mathcal{E}|)}(e,j) 
\end{align} by choosing $C_6=(C')^2$. This is because under $\bigcup_{j\in S^*, e\in \mathcal{E}} \mathcal{K}_{6,t + \log(2s^*|\mathcal{E}|)}(e,j)$, one has
\begin{align*}
    \sum_{j\in T} \sum_{e\in \mathcal{E}} \omega^{(e)} \left|\hat{\mathbb{E}}[x_j^{(e)}\varepsilon^{(e)}]\right|^2 &\le \sum_{j\in T} \sum_{e\in \mathcal{E}} (C')^2 \kappa_U \sigma_x^2\sigma_\varepsilon^2\left(\frac{t + \log(2s^*|\mathcal{E}|)}{n^{(e)}} \omega^{(e)}\right) \\
    &\le |T| (C')^2 \kappa_U \sigma_x^2\sigma_\varepsilon^2\frac{t + \log(2s^*|\mathcal{E}|)}{\bar{n}}
\end{align*} provided $t+\log(2s^*|\mathcal{E}|) \le n_{\min}$, this validates the claim \eqref{eq:low-dim-id-step6-claim1}. Therefore, we have
\begin{align*}
    \mathbb{P}(\mathcal{C}_{6,t}) = \mathbb{P} \left[\bigcup_{T\subseteq S^*} \mathcal{K}(T)\right] &\ge \mathbb{P} \left[ \bigcup_{j\in S^*, e\in \mathcal{E}} \mathcal{K}_{6,t + \log(2s^*|\mathcal{E}|)}(e,j) \right]\\
    &\ge 1 - \sum_{e\in \mathcal{E}, j\in S^*} \left(1-\mathbb{P}\left[\mathcal{K}_{6,t + \log(2s^*|\mathcal{E}|)}(e,j)\right]\right) \\
    &\ge 1-2s^*|\mathcal{E}|e^{-t - \log(2s^*|\mathcal{E}|)} \ge 1-e^{-t}.
\end{align*}

\noindent {\sc Step 7. Conclusion.} We are now ready to conclude the proof by combining results \eqref{eq:low-dim-id-bound1}, \eqref{eq:low-dim-id-bound2}, \eqref{eq:low-dim-id-bound3}, \eqref{eq:low-dim-id-bound4}, \eqref{eq:low-dim-id-bound5}, \eqref{eq:low-dim-id-bound6} we obtained from {\sc Step 1} to {\sc Step 6}. Plugging these upper bounds back into our decomposition in \eqref{eq:low-dim-id-decomposition}, we have that, under the event $\bigcap_{k=1}^6\mathcal{C}_{k,t}$, which occurs with probability at least $1-6e^{-t}$, the inequality in \eqref{eq:instance-dependent-bound-j-highdim} holds provided $n_{\min} \ge (s^* + \log(2|\mathcal{E}|) + t) \lor (\log(2|\mathcal{E}||S^*|) + t) = s^* + \log(2|\mathcal{E}|) + t$. This completes the proof. \qed

\subsection{Proof of \cref{lemma:instance-dependent-bound-r-highdim}}
\label{subsec:proof-instance-dependent-bound-r-highdim}

It follows from the definition of the pooled $L_2$ risk that
\begin{align*}
&~\mathsf{R}(\bbeta) - \mathsf{R}(\bbeta^*) - \hat{\mathsf{R}}(\bbeta) + \hat{\mathsf{R}}(\bbeta^*) \\
&~~~~~~~  = \sum_{e\in \mathcal{E}} \omega^{(e)} \left\{\mathsf{R}^{(e)}(\bbeta) - \mathsf{R}^{(e)}(\bbeta^*) - \left(\hat{\mathsf{R}}^{(e)}(\bbeta) - \hat{\mathsf{R}}^{(e)}(\bbeta^*)\right) \right\}\\
&~~~~~~~  = \sum_{e\in \mathcal{E}} \omega^{(e)} (\bbeta - \bbeta^*)^\top \left(\bSigma^{(e)} - \hat{\bSigma}^{(e)}\right) (\bbeta - \bbeta^*) - 2(\bbeta - \bbeta^*)^\top \left(\mathbb{E}[\bx^{(e)}\varepsilon^{(e)}] - \hat{\mathbb{E}}[\bx^{(e)} \varepsilon^{(e)}]\right) \\
&~~~~~~~  = (\bbeta - \bbeta^*)^\top \bA (\bbeta - \bbeta^*) - 2(\bbeta - \bbeta^*)^\top \bb.
\end{align*}

Recall that $\delta_1 = \sqrt{\frac{V(s,t)}{n_{\bomega}}} + \frac{V(s,t)}{n_*}$, we argue that the following two events
\begin{align*}
    \mathcal{C}_{1,t} &= \left\{ \sup_{(\bx, \by)\in \mathcal{B}_s \times \mathcal{B}_s, \|\bx\|_2=\|\by\|_2=1} \bx^\top \bA \by \le C_1 \kappa_U \sigma_x^2 \delta_1\right\} \\
    \qquad \mathcal{C}_{2,t} &= \left\{ \sup_{\bx \in \mathcal{B}_s, \|\bx\|_2=1} \bx^\top \bb \le C_2 \kappa_U^{1/2} \sigma_x\sigma_\varepsilon \delta_1\right\}
\end{align*} satisfies $\mathbb{P}[\mathcal{C}_{1,t}] \land \mathbb{P}[\mathcal{C}_{2,t}] \ge 1-0.5 e^{-t}$ for any $t>0$, where $C_1, C_2$ are some universal constants to be determined. If the two claims are verified, then we have, under the event $\mathcal{C}_{1,t} \cap \mathcal{C}_{2,t}$ which occurs with probability at least $1-e^{-t}$, the following holds
\begin{align*}
&~|\mathsf{R}(\bbeta) - \mathsf{R}(\bbeta^*) - \hat{\mathsf{R}}(\bbeta) + \hat{\mathsf{R}}(\bbeta^*)| \\
&~~~~~~ \le \left|(\bbeta - \bbeta^*)^\top \bA (\bbeta - \bbeta^*)\right| + 2\left|(\bbeta - \bbeta^*)^\top \bb \right| \\
&~~~~~~ \le \{C_1 \lor (2C_2) \}\left(\kappa_U \sigma_x^2 \delta_1 \|\bbeta - \bbeta^*\|_2^2 + \kappa_U^{1/2} \sigma_x\sigma_\varepsilon \delta_1 \|\bbeta - \bbeta^*\|_2\right)
\end{align*} for some universal constant $C'$. This completes the proof of \cref{lemma:instance-dependent-bound-r-lowdim}.

Now, we prove the two claims separately.

\noindent \emph{Proof of the Claim $\mathbb{P}[\mathcal{C}_{1,t}]\ge 1-0.5e^{-t}$. } The proof strategy is very similar to {\sc Step 3} in the proof of \cref{lemma:instance-dependent-bound-j-lowdim}. To be specific, similar to the derivation in \eqref{eq:step3-sup-argument-2}, there exist $N \le 81^{s^*} (81ep/s)^s$ pairs of $p$-dimensional unit vectors $(\bv_1, \bu_1),\ldots, (\bv_N, \bu_N)$ such that
\begin{align*}
    \sup_{(\bx, \by)\in \mathcal{B}_s \times \mathcal{B}_s, \|\bx\|_2=\|\by\|_2=1} \bx^\top \bA \by \le 4\sup_{k\in [N]} \bv_k^\top \bA \bu_k.
\end{align*}
For fixed $(\bv_k, \bu_k)$, it follows from Conditions \ref{cond0-model}, \ref{cond1-well-condition} and \ref{cond2-subgaussain-x} that $\bv_k^\top \bA \bu_k$ is the sum of independent variables $X_{e,i}-\mathbb{E}[X_{e,i}]$, and $X_{e,i}$ is the product of two zero-mean sub-Gaussian random variables with parameter $\kappa_U^{1/2}\sigma_x$. Then applying \cref{lemma:product-sub-gaussian} and \cref{lemma:sum-sub-exp} gives
\begin{align*}
    \mathbb{P}\left[|\bv_k^\top \bA \bu_k| \ge C' \kappa_U^{1/2} \sigma_x \left(\sqrt{\frac{u}{n_{\bomega}}} + \frac{u}{n_*}\right)\right] \le 2e^{-u}
\end{align*} for any $u>0$ and some universal constant $C'$. Using the union bounds over all the $k$, we have the following event
\begin{align*}
    \sup_{(\bx, \by)\in \mathcal{B}_s \times \mathcal{B}_s} \le 4\sup_{k\in [N]} \bv_k^\top \bA \bu_k \le 4C' \kappa_U^{1/2} \sigma_x \left(\sqrt{\frac{u}{n_{\bomega}}} + \frac{u}{n_*}\right)
\end{align*} will occurs with probability at least $1-2Ne^{-u}$. Letting $u=t+\log(4N) \le 6V(s,t)$ completes the proof.

\noindent \emph{Proof of the Claim $\mathbb{P}[\mathcal{C}_{2,t}]\ge 1-0.5e^{-t}$.} The proof strategy is also similar to the first part of {\sc Step 5} in the proof of \cref{lemma:instance-dependent-bound-j-lowdim}. To be specific, there exist $N=90^p$ $p$-dimensional unit vectors $\bu_1,\ldots, \bu_N$ such that $\sup_{\bx \in \mathcal{B}_s, \|\bx\|_2=1} \bx^\top \bb \le 2\sup_{k\in [N]} \bu_k^\top \bb$. Moreover, for fixed $\bu_k$, it follows from Conditions \ref{cond0-model} -- \ref{cond3-sub-gaussian-eps} such that $\bu_k^\top \bb$ is the sum of independent variables $X_{e,i} - \mathbb{E}[X_{e,i}]$, and $X_{e,i}$ is the product of two zero-mean sub-Gaussian random variables with parameters $\kappa_U^{1/2}\sigma_x$ and $\sigma_\varepsilon$, respectively. Following a similar procedure of applying \cref{lemma:product-sub-gaussian}, \cref{lemma:sum-sub-exp}, and the union bound as above concludes the proof.


\subsection{Proof of \cref{proposition:small-hat-s}}
\label{subsec:proof-small-hat-s}

We need the following lemma. 
\begin{lemma}
\label{lemma:lemma-qhat-beta-star}
    Assume Conditions \ref{cond0-model}--\ref{cond3-sub-gaussian-eps} hold. Define the following event
    \begin{align}
        \mathcal{A}_{5,t} = \left\{\hat{\mathsf{Q}}(\bbeta^*;\gamma,\bomega) \le \sigma_\varepsilon^2 \left( 1 + c_1 \sqrt{\frac{t}{n_*}} + c_1 \kappa_U\sigma_x^2 \gamma \frac{s^*( \log(2s^*|\mathcal{E}|) + t)}{\bar{n}}\right)\right\}
    \end{align} for some universal constant, we have $\mathbb{P}(\mathcal{A}_{5,t}) \ge 1-2e^{-t}$. 
\end{lemma}

\begin{proof}[Proof of \cref{lemma:lemma-qhat-beta-star}]
It follows from the definition that
\begin{align*}
    \hat{\mathsf{Q}}(\bbeta^*) &= \sum_{e\in \mathcal{E}} \omega^{(e)} \left\{\hat{\mathbb{E}} |\varepsilon^{(e)}|^2 + \gamma \sum_{j\in S^*} \left(\hat{\mathbb{E}} [x_j^{(e)} \varepsilon^{(e)}]\right)^2\right\} \\
    &= \sum_{e\in \mathcal{E}} \omega^{(e)} \mathbb{E} |\varepsilon^{(e)}|^2 + \sum_{e\in \mathcal{E}} \omega^{(e)} \left(\hat{\mathbb{E}} |\varepsilon^{(e)}|^2 - {\mathbb{E}} |\varepsilon^{(e)}|^2\right) + \gamma \sum_{e\in \mathcal{E}} \omega^{(e)} \left(\hat{\mathbb{E}} [x_j^{(e)} \varepsilon^{(e)}]\right)^2.
\end{align*} It follows from \cref{cond3-sub-gaussian-eps} that
\begin{align}
\label{eq:proof-qhat-beta-star-c0}
    \sum_{e\in \mathcal{E}} \omega^{(e)} \mathbb{E} |\varepsilon^{(e)}|^2 \le \sum_{e\in \mathcal{E}} \omega^{(e)} \sigma_\varepsilon^2 = \sigma_\varepsilon^2,
\end{align} and for any $t>0$,
\begin{align}
\label{eq:proof-qhat-beta-star-c1}
    \mathbb{P}\left[\sum_{e\in \mathcal{E}} \omega^{(e)} \left(\hat{\mathbb{E}} |\varepsilon^{(e)}|^2 - {\mathbb{E}} |\varepsilon^{(e)}|^2\right) \le C' \sigma_\varepsilon^2 \left(\sqrt{\frac{t}{n_{\bomega}}} + \frac{t}{n_{*}}\right)\right] \ge 1-e^{-x}
\end{align} for some universal constant $C'>0$. At the same time, for any fixed $e\in \mathcal{E}$, it follows from \cref{lemma:product-sub-gaussian} and \cref{lemma:sum-sub-exp} that, for any $x>0$ and $j\in [S^*]$,
\begin{align*}
    \mathbb{P}[\mathcal{K}_{t}(e,j)] = \mathbb{P}\left[ \left|\hat{\mathbb{E}} [x_j^{(e)} \varepsilon^{(e)}] - \mathbb{E}[x_j^{(e)} \varepsilon^{(e)}]\right| \ge C' \kappa_U^{1/2}\sigma_x\sigma_\varepsilon \left(\sqrt{\frac{x}{n^{(e)}}} + \frac{x}{n^{(e)}}\right)\right] \le 1-2e^{-x}.
\end{align*} Note $\mathbb{E}[x_j^{(e)} \varepsilon^{(e)}] = 0$. Then under $\mathcal{K}_t = \bigcap_{j\in S^*, e\in \mathcal{E}} \mathcal{K}_{t+\log(2s^*|\mathcal{E}|)}(e,j)$, which occurs with probability at least $1-e^{-t}$, we have
\begin{align}
\label{eq:proof-qhat-beta-star-c2}
    \sum_{j\in S^*} \sum_{e\in \mathcal{E}} \omega^{(e)} \left(\hat{\mathbb{E}} [x_j^{(e)} \varepsilon^{(e)}] \right)^2 \le (C')^2 \kappa_U \sigma_x^2 \sigma_\varepsilon^2 \frac{s^*(\log(2s^*|\mathcal{E}|) + t)}{\bar{n}}
\end{align} provided $t+\log(2s^*|\mathcal{E}|) \le n_{\min}$. Putting the three bounds \eqref{eq:proof-qhat-beta-star-c0}, \eqref{eq:proof-qhat-beta-star-c1} and \eqref{eq:proof-qhat-beta-star-c2} together completes the proof.
\end{proof}

Now we are ready to prove \cref{proposition:small-hat-s}.
\begin{proof}[Proof of \cref{proposition:small-hat-s}]
Observing that \begin{align*}
    \left\{|\hat{S}| \le 2s^*\right\} \subseteq \left\{ \forall \bbeta~\text{with}~\supp(\bbeta) > 2s^*,~~ \hat{\mathsf{Q}}(\bbeta) + \lambda \|\bbeta\|_0 > \hat{\mathsf{Q}}(\bbeta^*) + \lambda \|\bbeta^*\|_0 \right\} := \mathcal{C}_t,
\end{align*} it remains to show that $\mathbb{P}(\mathcal{C}_t) \ge 1-e^{-t}$. To this end, we use a peeling device. Let $\alpha_\ell=2^{\ell} s^*$, then 
\begin{align*}
    \mathcal{C}_t = \bigcup_{\ell=1}^{\lceil \log_2(p/s^*) \rceil - 1} \left\{\forall \bbeta \in \mathcal{B}_{\alpha_{\ell+1}} \setminus \mathcal{B}_{\alpha_{\ell}}, ~~ \hat{\mathsf{Q}}(\bbeta) + \lambda \|\bbeta\|_0 > \hat{\mathsf{Q}}(\bbeta^*) + \lambda \|\bbeta^*\|_0\right\} = \bigcup_{\ell=1}^{\lceil \log_2(p/s^*) \rceil - 1} \mathcal{C}_{t}(\ell).
\end{align*}
We claim that if $\lambda$ satisfies the conditions presented in the statement, then \begin{align}
\label{eq:proof-small-hats-peeling-argument}
    \bigcup_{\ell=1}^{\lceil \log_2(p/s^*) \rceil - 1} \mathcal{C}_{t}(\ell) \subseteq \mathcal{A}_{5,t} \cup \left\{\bigcup_{\ell=1}^{\lceil \log_2(p/s^*) \rceil - 1} \mathcal{A}_{3,u_{\ell}}(\alpha_{\ell+1} s^*) \cup \mathcal{A}_{4,u_{\ell}}(\alpha_{\ell+1} s^*)\right\}.
\end{align} with $u_\ell = t + \log (\lceil \log_2(p/s^*) \rceil) \le t + \log (ep/s^*)$.
Denote $\phi_* = s^* + \log(ep/s^*)+t$ and
\begin{align*}
    \tilde{s} = \frac{n_*}{2\{\log(ep/s^*) + t\}} \land \frac{\left(\frac{1}{12C_1} \frac{\kappa_L}{\gamma \kappa_U^2\sigma_x^2}\right)^2n_*}{\log(ep/s^*) + t} \land \frac{n_\dagger}{\{\log(ep/s^*)+t\} \sqrt{\phi_*} (\gamma/\kappa_L)\kappa_U^{3/2} \sigma_x^2 },
\end{align*} where $C_1$ is some universal constant to be determined. We prove this claim by considering the following two cases.

\noindent \emph{Case 1, $\alpha_{\ell+1} s^* \le \tilde{s}$.} Let $S=\supp(\bbeta)$. Denote $\phi_s=\alpha_{\ell+1} s^* \log(ep/s^*) + t$. In this case, there exists some universal constant $C_1$ such that if $\mathcal{A}_{3,u_{\ell}}(\alpha_{\ell+1}s^*) \cap \mathcal{A}_{4,u_{\ell}}(\alpha_{\ell+1}s^*)$ occurs, then we have, for any $\bbeta \in \mathcal{B}_{\alpha_{\ell+1}}$,
\begin{align*}
    \hat{\mathsf{Q}}(\bbeta) - \hat{\mathsf{Q}}(\bbeta^*) &= \hat{\mathsf{Q}}(\bbeta) - \mathsf{Q}(\bbeta) + \mathsf{Q}(\bbeta) - \mathsf{Q}(\bbeta^*) + \mathsf{Q}(\bbeta^*) - \hat{\mathsf{Q}}(\bbeta^*) \\
    &\overset{(a)}{\ge} \frac{\kappa_L}{2} \|\bbeta - \bbeta^*\|_2^2 + \frac{\gamma}{6} \mathsf{J}(\bbeta; \bomega) + \hat{\mathsf{R}}(\bbeta) - \mathsf{R}(\bbeta) + \mathsf{R} (\bbeta^*) - \hat{\mathsf{R}} (\bbeta^*) \\
    & ~~~~~~~~~~~~ - \gamma \left(\mathsf{J}(\bbeta) - \hat{\mathsf{J}}(\bbeta) - \mathsf{J}(\bbeta^*) + \hat{\mathsf{J}}(\bbeta^*)\right)\\
    &\overset{(b)}{\ge} \left(\frac{\kappa_L}{2} - C_1 \kappa_U^2\sigma_x^2 \gamma \sqrt{\frac{\phi_s}{n_*}}\right) \|\bbeta - \bbeta^*\|_2^2 \\
    & ~~~~~~~~~~~~ + \frac{\gamma}{6} \mathsf{J}(\bbeta;\bomega) - C_1 \gamma \sqrt{\sum_{e\in \mathcal{E}} \omega^{(e)} \left\|\mathbb{E}[\bx_{S}^{(e)} \varepsilon^{(e)}]\right\|_2^2} \kappa_U^{1/2} \sigma_x\sigma_\varepsilon \sqrt{\frac{\phi_s}{n_*}} \\
    & ~~~~~~~~~~~~ - C_1 \|\bbeta - \bbeta^*\|_2 \times \gamma \kappa_U^{3/2} \sigma_x^2 \sigma_\varepsilon \sqrt{\frac{\phi_s}{n_*}} \\
    & ~~~~~~~~~~~~ - C_1 \|\bbeta - \bbeta^*\|_2 \times \gamma \kappa_U^{3/2} \sigma_x^3 \sigma_\varepsilon \left(\frac{\sqrt{\phi_*} \sqrt{\phi_s}}{\bar{n}} + \frac{\phi_s \sqrt{\phi_*}}{n_\dagger} \right) \\
    & ~~~~~~~~~~~~ - C_1 \gamma |S^*\setminus S| \kappa_U \sigma_x^2 \sigma_\varepsilon^2 \frac{t+ \log p}{\bar{n}} - C_1 \gamma \kappa_U \sigma_x \sigma_\varepsilon^2 \frac{\phi_s}{n_*} \\
    &= \mathsf{I}_1 + \mathsf{I}_2 + \mathsf{I}_3 + \mathsf{I}_4 + \mathsf{I}_5,
\end{align*} where $(a)$ follows from \cref{prop:population-lb}, $(b)$ follows from the bounds in \cref{lemma:instance-dependent-bound-j-highdim} \& \ref{lemma:instance-dependent-bound-r-highdim} together with the facts that $\log(|\mathcal{E}|) \lesssim \log p$, $\phi_s/n_* \le \tilde{s}(\log(ep/s^*)+t)/ n_* \le 1$, and $|S\setminus S^*| \log(ep/|S\setminus S^*|) + s^* \le 2|S\setminus S^*| \log(ep/s^*) $ because $|S\setminus S^*| \ge s^*$. For $\mathsf{I}_1$, we have $\mathsf{I}_1 \ge \frac{5}{12}\kappa_L \|\bbeta - \bbeta^*\|_2^2$ whenever
\begin{align*}
    C_1 \kappa_U^2\sigma_x^2 \gamma \sqrt{\frac{\phi_s}{n_*}} \le C_1 \kappa_U^2 \sigma_x^2 \gamma \sqrt{\frac{\tilde{s} \{\log(ep/s^*) + t\}}{n_*}} \le \frac{\kappa_L}{12}.
\end{align*} Denote $\Diamond=0.5 (6C_1)^2 (\gamma/\kappa_L)^2 \kappa_U^3 \sigma_x^4\sigma_\varepsilon^2$. 
For $\mathsf{I}_2$, it follows from \cref{lemma:removecondition7} that
\begin{align*}
    \mathsf{I}_2 &\ge \frac{\gamma}{6} \mathsf{J}(\bbeta; \bomega) - C_1 \gamma \sqrt{2\kappa_U^2\|\bbeta - \bbeta^*\|_2^2 + 2\mathsf{J}(\bbeta; \bomega)} \kappa_U^{1/2} \sigma_x \sigma_\varepsilon \sqrt{\frac{\phi_s}{n_*}} \\
    &\ge \frac{\gamma}{6} \mathsf{J}(\bbeta;\bomega) - \frac{\kappa_L}{12} \|\bbeta-\bbeta^*\|_2^2 - 4\Diamond \cdot \frac{\phi_s}{n_*} - \frac{\gamma}{6} \mathsf{J}(\bbeta;\bomega) \\
    &\ge - \frac{\kappa_L}{12} \|\bbeta-\bbeta^*\|_2^2 - 4\Diamond \cdot \frac{(\alpha_{\ell+1} s^*) \{\log(ep/s^*)+t\}}{n_*}
\end{align*}
For $\mathsf{I}_3$ and $\mathsf{I}_4$, it follows from the fact $xy \le \frac{1}{2} (x^2 + y^2)$ that
\begin{align*}
    \mathsf{I}_3 \ge -\frac{\kappa_L}{12} \|\bbeta - \bbeta^*\|_2^2 - \Diamond \cdot \frac{(\alpha_{\ell+1} s^*) \{\log(ep/s^*) + t\} }{n_*} 
\end{align*} and
\begin{align*}
    \mathsf{I}_4 \ge & ~-\frac{\kappa_L}{12} \|\bbeta - \bbeta^*\|_2^2 - 2\Diamond \cdot \frac{(\alpha_{\ell+1} s^*) \cdot \phi_* \left\{\log(ep/s^*) + t\right\}}{\bar{n}^2} \\
    &~~~~~~~~~~~~ - 2\Diamond \cdot \frac{(\alpha_{\ell+1} s^*) \cdot \phi_* \tilde{s} \{\log(ep/s^*) + t\}^2}{n_\dagger^2}.
\end{align*} For $\mathsf{I}_5$, it follows from \eqref{eq:low-dim-bound-i4-1} that
\begin{align*}
\mathsf{I}_5 &\ge -C_1 \|\bbeta - \bbeta_*\|_2^2  \gamma \left\{\frac{\kappa_U \sigma_x^2\sigma_\varepsilon^2}{ \min_{j\in S^*} |\bbeta_j^ *|^2}\frac{t + \log(es^*p)}{\bar{n}}\right\} - C_1 \gamma \kappa_U\sigma_x\sigma_\varepsilon^2 \frac{(\alpha_{\ell+1} s^*) \{\log(ep/s^*) + t\}}{n_*} \\
&\ge -\frac{\kappa_L}{12} \|\bbeta - \bbeta^*\|_2^2 - C_1 \gamma \kappa_U\sigma_x\sigma_\varepsilon^2 \frac{(\alpha_{\ell+1} s^*) \{\log(ep/s^*) + t\}}{n_*}
\end{align*} if $\bar{n} \ge 12C_1 (\gamma/\kappa_L) \kappa_U\sigma_x^2\sigma_\varepsilon^2 (t + \log p) / \min_{j\in S^*} |\beta_j^*|^2$. 
Recall that $\zeta = t + \log(ep/s^*)$. 
Putting all the pieces together and plugging $\zeta$ in, we obtain that, for any $\bbeta \in \mathcal{B}_{\alpha_{\ell+1}s^*} \setminus \mathcal{B}_{\alpha_{\ell}s^*}$
\begin{align*}
    & ~\lambda \|\bbeta\|_0 - \lambda \|\bbeta^*\|_0 + \hat{\mathsf{Q}}(\bbeta) - \hat{\mathsf{Q}}(\bbeta^*) \\
    & ~~~~ \overset{(a)}{\ge} \lambda \left(\alpha_{\ell} -  1\right) s^* + \hat{\mathsf{Q}}(\bbeta) - \hat{\mathsf{Q}}(\bbeta^*) \ge  \frac{\lambda}{4} \alpha_{\ell+1} s^* + \hat{\mathsf{Q}}(\bbeta) - \hat{\mathsf{Q}}(\bbeta^*) \\
    & ~~~~ \ge \frac{1}{2}\left(\kappa_L - \frac{\kappa_L}{6} \times 5\right) \|\bbeta - \bbeta^*\|_2^2 \\
    & ~~~~~~~~~~~~ + \alpha_{\ell+1} s^* \Bigg\{\frac{\lambda}{4} - 6\Diamond \cdot \left(\frac{\zeta}{n_*} + \frac{\phi_* \cdot \zeta}{\bar{n}^2} + \frac{\phi_* \tilde{s}\cdot \zeta^2}{n_\dagger^2}\right)\Bigg\},
\end{align*} where $(a)$ follows from the fact that $\bbeta \notin \mathcal{B}_{\alpha_\ell s^*}$. Recall our choice of $\tilde{s}$. We can conclude that if
\begin{align}
\label{eq:proof-lambda-req2}
\begin{split}
    \lambda &\ge 432C_1^2 \kappa_U^{3} \sigma_x^4 (\gamma/\kappa_L)^2 \sigma_{\varepsilon}^2 \left(\frac{\zeta}{n_*} + \frac{\zeta \phi_*}{\bar{n}^2}\right) + 432C_1^2 \kappa_U^{3/2} \sigma_x^2 (\gamma/\kappa_L) \sigma_\varepsilon^2 \frac{\zeta \sqrt{\phi_*}}{n_\dagger} \\
    &\ge 24\Diamond \cdot \left(\frac{\zeta}{n_*} + \frac{\phi_* \cdot \zeta}{\bar{n}^2} + \frac{\phi_* \tilde{s}\cdot \zeta^2}{n_\dagger^2} \right),
\end{split}
\end{align} then 
\begin{align*}
    \underbrace{\left\{\forall \bbeta \in \mathcal{B}_{\alpha_{\ell+1} s^*} \setminus \mathcal{B}_{\alpha_{\ell} s^*}, ~~ \hat{\mathsf{Q}}(\bbeta) + \lambda \|\bbeta\|_0 - \hat{\mathsf{Q}}(\bbeta^*) - \lambda \|\bbeta^*\|_0 > 0\right\}}_{C_t(\ell)} \subseteq \mathcal{A}_{3,u_{\ell}}(\alpha_{\ell+1} s^*) \cup \mathcal{A}_{4,u_{\ell}}(\alpha_{\ell+1} s^*).
\end{align*}

\noindent \emph{Case 2, $\alpha_{\ell+1} s^* \ge \tilde{s}$.} Observe that $\alpha_{\ell} s^* \ge \frac{1}{2} \alpha_{\ell+1} s^* \ge \frac{1}{2}\tilde{s}$ in this case. Hence the following holds
\begin{align*}
    \lambda \|\bbeta\|_0 - \lambda \|\bbeta^*\| + \hat{\mathsf{Q}}(\bbeta) - \hat{\mathsf{Q}}(\bbeta^*) \ge \frac{\lambda}{2} \tilde{s} - \hat{\mathsf{Q}}(\bbeta^*) - \lambda s^* \ge \frac{\lambda}{4} \tilde{s} - \hat{\mathsf{Q}}(\bbeta^*)
\end{align*} provided $4s^* \le \tilde{s}$. At the same time, if $t\le n_*$ and $\bar{n} \ge \kappa_U\sigma_x^2 \gamma s^*( \log(2s^*|\mathcal{E}|) + t)$, then it follows from \cref{lemma:lemma-qhat-beta-star} that under $\mathcal{A}_{5,t}$, 
\begin{align*}
    \hat{\mathsf{Q}}(\bbeta^*) \le (2 + C_2) \sigma_\varepsilon^2,
\end{align*} where $C_2 = c_1$ is the constant in \cref{lemma:lemma-qhat-beta-star}. At the same time, it follows from our construction of $\tilde{s}$ that
\begin{align*}
    (\tilde{s})^{-1} \le (12C_1)^2 \kappa_U^4\sigma_x^4 (\gamma/\kappa_L)^2 \frac{\zeta}{n_*} + \kappa_U^{3/2} \sigma_x^2 (\gamma/\kappa_L) \frac{\zeta \sqrt{\phi_*} }{n_\dagger}.
\end{align*}
Therefore, we can argue that $\mathcal{C}_{t}(\ell) \subseteq \mathcal{A}_{5,t}$ when 
\begin{align}
\label{eq:lambda-req-3}
\lambda \ge 8(2+C_2) (16C_1)^2 \kappa_U^4\sigma_x^4 \sigma_\varepsilon^2 (\gamma/\kappa_L)^2 \frac{\zeta}{n_*} + 8(2+C_2) \kappa_U^{3/2} \sigma_x^2 \sigma_\varepsilon^2 (\gamma/\kappa_L) \frac{\zeta \sqrt{\phi_*} }{n_\dagger} \ge 8(2+C_2) \sigma_\varepsilon^2 \tilde{s}^{-1}.
\end{align}

Now, we combine the two cases. According to the above discussion, we find that under the conditions 
\begin{align*}
    \tilde{s} \ge 4s^*\qquad \text{and} \qquad \bar{n} \ge C' \kappa_U\sigma_x^2 \gamma (t + \log p) \left\{ s^* +  \sigma_\varepsilon^2 / (\kappa_L \min_{j\in S^*} |\beta_j^*|^2)\right\},
\end{align*} we have \eqref{eq:proof-small-hats-peeling-argument} holds if our choice of $\lambda$ satisfies \eqref{eq:proof-lambda-req2} and \eqref{eq:lambda-req-3}. It then follows from the union bound that
\begin{align*}
    \mathbb{P}\left[\mathcal{C}_t^c\right] \le \sum_{\ell=1}^{\lceil \log_2(p/s^*) \rceil - 1} \left\{\mathbb{P}[\mathcal{A}_{3,u_\ell}(\alpha_{\ell+1} s^*)^c] + \mathbb{P}[\mathcal{A}_{4,u_\ell} (\alpha_{\ell+1} s^*)] \right\}+ \mathbb{P}(\mathcal{A}_{5,t}^c) \le 3e^{-t}.
\end{align*} This completes the proof.

\end{proof}


\subsection{Technical Lemmas}

\begin{lemma}
\label{lemma:removecondition7}
Under Conditions \ref{cond0-model} and \ref{cond1-well-condition}, we have
\begin{align*}
    \forall \bbeta \in \mathbb{R}^p \qquad \sum_{e\in \mathcal{E}} \omega^{(e)} \left\|\mathbb{E}[\bx_{\supp(\bbeta)}^{(e)} \varepsilon^{(e)}] \right\|_2^2 \le 2 \mathsf{J}(\bbeta;\bomega) + 2\kappa_U^2 \|\bbeta - \bbeta^*\|_2^2
\end{align*}
\end{lemma}
\begin{proof}[Proof of \cref{lemma:removecondition7}]
Let $S=\supp(\bbeta)$, it follows from the definition of $\mathsf{J}$ that
\begin{align*}
    \mathsf{J}(\bbeta;\bomega) = \sum_{e\in \mathcal{E}} \omega^{(e)} \|\mathbb{E}[(y^{(e)} - \bbeta^\top \bx^{(e)}) \bx_S^{(e)}]\|_2
    &= \sum_{e\in \mathcal{E}} \omega^{(e)} \left\|\mathbb{E}\left[(\varepsilon^{(e)} + (\bbeta^*)^\top \bx^{(e)} - \bbeta^\top \bx^{(e)}) \bx_S^{(e)}\right]\right\|_2.
\end{align*} Denote $\overline{S} = S^* \cup S$. Then it follows from the fact $(a+b)^2 \le 2(a^2+b^2)$ and \cref{cond1-well-condition} that
\begin{align*}
    \qquad \sum_{e\in \mathcal{E}} \omega^{(e)} \left\|\mathbb{E}[\bx_{S}^{(e)} \varepsilon^{(e)}] \right\|_2^2 &= \sum_{e\in \mathcal{E}} \omega^{(e)} \left\|\mathbb{E}[\bx_S^{(e)} \varepsilon^{(e)} + \bx_S^{(e)} (\bx^{(e)}_{\overline{S}})^\top  (\bbeta^*-\bbeta)_{\overline{S}} ] - \bSigma^{(e)}_{S,\overline{S}} (\bbeta^*-\bbeta)_{\overline{S}}\right\|_2^2 \\
    &\le 2 \sum_{e\in \mathcal{E}} \omega^{(e)} \left\|\mathbb{E}[\bx_S^{(e)} \varepsilon^{(e)} + \bx_S^{(e)} (\bx^{(e)}_{\overline{S}})^\top  (\bbeta^*-\bbeta)_{\overline{S}} ]\right\|_2^2 + 2\kappa_U^2 \|\bbeta - \bbeta^*\|_2^2 \\
    &\le 2\mathsf{J}(\bbeta;\bomega) + 2\kappa_U^2 \|\bbeta - \bbeta^*\|_2^2.
\end{align*} This completes the proof.
\end{proof}


\begin{lemma}
\label{lemma:x-product-epsilon}
Under \cref{cond1-well-condition} and \cref{cond3-sub-gaussian-eps}, we have that, for any $e\in \mathcal{E}$,
\begin{align*}
\left\| \mathbb{E}[\varepsilon^{(e)} \bx^{(e)}]	\right\|_2 \le \sigma_\varepsilon \kappa_U^{1/2}.
\end{align*}
\end{lemma}

\begin{proof}[Proof of \cref{lemma:product-sub-gaussian}]
	Let $\bz = ((\bx^{(e)})^\top, \varepsilon^{(e)})^\top$ be a $p+1$-dimensional random vector. Define the matrix $\bA = \mathbb{E}[\bz \bz^\top]$, we have
    \begin{align*}
		\bA = \begin{bmatrix}	
			\bSigma^{(e)} & \mathbb{E}[\bx^{(e)} \varepsilon^{(e)}] \\
			(\mathbb{E}[\bx^{(e)} \varepsilon^{(e)}])^{\top} & \mathbb{E}[|\varepsilon^{(e)}|^2]
 		\end{bmatrix}.
	\end{align*}
    The matrix is positive semi-definite. Combining this with the fact that $\bSigma^{(e)}$ is invertible indicates that the Schur complement is non-negative, that is,
	\begin{align*}
		\mathbb{E}[|\varepsilon^{(e)}|^2] - \left(\mathbb{E}[\varepsilon^{(e)} \bx^{(e)}]\right)^\top \left(\bSigma^{(e)} \right)^{-1} \left(\mathbb{E}[\varepsilon^{(e)} \bx^{(e)}] \right) \ge 0.
	\end{align*} Hence we have
	\begin{align*}
		\left\|\mathbb{E}[\varepsilon^{(e)} \bx^{(e)}]\right\|^2 \kappa_U^{-1} &\overset{(a)}{\le} \left\|\mathbb{E}[\varepsilon^{(e)} \bx^{(e)}]\right\|^2 \lambda_{\min}\{(\bSigma^{(e)})^{-1}\} \\
		&\le \left(\mathbb{E}[\varepsilon^{(e)} \bx^{(e)}]\right)^\top \left(\bSigma^{(e)} \right)^{-1} \left(\mathbb{E}[\varepsilon^{(e)} \bx^{(e)}] \right) \le \mathbb{E}[|\varepsilon^{(e)}|^2] \overset{(b)}{\le} \sigma_\varepsilon^2.
	\end{align*} where $(a)$ follows from \cref{cond1-well-condition} and $(b)$ follows from \cref{cond3-sub-gaussian-eps}. This completes the proof.
\end{proof}


\iffalse
\begin{lemma}
\label{lemma:removecondition7}
Under \cref{cond0-model} and \ref{cond1-well-condition}, we have
\begin{align*}
    \sum_{e\in \mathcal{E}} \omega^{(e)} \|\mathbb{E}[\bx_S^{(e)} \varepsilon^{(e)}] \|_2^2 \le \Bigg\|\sum_{e\in \mathcal{E}} \omega^{(e)} \mathbb{E}[\bx_S^{(e)} \varepsilon^{(e)}] \Bigg\|_2^2 + 2\kappa_U^2\|\bbeta_{S^*\setminus S}^*\|_2^2 + 2\kappa_U^2 \sum_{e\in \mathcal{E}} \omega^{(e)} \|\bbeta^{(e,S)} - \bar{\bbeta}^{(S)}\|_2^2
\end{align*}
\end{lemma}
\begin{proof}[Proof of \cref{lemma:removecondition7}]
Denote $\bu^{(e)} = \mathbb{E}[\bx_S^{(e)} \varepsilon^{(e)}] \in \mathbb{R}^{|S|}$ and $\bar{\bu} = \sum_{e\in \mathcal{E}} \omega^{(e)} \bu^{(e)}$. We have
\begin{align*}
    \sum_{e\in \mathcal{E}} \omega^{(e)} \|\mathbb{E}[\bx_S^{(e)} \varepsilon^{(e)}] \|_2^2 &= \sum_{e\in \mathcal{E}} \omega^{(e)} \|\bu^{(e)}-\bar{\bu} + \bar{\bu}\|_2^2 \\
    &= \sum_{e\in \mathcal{E}} \omega^{(e)} \left\{\|\bu^{(e)} - \bar{\bu}\|_2^2 + \|\bar{\bu}\|_2^2 + 2\bar{\bu}^\top (\bu^{(e)} - \bar{\bu})\right\} \\
    &= \|\bar{\bu}\|_2^2 + \sum_{e\in \mathcal{E}} \omega^{(e)} \|\bu^{(e)} - \bar{\bu}\|_2^2
\end{align*} where the last identity follows from the fact that $\sum_{e\in \mathcal{E}} \omega^{(e)}=1$ and the definition of $\bar{\bu}$. Denote $T = S^* \setminus S$, observe that 
\begin{align*}
    \bu^{(e)} = \bSigma^{(e)}_S \bbeta_S^{(e,S)} - \bSigma_{S,T}^{(e)} \bbeta_T^*
\end{align*} using the representation of $\bbeta_S^{(e,S)}$. Then we have
\begin{align*}
    \sum_{e\in \mathcal{E}} \omega^{(e)} \|\bu^{(e)} - \bar{\bu}\|_2^2 &= \inf_{\bv \in \mathbb{R}^{|S|}} \sum_{e\in \mathcal{E}} \omega^{(e)} \|\bu^{(e)} - \bv\|_2^2 \\
    &= \inf_{\bv \in \mathbb{R}^{|S|}} \sum_{e\in \mathcal{E}} \omega^{(e)} \|\bSigma^{(e)}_S \bbeta_S^{(e,S)} - \bSigma_{S,T}^{(e)} \bbeta_T^* - \bv\|_2^2 \\
    &\le \inf_{\bv \in \mathbb{R}^{|S|}} \sum_{e\in \mathcal{E}} \omega^{(e)} \left\{2\|\bSigma^{(e)}_S \bbeta_S^{(e,S)} - \bv\|_2^2 + 2\|\bSigma_{S,T}^{(e)} \bbeta_T^* \|_2^2\right\} \\
    &\le 2\kappa_U^2 \| \bbeta_T^* \|_2^2 + 2\kappa_U^2 \inf_{\bv \in \mathbb{R}^{|S|}} \sum_{e\in \mathcal{E}} \omega^{(e)} \| \bbeta_S^{(e,S)} - \bv\|_2^2 \\
    &= 2\kappa_U^2 \|\bbeta_T^*\|_2^2 + 2\kappa_U^2 \sum_{e\in \mathcal{E}} \omega^{(e)} \| \bbeta_S^{(e,S)} - \bar{\bbeta}^{(S)}\|_2^2
\end{align*} Putting all the pieces together, we have
\begin{align*}
\sum_{e\in \mathcal{E}} \omega^{(e)} \|\mathbb{E}[\bx_S^{(e)} \varepsilon^{(e)}] \|_2^2 \le \|\bar{\bu}\|_2^2 + 2\kappa_U^2 \|\bbeta_T^*\|_2^2 + 2\kappa_U^2 \sum_{e\in \mathcal{E}} \omega^{(e)} \| \bbeta_S^{(e,S)} - \bar{\bbeta}^{(S)}\|_2^2.
\end{align*} This completes the proof.

\end{proof}
\fi

\section{Omitted Discussions in Main Text}
\label{sec:discussion}

\subsection{Comparison with \cite{fan2014endogeneity}}
\label{subsec:discussionliao}

Both \cite{fan2014endogeneity} and our work aim to deal with the problem of endogeneity from a technical perspective. Our constructed focused linear invariance regularizer is similar to their developed FGMM criterion function from the view of the \emph{over-identification} idea they proposed. 

Briefly speaking, we say a parameter $\tilde{\bbeta}$ is over-identified if there are more restrictions than the degree of freedom. When $|\mathcal{E}|=1$, there will be exponential number of distinct $\bbeta$ satisfying $\mathsf{J}(\bbeta;\bomega)=0$. To see this, for any $S \subseteq [p]$, one can find some $\bbeta$ with $\supp(\bbeta) \subseteq S$ satisfying $|S|$ constraints that $\mathbb{E}[(y^{(1)} - \bbeta_S^\top \bx_S^{(1)}) x_j^{(1)}] = 0$ for any $j \in [S]$ because the degree of freedom and the number of constraints are both $|S|$. However, things may be different when $|\mathcal{E}| \ge 2$. In this case, for a given $S \subseteq [p]$, one need to find some $\bbeta$ with a degree of freedom $|S|$ satisfies $|\mathcal{E}| \cdot |S|$ constraints 
\begin{align*}
    \forall e \in \mathcal{E}, j\in [S], ~~~~~~ \mathbb{E}\left[x^{(e)}_j (y^{(e)} - \bbeta_S^\top \bx_S^{(e)})\right] = 0
\end{align*}
simultaneously to let $\mathsf{J}(\bbeta;\bomega)=0$, which does not hold in general.

According to the above discussion, the focused linear invariance regularizer shares a similar spirit with their proposed FGMM from a technical viewpoint since the over-identification of our regularizer comes from multiple environments. In contrast, theirs come from two instrumental variables or two marginal features of the covariate, for example, $x_j$ and $x_j^2$. However, the statistical models the two papers work on are different. We briefly remark on the differences between our method and theirs using marginal nonlinear features as follows:
\begin{itemize}
    \item[1.] We are working with multiple environment settings. The necessity of heterogeneous environments and the potential violation of their identification condition are illustrated by \cref{prop:necessary-multiple-envs}. In particular, if two marginal features of the covariate are used in their method, the corresponding identification condition is a sufficient condition of the condition that $S^*$ is the only CE-invariant set among $\mathcal{E}=\{1\}$ other than $\emptyset$.
    \item[2.] We use a linear combination of invariance regularizer and the $L_2$ loss to avoid collapsing to conservative solutions. At the same time, theirs will suffer from the case where two groups of important variables are independent.
    \item[3.] We provide a more explainable identification condition in the context of multiple environments and a non-asymptotic upper bound on the critical threshold of the regularization hyper-parameter $\gamma$; see the intuition explanations in \cref{subsubsec:debias} and a formal presentation in \cref{subsec:strong-convexity}. While \cite{fan2014endogeneity} has few discussions on the population-level conditions.
    \item[4.] The finite sample analyses are completely different because (1) the objective functions are different; and (2) we establish the variable selection consistency for the global minimizer while they only show that some good local minimizer satisfying the variable selection consistency exists. We should carefully deal with the dependence on the number of environments and apply a novel localization argument to obtain a fast rate or a weak condition for variable selection consistency. 
\end{itemize}

\subsection{Comparison with IRM}
\label{sec:irm}

Our constructed EILLS objective is similar to the invariant risk minimization criterion \citep{arjovsky2019invariant} by letting $\gamma \to \infty$. To see this, when $\gamma\to \infty$, the population-level EILLS objective becomes
\begin{align*}
    \min_{\bbeta \in \mathbb{R}^p} \mathsf{R}(\bbeta;\bomega) ~~~~~~s.t.~~~~~~ \mathsf{J}(\bbeta;\bomega)=0.
\end{align*} This optimization problem finds the most effective solution in the sense of small $L_2$ risk among all the LLS-invariant solutions as discussed in \cref{subsec:focused-regularizer}. However, it is unclear (1) what it implies when $\mathsf{J}(\bbeta;\bomega) \approx 0$; and (2) how large $\gamma$ must be to estimate $\beta^*$ consistently. Such two concerns are important for the non-asymptotic analysis because of finite-sample data and finite choice of $\gamma$. We provide an intuitive explanation in \cref{subsubsec:debias} addressing the above concerns, which can be treated as an informal presentation of our theory in \cref{subsec:strong-convexity}.



\subsection{Debiasing by Regularizer via Bias Differences} 
\label{subsubsec:debias} 

\cref{subsec:focused-regularizer} and \cref{sec:irm} provide a sketchy glimpse at the effects of two losses $\mathsf{J}$ and $\mathsf{R}$ on the solution. The following concerns remain: (1) we only analyze the implication of focused linear invariance regularizer $\mathsf{J}(\bbeta;\bomega)$ on the solution when $\mathsf{J}(\bbeta;\bomega)=0$, but we can not expect $\mathsf{J}=0$ in practice especially when it comes to the analysis of its empirical analogs with finite sample data; (2) the above paragraph characterizes that property of proposed EILLS estimator in $\gamma \to \infty$, while it is still unknown how large $\gamma$ is enough. In this part, we will provide a different perspective. In a high-level viewpoint, the $L_2$ risk will have some bias in the presence of the \emph{pooled linear spurious variables} satisfying $\sum_{e\in \mathcal{E}} \omega^{(e)} \mathbb{E} [\varepsilon^{(e)} x_j^{(e)}] \neq 0$. This is where our proposed regularizer $\mathsf{J}$ comes into play: it debiases the pooled least squares loss using the difference of biases between heterogeneous environments, as illuminated below. Such insight also answers the question of how large $\gamma$ is enough.

\begin{figure}[!t]
    \centering
    \begin{tabular}{cc}
    \subfigure[Small $\gamma^*$: small bias, small bias-difference]{
        \begin{tikzpicture}[scale=1.2]
        \draw[->] (0,0) -- (0,3);
        \draw[->] (0,0) -- (3,0);
        \draw (3, -0.3) node{$x_1$};
        \draw (-0.3, 3) node{$x_2$};
        \draw[mylightblue, dashed] (2, 0) -- (1.8, 0.25) node{\scriptsize $d_1$} -- (1.6, 0.5);
        \draw[myblue, dashed] (2, 0) -- (2.25, 0.2) node{\scriptsize $d_2$} -- (2.5, 0.4);
        \draw[myorange, dashed] (1.6, 0.5) -- (2.05, 0.45) node{\scriptsize $\Delta$} -- (2.5, 0.4);
        \draw[myred, fill] (2, 0) circle[radius=0.06];
        \draw[myred] (2,-0.3) node{${\bbeta}^*$};
        \draw[mylightblue, fill] (1.6,0.5) circle[radius=0.06];
        \draw[mylightblue] (1.2, 0.5) node{$\bbeta^{(1)}$};
        \draw[myblue, fill] (2.5, 0.4) circle[radius=0.06];
        \draw[myblue] (3, 0.4) node {$\bbeta^{(2)}$};
        \end{tikzpicture} 
    }
    &
        \subfigure[Small $\gamma^*$: large bias, large bias-difference]{
        \begin{tikzpicture}[scale=1.2]
        \draw[->] (0,0) -- (0,3);
        \draw[->] (0,0) -- (3,0);
        \draw (3, -0.3) node{$x_1$};
        \draw (-0.3, 3) node{$x_2$};
        \draw[mylightblue, dashed] (2, 0) -- (1.5, 0.9) node{\scriptsize $d_1$} -- (1, 1.8);
        \draw[myblue, dashed] (2, 0) -- (2.4, 0.65) node{\scriptsize $d_2$} -- (2.8, 1.3);
        \draw[myorange, dashed] (1, 1.8) -- (1.9, 1.55) node{\scriptsize $\Delta$} -- (2.8, 1.3);
        \draw[myred, fill] (2, 0) circle[radius=0.06];
        \draw[myred] (2,-0.3) node{${\bbeta}^*$};
        \draw[mylightblue, fill] (1,1.8) circle[radius=0.06];
        \draw[mylightblue] (0.6, 1.8) node{$\bbeta^{(1)}$};
        \draw[myblue, fill] (2.8, 1.3) circle[radius=0.06];
        \draw[myblue] (3, 1) node {$\bbeta^{(2)}$};
        \end{tikzpicture} 
    }
    \\
    \subfigure[Large $\gamma^*$: large bias, small bias-difference]{
        \begin{tikzpicture}[scale=1.2]
        \draw[->] (0,0) -- (0,3);
        \draw[->] (0,0) -- (3,0);
        \draw (3, -0.3) node{$x_1$};
        \draw (-0.3, 3) node{$x_2$};
        \draw[myblue, dashed] (2, 0) -- (1.4, 1.3) node{\scriptsize $d_2$} -- (0.8, 2.6);
        \draw[mylightblue, dashed] (2, 0) -- (1.25, 1.25) node{\scriptsize $d_1$} -- (0.5, 2.5);
        \draw[myorange, dashed] (0.5, 2.5) -- (0.8, 2.6);
        \draw[myred, fill] (2, 0) circle[radius=0.06];
        \draw[myred] (2,-0.3) node{${\bbeta}^*$};
        \draw[mylightblue, fill] (0.5, 2.5) circle[radius=0.06];
        \draw[mylightblue] (0.3, 2.1) node{$\bbeta^{(1)}$};
        \draw[myblue, fill] (0.8, 2.6) circle[radius=0.06];
        \draw[myblue] (1.2, 2.6) node {$\bbeta^{(2)}$};
	\draw[myorange] (0.6, 2.7) node{\scriptsize $\Delta$};
        \end{tikzpicture} 
    }
    &
    \subfigure[Impossible: non-zero bias, zero bias-difference]{
        \begin{tikzpicture}[scale=1.2]
        \draw[->] (0,0) -- (0,3);
        \draw[->] (0,0) -- (3,0);
        \draw[mylightblue, dashed] (2-0.01, 0-0.01) -- (1.5-0.01, 1.5-0.01);
        \draw[myblue, dashed] (2+0.01, 0+0.01) -- (1.5+0.01,1.5+0.01);
        
        \draw (3, -0.3) node{$x_1$};
        \draw (-0.3, 3) node{$x_2$};
        \draw[myred, fill] (2, 0) circle[radius=0.06];
        \draw[myred] (2,-0.3) node{${\bbeta}^*$};
        \draw[myblue, fill] (1.5,1.5) circle[radius=0.06];
        \draw[mylightblue] (1.1, 1.5) node{$\bbeta^{(1)}$};
        \draw[mylightblue, fill=mylightblue] (1.5,1.5+0.06) arc [start angle=90, end angle=270, radius=0.06];
        \draw[myblue] (2.1, 1.5) node{$\bbeta^{(2)}$};
        \draw[mylightblue] (1.7, 0.5) node{\scriptsize $d_1$};
        \draw[myblue] (1.9, 0.9) node{\scriptsize $d_2 $};
        \end{tikzpicture} 
    }
    \end{tabular}
    \caption{A geometric illustration of the bias-difference debiasing idea. We consider the same case where $|\mathcal{E}|=p=2$, $x_1$ is the important variable, and $x_2$ is the pooled linear spurious variable. In each subplot, ${\color{myred} \bbeta^*}$ is the true parameter and $\bbeta^{(e)}$ with $e\in \{{\color{mylightblue} 1}, {\color{myblue} 2}\}$ is the population risk minimizer in each environment. Following the discussion in the text, $d_e = \|\bbeta^* - \bbeta^{(e)}\|_2 \asymp \|\bb^{(e)}\|_2$ quantifies the bias of each environment and $\Delta = \|\bbeta^{(1)} - \bbeta^{(2)}\|_2$ represents the bias-difference. The four plots demonstrate four cases in which the magnitudes of bias and bias-difference vary, leading to different thresholds $\gamma^*$ satisfying $\gamma^* \asymp \left(\frac{ {\color{mylightblue} d_1 } + { \color{myblue} d_2} }{{\color{myorange} \Delta}}\right)^{2}$. The above two plots (a) and (b) are the cases where $\gamma^*$ is of reasonable, constant order. We can see when both bias and bias-difference are relatively small in plot (a) or relatively large in plot (b), and the ratio of the two quantities is within constant order that ${\color{mylightblue} d_1} + {\color{myblue} d_2} \asymp {\color{myorange} \Delta}$, the choice of $\gamma^*$ is also of constant order. However, when the bias is much larger than the bias difference that ${\color{mylightblue} d_1} + {\color{myblue} d_2} \gg {\color{myorange} \Delta}$ in (c), one needs to use a large $\gamma^*$ to accommodate the gain in loss decrease from selecting pooled linear spurious variable $x_2$. (d) present a case where the variable set $\{1,2\}$ is also LLS-invariance across the two environments because ${\color{mylightblue} \bbeta^{(1)}}$ and ${\color{myblue} \bbeta^{(2)}}$ coincides. In this case, our proposed EILLS approach will fail and converge to the spurious solution ${\color{mylightblue} \bbeta^{(1)}}={\color{myblue} \bbeta^{(2)}}$ instead of recovering ${\color{myred} \bbeta^*}$.}
    \label{fig:bias-difference-debiasing}
\end{figure}

We illustrate the phenomenon of debiasing at the population level by considering the simplest case of $|\mathcal{E}|=p=2$ and defer a thorough and rigorous analysis to \cref{subsec:strong-convexity}. Here we let $\omega^{(1)}=\omega^{(2)}=1/2$ and $\bbeta^*=(1,0)$ such that the first variable $x_1$ is an important and invariant variable and the second variable $x_2$ is a linear spurious variable. Moreover, suppose $\lambda(\bSigma^{(e)}) \asymp 1$ where $\lambda(\bS)$ represent the eigenvalues of the symmetric matrix $\bS$. 


Under the above toy model, the population-level excess pooled $L_2$ risk is given 
\begin{align*}
    \mathsf{R}(\bbeta) - \mathsf{R}(\bbeta^*) = (\bbeta - \bbeta^*)^\top \bar{\bSigma} (\bbeta - \bbeta^*) - 2\times (\bbeta - \bbeta^*)^\top \bb
\end{align*} 
where $\bar{\bSigma} = \frac{1}{2} (\bSigma^{(1)} + \bSigma^{(2)})$ and  $\bb = (0, b_2)$ with  $b_2 = \frac{1}{2} \{\mathbb{E}[\varepsilon^{(1)} x_2^{(1)}] + \mathbb{E}[\varepsilon^{(2)} x_2^{(2)}]\}$.
This indicates that $\bbeta^*$ is not the minimizer of $\mathsf{R}(\bbeta)$ when $b_2 \not = 0$
and one can further decrease the loss in the direction of $(\bar{\bSigma})^{-1} \bb$. Furthermore, the  decrease in $\mathsf{R}(\bbeta)$ by moving towards this direction is  bounded from below by
\begin{align*}
    \mathsf{R}(\bbeta) - \mathsf{R}(\bbeta^*) \ge C_1^{-1} \|\bbeta - \bbeta^*\|_2^2 - C_1 \|\bb\|_2^2
\end{align*} for some constant $C_1>0$. Let us see how the invariance regularizer $\mathsf{J}(\bbeta)$ can compromise the bias. Denote $\bb^{(e)}=(0,\mathbb{E}[\varepsilon^{(e)} x_2^{(e)}])$ be the bias vector in each environment $e$. When $\bbeta$ is supported on $\{1,2\}$, we find
\begin{align*}
    \mathsf{J}(\bbeta) - \mathsf{J}(\bbeta^*) &= \frac{1}{2} \|\bSigma^{(1)} (\bbeta - \bbeta^*) - \bb^{(1)} \|_2^2 + \frac{1}{2} \|\bSigma^{(2)} (\bbeta - \bbeta^*) - \bb^{(2)} \|_2^2 \\
    & \ge 2 C_2^{-1} \Big \{\| \bbeta - \bbeta^* - (\bSigma^{(1)})^{-1} \bb^{(1)} \|_2^2 +  \|\bbeta - \bbeta^* - (\bSigma^{(2)})^{-1}\bb^{(2)} \|_2^2 \Big \}\\
    &\ge C_2^{-1} \|(\bSigma^{(1)})^{-1} \bb^{(1)} - (\bSigma^{(2)})^{-1} \bb^{(2)}\|_2
\end{align*} for another constant $C_2>0$. This demonstrates that one needs to pay an extra cost of bias-difference $\geq \gamma C_2^{-1}\|(\bSigma^{(1)})^{-1} \bb^{(1)} - (\bSigma^{(2)})^{-1} \bb^{(2)}\|_2$ when adopting pooled linear spurious variables. Such a cost can compromise the gain of the decrease in $\mathsf{R}(\bbeta)$ when a large $\gamma$ is used. Formally, we have
\begin{align*}
    \forall \bbeta ~~\text{with}~~ \|\bbeta \|_0 = 2, ~~~~~~\mathsf{Q}(\bbeta) - \mathsf{Q}(\bbeta^*) \ge C_1^{-1}\| \bbeta - \bbeta^*\|_2^2
\end{align*} provided
\begin{align}
\label{eq:gamma-star-sketch}
    \gamma \ge \gamma^* = C_1 C_2 \frac{\|\bb\|_2^2}{\|(\bSigma^{(1)})^{-1} \bb^{(1)} - (\bSigma^{(2)})^{-1} \bb^{(2)}\|_2} \asymp \frac{\|(\bbeta^{(1)} - \bbeta^*) + (\bbeta^{(2)} - \bbeta^*)\|_2^2}{\|(\bbeta^{(1)} - \bbeta^*) - (\bbeta^{(2)} - \bbeta^*)\|_2^2}
\end{align} where $\bbeta^{(e)} = \argmin_{\bbeta \in \mathbb{R}^2} \mathsf{R}^{(e)}(\bbeta)$ is the  population risk minimizer for environment $e\in \mathcal{E}$. The R.H.S. of \eqref{eq:gamma-star-sketch} follows from the fact that $\bbeta^{(e)} - \bbeta^* = (\bSigma^{(e)})^{-1} \bb^{(e)}$. We present a geometric illustration of the above discussion and how the bias and bias-difference jointly affect the critical threshold $\gamma^*$ in \cref{fig:bias-difference-debiasing}. 




\subsection{Implementation Details of Simulations in \cref{sec:simulation}}
\label{subsec:implementation}

The structural assignment of the SCM in $e=1$ and $e=2$ are as follows
\begin{align*}
    x^{(e)}_1 &\gets u_1^{(e)} \\
    x^{(1)}_4 &\gets u_4^{(1)} \qquad \qquad \qquad \qquad &x^{(2)}_4 \gets (u_4^{(2)})^2-1\\
    x^{(e)}_2 &\gets \sin(x_4^{(e)}) + u_2^{(e)} \\
    x^{(e)}_3 &\gets \cos(x_4^{(e)}) + u_3^{(e)} \\
    x^{(e)}_5 &\gets \sin(x_3^{(e)} + u_5^{(e)}) \\
    x^{(e)}_{10} &\gets 2.5 x_1^{(e)} + 1.5 x_2^{(e)} + u_{10}^{(e)} \\
    y^{(e)} &\gets 3 x_1^{(e)} + 2x_2^{(e)} - 0.5 x_3^{(e)} + u_{13}^{(e)} \\
    x_6^{(e)} &\gets 0.8 y^{(e)} u_6^{(e)} \\
    x_7^{(1)} &\gets 0.5 x_3^{(1)} + y^{(1)} + u_7^{(1)} & x_7^{(2)} \gets 4 x_3^{(2)} + \tanh(y^{(2)}) + u_7^{(2)} \\
    x_8^{(e)} &\gets 0.5 x_7^{(e)} - y^{(e)} + x_{10}^{(e)} + u_8^{(e)} \\
    x_9^{(e)} &\gets \tanh(x_7^{(e)}) + 0.1 \cos(x_8^{(e)}) + u_9^{(e)} \\
    x^{(e)}_{11} &\gets 0.4 (x_7^{(e)} + x_8^{(e)}) * u_{11}^{(e)} \\
    x_{12}^{(e)} &\gets u^{(e)}_{12}
\end{align*} 
Here $u^{(e)}_1,\ldots, u^{(e)}_{13} \sim \mathcal{N}(\bm{0}, I_{13\times 13})$ for all the $e\in \mathcal{E}$.

\noindent \textbf{Implementation.} We use brute force search to calculate $\hat{\bbeta}_{\mathsf{Q}}$. To be specific, we enumerate all the possible support set $S\in 2^{[p]}$, for each fixed $S$, the EILLS objective \eqref{eq:eills-objective-q} is a quadratic function of $[\bbeta]_S$, whose minimum value $\hat{\bbeta}^{(S)}$ can be found by setting the first order condition to be held, that is, 
\begin{align*}
        \hat{\bbeta}^{(S)} &= \mathop{\mathrm{argmin}}_{\mathrm{supp}(\bbeta) \subseteq S} \sum_{e\in \mathcal{E}} \omega^{(e)} \hat{\mathbb{E}}[|y^{(e)} - \bbeta^\top \bx^{(e)}|^2] + \gamma \sum_{e\in \mathcal{E}} \left\|\hat{\mathbb{E}}[\{y^{(e)} - \bbeta^\top \bx^{(e)}\}\bx_{S}^{(e)}]\right\|_2^2 \\
        &= \left[\hat{\bbeta}^{(S)}_S, \hat{\bbeta}^{(S)}_{S^c}\right] \\
        &= \left[\left\{\sum_{e\in \mathcal{E}} \omega^{(e)} \hat{\bSigma}^{(e)}_S + \gamma \sum_{e\in \mathcal{E}} \omega^{(e)} (\hat{\bSigma}^{(e)}_S)^2\right\}^{-1} \left\{ \sum_{e\in \mathcal{E}} \omega^{(e)} \hat{\mathbb{E}}[\bx_S^{(e)} y^{(e)}] + \gamma \sum_{e\in \mathcal{E}} \omega^{(e)} \hat{\bSigma}_{S} \hat{\mathbb{E}}[\bx_S^{(e)} y^{(e)}]\right\}, \bm{0} \right].
\end{align*}
Then $\hat{\bbeta}_{\mathsf{Q}}$ is assigned to be $\hat{\bbeta}^{(S)}$ with the minimum empirical objective value $\hat{\mathsf{Q}}$. The total computational complexity is $\mathcal{O}(2^p p^3 + \sum_{e\in \mathcal{E}} n^{(e)} p^2)$.

We argue here it is possible to do some relaxation on the objective \eqref{eq:eills-objective-q} such that it may be efficiently solved via gradient descent or other methods. However, minimizing $\hat{\mathsf{Q}}$ is still a non-convex problem. We leave an efficient implementation as future work since this paper focused on a thorough statistical analysis for the multiple environment linear regression.

\noindent \textbf{Implementation for Other Invariance Based Methods.} It is noteworthy that all the invariance-based methods have hyper-parameters related to ``invariance'' besides. For example, the hyper-parameters balancing least squares and invariance regularizer in IRM and Anchor regression, and the Type-I error threshold $\alpha$ in ICP. The criterion for choosing this type of hyper-parameter is fundamentally different from that for hyper-parameters controlling the statistical complexity, such as $L_1$/$L_2$ regularization. We can use a train/valid split for the latter and choose the hyper-parameters using the validation dataset. On the contrary, the optimal hyper-parameter for the former should be tuned using the ``test dataset''. In the comparison experiment, the choice of hyper-parameters for ICP, IRM, and Anchor regression are picked in an oracle manner, that is, we enumerate all the possible hyper-parameters and choose the one that minimizes the $\ell_2$ prediction error $\|\bar{\bSigma}^{1/2}(\hat{\bbeta} - \bbeta^*)\|_2^2$. We set the candidate set to be $\{0,1e-5,1e-4,1e-3,1e-2,1e-1,1\}$ for IRM, $\{0, 1, 2, 4, 8, 10, 15, 20, 30, 40, 60, 80, 90, 100, 150, 200, 500, 1000, 5000, 10000\}$ for Anchor regression. For ICP, we try the Type-I error parameter in $\{0.9, 0.95, 0.99, 0.995\}$. We use the dummy variable $1\{E=1\}$ as the anchor variable for Anchor regression.

\noindent \textbf{Configurations for figures.} \cref{fig:numerical} (a) plots the estimated coefficients $\hat{\beta}_j$ over $\gamma \in \{0, 1, 2, 3, 8\} \cup \{5k: k\in [19]\} \cup \{100j:j\in [9]\} \cup \{1000\ell: \ell\in [10]\}$. \cref{fig:numerical} (b) (c) enumerates $n$ in $\{100k: k\in [10]\} \cup \{1500, 2000\}$ and all the plotted quantities are the average over 500 replications.

\subsection{A Diagnosis on Why Refitting is Worse than Vanilla EILLS} 
\label{sec:diagnosis}

In \cref{fig:numerical} (b), we see that the performance of EILLS with refitting is slightly worse than that of the vanilla EILLS estimator. Here, we provide a possible explanation for why there is a noticeable gap when $n=100$ from the observations in simulation studies and leave the finer analysis for future studies. \cref{fig:refit} visualizes the solution EILLS and EILLS with refitting produced in 100 replications when $n=100$. In the plot, each point $(x, y)$ with marker $m$ represents a solution $\hat{\bbeta}$ that the method $m$ produces. Here, $x$ denotes the relative $\ell_2$ norm restricted to the true important variables set $S^*$, calculated as $\|\hat{\bbeta}_{S^*}\|_2/\|{\bbeta}^*_{S^*}\|_2$; and $y$ represents the relative $\ell_2$ norm restricted to the pooled linear spurious variables set $G$, expressed as $\|\hat{\bbeta}_{G}\|_2/\|\bar{\bbeta}_{G}\|_2$ with $\bar{\bbeta} = \bar{\bbeta}^{([12])}$.  The solutions in each trial are connected by an arrow from EILLS to EILLS with refitting, indicating that the latter is the refitted version of the former. 

As shown in \cref{fig:refit}, some solutions falsely select variables in $G=\{7,8,9\}$. For these cases, the EILLS solutions tend to be more dependent on the variables in $S^*$ and less dependent on the variables in $G$ than the EILLS with refitting solution. In other words, the EILLS method yields a more robust solution when the variable selection property fails to hold. We guess that will contribute to the $\ell_2$ estimation error gap between the two methods.

    \begin{figure*}[t!]
    \centering
        \includegraphics[scale=0.5]{figures/refit.pdf}
    \caption{Visualization of solutions EILLS and EILLS with refitting produce in 100 replications when $n=100$.}
    \label{fig:refit}
    \end{figure*}



\subsection{EILLS by Gumbel Approximation}
\label{sec:gumbel}

In this section, we briefly describe how to use stochastic approximation and Gumbel approximation to handle the combinatorial nature of optimization and let a variant of gradient descent continue to work. This is supported by simulation results when $p=70$. We follow the notations in \citep{gu2024causality}.

The original EILLS objective can be written as 
\begin{align}
\label{fan5}
    (\hat{\beta}, \hat{a}) \in \mathop{\mathrm{argmin}}_{\beta \in \mathbb{R}^p, a\in \{0,1\}^p} \underbrace{\sum_{e\in \mathcal{E}} \hat{\mathbb{E}}[|Y^{(e)} - (\beta \odot a)^\top X^{(e)}|^2] + \gamma \sum_{j=1}^p a_j \left|\hat{\mathbb{E}}[\{Y^{(e)} - (\beta \odot a)^\top X^{(e)}\} X_j] \right|^2}_{\hat{\mathsf{Q}}_\gamma(\beta, a)},
\end{align} where $\odot$ is the point-wise multiplication, i.e., $[x \odot y]_j = x_j y_j$ for any $x, y\in \mathbb{R}^d$. This step is used to disentangle the effect of variable selection and parameter estimation, but the combinatorial nature remains. To avoid this, we first rewrite the optimization as a ``continuous'' optimization:
\begin{align*}
    (\hat{\beta}, \hat{w}) \in \argmin_{\beta \in \mathbb{R}^p, w\in \mathbb{R}^p} \mathbb{E}_{B(w)} \left [ \hat{\mathsf{Q}}_\gamma(\beta, B(w)) \right ],
\end{align*}
where the $j^{th}$ component of $B(w) \in \{0, 1\}^d$ follows an independent Bernoulli with probability of success $\sigma(w_j) = \exp(w_j)/(1+\exp(w_j))$.  This is easily seen by taking $\hat{w} = \mathrm{logit}(\hat a)= \log (\frac{\hat a}{1 - \hat a})$ (taking values $\pm \infty$).  Note that $B_j(w_j) = I(\mathrm{logit} (U_j) \leq w_j)$ is discontinuous in $w_j$ where $U_j \sim$ Uniform[0,1], but can be approximated as
\begin{equation} \label{fan6}
    B_j(w_j) \approx \frac{1}{1 + e^{ (\mbox{\scriptsize logit}(U_j) - w_j))/\tau}} \equiv V_\tau(U_j, w_j) ~~~~ \text{as} ~~~~\tau \to 0^+,
\end{equation}
for which its gradient can be taken.  Let
\begin{align*}
    A_\tau(U, w) = (V_\tau(U_{1}, w_1),\ldots, V_\tau(U_{d}, w_d))^\top \in \mathbb{R}^d
\end{align*} with $\{U_j\}_{j=1}^d$ being i.i.d. uniform random variables.
One can approximate the original objective \eqref{fan5} by
\begin{align}
\label{eq:final}
    (\hat{\beta}, \hat{w}) \argmin_{\beta \in \mathbb{R}^p, w\in \mathbb{R}^p} \mathbb{E}_{U} \left [ \hat{\mathsf{Q}}_\gamma(\beta, A_\tau(U, w)) \right ].
\end{align} 
Since $\mathrm{logit}(U_j) \stackrel{d}{=} U_{j,1}- U_{j, 2}$ with $\{U_{j,1}, U_{j,2}\}_{j=1}^d$ being i.i.d. Gumbel(0,1) random variables, the approximation \eqref{fan6} is also referred to as the Gumbel approximation. Given \eqref{eq:final}, one can adopt the following variants of gradient descent in Algorithm \ref{algo1}.


\begin{algorithm}
\caption{Scaling EILLS to High-dimension by Gumbel Trick}
\begin{algorithmic}[1]
\State \textbf{Hyper-parameter: } number of iterations $T$, hyper-parameter $\gamma$
\State \textbf{Gumbel parameters: } initial/final temperature $(\tau_0,\tau_T)$, anneal rate $\rho$, anneal iteration $T_\tau$.
\State \textbf{Input:} data $\{(X^{(e)}_i, Y^{(e)}_i)\}_{i\in [n], e\in \mathcal{E}}$
\State Initialize $\beta, w$
\State Set $\tau = \tau_0$
\For {$t \in \{1,\ldots, T\}$} 
    \State $\tau = \max(\tau_T, \tau \times \rho)$ \textbf{if} $t ~\mathrm{mod}~ T_\tau = 0$.
    \State Sample $\{U_{j,1}, U_{j,2}\}_{j=1}^p$ from Gumbel(0,1).
    \State Update $\beta, w$ by descending its gradient
        \begin{align*}
                \nabla_{(\beta, w)} \left[\hat{\mathsf{Q}}_\gamma(\beta, A_\tau(U, w))\right]
        \end{align*}
\EndFor
\State \textbf{Output:} estimate $\beta \odot \sigma(w)$. 
\end{algorithmic}
\label{algo1}
\end{algorithm}

\subsubsection{Simulation Results}

This section is to illustrate the performance of the above Gumbel-trick optimized EILLS estimator under moderate high dimension $p=70$, where brute force search is not feasible. The data-generating process and the experimental setting are the same as Section 5.2.1 in \cite{gu2024causality}. 

\noindent \textbf{Data Generating Process.} We consider the case where $|\mathcal{E}|=2$ and the data $(X^{(e)}, Y^{(e)})$ in each environment $e\in \{0,1\}$ are generated from two SCMs sharing the same causal relationship between variables. For each trial, we first generate the parent-children relationship among the variables. We enumerate all the $i\in [p+1]$. For each $i\in [p+1]$, we randomly pick at most $4$ parents for the variable $Z_i$ from $\{Z_1,\ldots, Z_{i-1}\}$, this step ensures that the induced graph is a DAG. We use fixed $p=70$, and let the variable $Z_{36}$ be $Y$ and the rest variables constitute the covariate $X$, that is, we let $(Z_1,\ldots, Z_{35}, Z_{36}, Z_{37}, \ldots, Z_{71})=(X_1,\ldots, X_{35}, Y, X_{36}, \ldots, X_{70})$. We also enforce that $Y$ has at least $5$ parents and at least $5$ children by adding parents and children when needed.  The structural assignment for each variable $Z_j$ is defined as
\begin{align*}
    Z_j^{(e)} \gets \sum_{k\in \mathtt{pa}(j)} C_{j,k}^{(e)} f_{j,k}^{(e)}(Z_k^{(e)}) + C_{j, j}^{(e)} \varepsilon_j
\end{align*} where $(\varepsilon_1,\ldots, \varepsilon_{71})$ are independent standard normal random variables. For $j\neq 36$, $f_{j,k}^{(e)}$ are sampled randomly from the candidate functions $\{\cos(x), \sin(x), \sin(\pi x), x, 1/(1+e^{-x})\}$, $C_{j,k}^{(e)}$ are sampled from the uniform distribution on $[-1.5, 1.5]$ with $|C_{j,j}^{(e)}| \ge 0.5$. For $j=36$ and $k<j$, we have $f_{36,k}^{(e)}(x)=x$ and $C_{36,k}^{(0)}\equiv C_{36,k}^{(1)}$ for linearity and invariance. The above data-generating process can be regarded as one observation environment $e=0$ and an interventional environment $e=1$ where the random and simultaneous interventions are applied to all the variables other than the variable $Y$, while the assignment from $Y$'s parent to $Y$ remains and furnishes the target regression function $f^*(x) = \sum_{k\in \mathtt{pa}(36)} C_{36, k}^{(e)} x_k$ in pursuit. In this case, we let $S^* = \mathtt{pa}(36)$ and $\beta^*$ with support set $S^*$ be such that $\beta_j^*=C^{(0)}_{36, k}=C^{(1)}_{36, k}$ for any $k\in S^*$. We also let the noise variance be different for the two environments, i.e., $C_{36,36}^{(0)} \neq C_{36,36}^{(1)}$. Now, the model only has conditional expectation invariance rather than the full conditional distribution invariance. \cref{fig:visualize-linear} (a) visualizes the induced graph in one trial. The complex cause-effect relationships in high-dimensional variables make the problem of estimating $\beta^*$ very challenging.

\begin{figure}[!t]
\centering
\begin{tabular}{cc}
\subfigure[]{
\includegraphics[scale=0.4]{figures/scm-eills.pdf}
}&
\subfigure[]{
\includegraphics[scale=0.3]{figures/eills-gb-seed1.pdf}
}
\end{tabular}
\caption{The visualization of (a) the SCM and (b) the $\sigma({w})$ during training in one trial for the FAIR-Linear estimator. We use different colors to represent the different relationships with $Y$: \myblue{blue} = parent, \myred{red} = child, \myorange{orange} = offspring, \mylightblue{lightblue} = other.}
\label{fig:visualize-linear}
\end{figure}

\medskip
\noindent \textbf{Implementation.} For the EILLS estimator realized by gradient descent with Gumbel trick, we run gradient descent ascent using Adam optimizer with a learning rate of 1e-3, batch size $64$ for $50k$ iterations. We adopt a fixed hyper-parameter $\gamma=10$ and report the performance of the following estimators using the median of the estimation error $\|\hat{\beta} - \beta^*\|_2^2$ over $50$ replications and varying $n\in \{200, 500, 1000, 2000, 5000\}$.
\begin{itemize}
    \item[(1)] Pool-LS: it simply runs least squares on the full covariate $X$ using all the data. 
    \item[(2)] EILLS-GB: Our EILLS estimator with Gumbel approximation that outputs $\beta \odot \sigma(w)$.
    \item[(3)] EILLS-RF: it selects the variables $x_j$ with $\sigma({w_j})>0.7$ of the fitted model in (2), i.e., $\hat{S}=\{j: \sigma(w_j)>0.7\}$, and refits least squares again on $X_{\hat{S}}$ using all the data.
    \item[(4)] Oracle: it runs least squares on $X_{S^*}$ using all the data.
    \item[(5)] Semi-Oracle: it runs least squares on $X_{G^c}$ using all the data, where $G$ is the set of all the descendants of $Y$. Compared with the ERM, it manually removes all the variables that will lead to a biased estimation, but it will also keep uncorrelated variables compared with the full Oracle estimation.  
\end{itemize}

\cref{fig:visualize-linear} (b) visualizes how the Gumbel gate values for different covariables $\sigma(w)$ evolve during training in one trial. We can see that $\sigma(w_j)$ for $j\in S^*$ quickly increases and dominates the values for other variables like children/offspring of $Y$ during the whole training process. 

\medskip
\noindent \textbf{Results. } The results are shown in \cref{fig:simulation-result}. We can see that the square of the $\ell_2$ estimation error $\|\hat{\beta} - \beta^*\|_2^2$ for the pooled least squares estimator (\myyellow{$\times$}) does not decrease and remains to be very large ($\approx 1.5$) as $n$ increases, indicating that it converges to a biased solution. At the same time, the estimation error for EILLS-GB/EILLS-RF (\mylightblue{$\blacklozenge$}/\myblue{$\blacksquare$}) decays as $n$ grows and lies in between that for least squares on $X_{G^c}$ (Semi-Oracle \myorange{$\blacktriangledown$}) and least squares on $X_{S^*}$ (Oracle \myred{$\blacktriangle$}). 


\begin{figure}[!t]
\centering
\includegraphics[scale=0.47]{figures/p70result.pdf}
\caption{The simulation results for linear models with $p=70$. It depicts how the median estimation errors (based on $50$ replications, shown in log scale) for different estimators (marked with different shapes and colors) change when $n$ varies in $\{200, 500, 1000, 2000, 5000\}$ respectively. 
}
\label{fig:simulation-result}
\end{figure}

\subsection{Choice of $\gamma$ in Practice}
\label{sec:gamma-tune}

We argue that the EILLS estimator is not very sensitive to the choice of $\gamma$ when the true causal signal is strong enough, thus one can, for example, adopt $\gamma = 36$, or $\gamma=100$. Such insensitivity is attributed to the discontinuity nature of the loss discussed in \cref{sec:small-gamma}. For example, in the simulation considered in this paper where there exists weak signal $\beta_3 = -0.5$, EILLS discards correct variables only when $\gamma$ is relatively large $\gamma > 3\times 10^3$ and it works pretty well in a wide range of $\gamma \in [15, 3\times 10^3]$ as shown in \cref{fig:numerical} (a).
    
    When the target is to produce the best predictions on unseen future data, we can choose a proper $\gamma$ from a set of candidate $\gamma$-values that has the best out-of-sample performance in certain validation environment(s). To be specific, we consider the two cases, for the first case, we have some training environments $\mathcal{E}_{train}$ and also some validation environments $\mathcal{E}_{valid}$ whose associations between $y$ and $\bx$ are similar to those to be confronted in the future by our belief. We can adopt the following procedure in this case:

    \begin{itemize} 
        \item[Step 1] Set candidate hyper-parameter set $\Gamma$.
        \item[Step 2] For each $\gamma \in \Gamma$, run EILLS estimator using data in $\mathcal{E}_{train}$ with $\gamma$ and get the estimated $\hat{\bbeta}^\gamma$, calculate the worse-case out-of-sample empirical $L_2$ risk among the validation set as
        \begin{align*}
            \hat{R}_\gamma = \max_{e\in \mathcal{E}_{valid}} \frac{1}{n^{(e)}} \sum_{i=1}^{n^{(e)}} \left\{y_i^{(e)} - (\hat{\bbeta}^\gamma)^\top \bx_i^{(e)} \right\}^2.
        \end{align*}
        \item[Step 3] We choose $\gamma$ from $\Gamma$ that corresponds to the minimum out-of-sample empirical $L_2$ risk, that is, $\hat{\gamma} = \argmin_{\gamma \in \Gamma} \hat{R}_\gamma$.
    \end{itemize}

    For the second case, suppose we only have $\mathcal{E}_{train}$, we can adopt the following leave-one-out procedure. 

    \begin{itemize} 
        \item[Step 1] Set candidate hyper-parameter set $\Gamma$.
        \item[Step 2] For fixed $\gamma \in \Gamma$, enumerate all the $e\in \mathcal{E}_{train}$. For each $e \in \mathcal{E}_{train}$, run EILLS estimator using data in $\mathcal{E}_{train} \setminus \{e\}$ with $\gamma$ and get the estimated $\hat{\bbeta}^{\gamma, e}$, calculate its out-of-sample empirical $L_2$ risk in $e$ as
        \begin{align*}
            \hat{R}_{\gamma, e} = \frac{1}{n^{(e)}} \sum_{i=1}^{n^{(e)}} \left\{y_i^{(e)} - (\hat{\bbeta}^{\gamma,e})^\top \bx_i^{(e)} \right\}^2,
        \end{align*} and get its maximum value $\hat{R}_{\gamma} = \max_{e\in \mathcal{E}_{train}} \hat{R}_{\gamma, e}$. 
        \item[Step 3] We choose $\gamma$ from $\Gamma$ that corresponds to the minimum out-of-sample empirical $L_2$ risk, that is, $\hat{\gamma} = \argmin_{\gamma \in \Gamma} \hat{R}_\gamma$.
    \end{itemize}


\subsection{The Interpretation of Small $\gamma$}
\label{sec:small-gamma}

When $\gamma$ is not large enough, i.e., $\gamma < \gamma^*$, the EILLS objective is somewhat similar to running (penalized) pooled least squares on variables excluding some spurious variable whose spuriousness to heterogeneity ratio is smaller than $\gamma$. We use the following variant of the thought experiment to illustrate the idea. Suppose we still do the cow/camel classification using three features $x_1=$~shape, $x_2=$~backgrouond, and $x_3=$~whether the object stands or not. In the two-environment training dataset, 95\% camels/cows on sand/grass, and 95\% camels/cows sit/stand in $\mathcal{D}_1$. Moreover, 90\% camels/cows on sand/grass, and 70\% camels/cows sit/stand in $\mathcal{D}_2$. Under mild conditions akin to \cref{cond-ident-main}, the regularization path of the population-level minimizer of EILLS $\bbeta_\gamma = \argmin_{\bbeta\in \mathbb{R}^p} \mathsf{Q}(\bbeta;\gamma,\bomega)$ can be interpreted as follows. There are change points $\gamma_2>\gamma_1>0$. When $\gamma \in [0,\gamma_1)$, $\bbeta_\gamma$ will be similar to regressing $Y$ on $(x_1,x_2, x_3)$ using all the data; it will threshold $X_3$ and thus be similar to regressing $y$ on $(x_1,x_2)$ using all the data as $\gamma \in (\gamma_1, \gamma_2)$; it will finally recover the ground-truth causal parameter $\bbeta^*$ with $\mathrm{supp}(\bbeta^*)=\{1\}$ when $\gamma > \gamma_2$.

% END INCLUDED SOURCE: tex/supplement_text.tex



\end{document}
